% concatenated from arxiv_parameterfree.tex
\documentclass{article}


% if you need to pass options to natbib, use, e.g.:
%     \PassOptionsToPackage{numbers, compress}{natbib}
% before loading neurips_2025


% ready for submission
%\usepackage{neurips_2025}


% to compile a preprint version, e.g., for submission to arXiv, add add the
% [preprint] option:
%\usepackage[preprint]{neurips_2025}


% to compile a camera-ready version, add the [final] option, e.g.:
\usepackage[final]{neurips_2025}


% to avoid loading the natbib package, add option nonatbib:
%    \usepackage[nonatbib]{neurips_2025}
\renewcommand{\citet}{\citep}

\usepackage[utf8]{inputenc} % allow utf-8 input
\usepackage[T1]{fontenc}    % use 8-bit T1 fonts
\usepackage{hyperref}       % hyperlinks
\usepackage{url}            % simple URL typesetting
\usepackage{booktabs}       % professional-quality tables
\usepackage{amsfonts}       % blackboard math symbols
\usepackage{nicefrac}       % compact symbols for 1/2, etc.
\usepackage{microtype}      % microtypography
\usepackage{xcolor}         % colors


\title{Parameter-free Algorithms for the Stochastically Extended Adversarial Model}


% The \author macro works with any number of authors. There are two commands
% used to separate the names and addresses of multiple authors: \And and \AND.
%
% Using \And between authors leaves it to LaTeX to determine where to break the
% lines. Using \AND forces a line break at that point. So, if LaTeX puts 3 of 4
% authors names on the first line, and the last on the second line, try using
% \AND instead of \And before the third author name.

\iffalse
\author{%
    Shuche Wang \\
    Institute of Operations Research and Analytics\\
    National University of Singapore\\
    Singapore, 117602 \\
    \texttt{shuche.wang@u.nus.edu} \\
    \And
    Adarsh Barik\\
    Department of Computer Science and Engineering\\
    Indian Institute of Technology Delhi\\
    Delhi, 110016, India\\
    \texttt{darshbarik1@iitd.ac.in}\\
    \And
    Peng Zhao\\
    National Key Laboratory for Novel Software Technology\\
    Nanjing University\\
    Nanjing, 210023, China\\
    \texttt{zhaop@lamda.nju.edu.cn}\\
    \And
    Vincent Y.~F. Tan\\
    Department of Mathematics\\
    Department of Electrical and Computer Engineering\\
    Institute of Operations Research and Analytics\\
    \texttt{vtan@nus.edu.sg}
}
\fi

\author{
  Shuche Wang$^{1}$ \quad
  Adarsh Barik$^{2}$ \quad
  Peng Zhao$^{3,4}$ \quad
  Vincent Y.~F. Tan$^{1,5,6}$\\[4pt]
  \small $^{1}$ Institute of Operations Research and Analytics, National University of Singapore \\
  \small $^{2}$ Department of Computer Science and Engineering, Indian Institute of Technology Delhi \\
  \small $^{3}$ National Key Laboratory for Novel Software Technology, Nanjing University \\
  \small $^{4}$ School of Artificial Intelligence, Nanjing University \\
  \small $^{5}$ Department of Mathematics, National University of Singapore \\
  \small $^{6}$ Department of Electrical and Computer Engineering, National University of Singapore \\[3pt]
    \small \texttt{shuche.wang@u.nus.edu, adarshbarik1@iitd.ac.in,}\\[-2pt]
    \small \texttt{zhaop@lamda.nju.edu.cn, vtan@nus.edu.sg}
}

    
  %David S.~Hippocampus\thanks{Use footnote for providing further information
  %  about author (webpage, alternative address)---\emph{not} for acknowledging
  %  funding agencies.} \\
  %Department of Computer Science\\
  %Cranberry-Lemon University\\
  %Pittsburgh, PA 15213 \\
  %\texttt{hippo@cs.cranberry-lemon.edu} \\
  % examples of more authors
  % \And
  % Coauthor \\
  % Affiliation \\
  % Address \\
  % \texttt{email} \\
  % \AND
  % Coauthor \\
  % Affiliation \\
  % Address \\
  % \texttt{email} \\
  % \And
  % Coauthor \\
  % Affiliation \\
  % Address \\
  % \texttt{email} \\
  % \And
  % Coauthor \\
  % Affiliation \\
  % Address \\
  % \texttt{email} \\


\usepackage{algorithm}
\usepackage{algorithmic}

% For theorems and such
\usepackage{amsmath}
\usepackage{amssymb}
\usepackage{mathtools}
\usepackage{amsthm}
\usepackage{tabularx}
\usepackage{booktabs}

\usepackage{bbm}
\usepackage{multirow}
\usepackage{placeins}
\usepackage{caption,enumitem}
\usepackage{colortbl}
\definecolor{LightCyan}{rgb}{0.8, 0.9, 1}

% if you use cleveref..
\usepackage[capitalize,noabbrev]{cleveref}

%%%%%%%%%%%%%%%%%%%%%%%%%%%%%%%%
% THEOREMS
%%%%%%%%%%%%%%%%%%%%%%%%%%%%%%%%
\theoremstyle{plain}
\newtheorem{theorem}{Theorem}[section]
\newtheorem{proposition}[theorem]{Proposition}
\newtheorem{lemma}[theorem]{Lemma}
\newtheorem{corollary}[theorem]{Corollary}
\theoremstyle{definition}
\newtheorem{definition}[theorem]{Definition}
\newtheorem{assumption}[theorem]{Assumption}
\theoremstyle{remark}
\newtheorem{remark}[theorem]{Remark}

%\usepackage{mathtools}
% \mathtoolsset{showonlyrefs}
% Todonotes is useful during development; simply uncomment the next line
%    and comment out the line below the next line to turn off comments
%\usepackage[disable,textsize=tiny]{todonotes}
\usepackage[textsize=tiny]{todonotes}


\DeclareMathOperator*{\argmax}{arg\,max}
\DeclareMathOperator*{\argmin}{arg\,min}

\usepackage{pifont}
\def \Yes {\ding{51}}
\def \No {\ding{55}}



%--------------- Calligraphy \newcommand Declarations -------------------

\newcommand{\cA}{\mathcal{A}}
\newcommand{\cB}{\mathcal{B}}
\newcommand{\cC}{\mathcal{C}}
\newcommand{\cD}{\mathcal{D}}
\newcommand{\cE}{\mathcal{E}}
\newcommand{\cF}{\mathcal{F}}
\newcommand{\cG}{\mathcal{G}}
\newcommand{\cH}{\mathcal{H}}
\newcommand{\cI}{\mathcal{I}}
\newcommand{\cJ}{\mathcal{J}}
\newcommand{\cK}{\mathcal{K}}
\newcommand{\cL}{\mathcal{L}}
\newcommand{\cM}{\mathcal{M}}
\newcommand{\cN}{\mathcal{N}}
\newcommand{\cO}{\mathcal{O}}
\newcommand{\cP}{\mathcal{P}}
\newcommand{\cQ}{\mathcal{Q}}
\newcommand{\cR}{\mathcal{R}}
\newcommand{\cS}{\mathcal{S}}
\newcommand{\cT}{\mathcal{T}}
\newcommand{\cU}{\mathcal{U}}
\newcommand{\cV}{\mathcal{V}}
\newcommand{\cW}{\mathcal{W}}
\newcommand{\cX}{\mathcal{X}}
\newcommand{\cY}{\mathcal{Y}}
\newcommand{\cZ}{\mathcal{Z}}



\newcommand{\ba}{\mathbf{a}}
%\newcommand{\bA}{\mathbf{A}}
\newcommand{\bb}{\mathbf{b}}
\newcommand{\bc}{\mathbf{c}}
\newcommand{\bd}{\mathbf{d}}
\newcommand{\be}{\mathbf{e}}
\newcommand{\bbf}{\mathbf{f}}
\newcommand{\bg}{\mathbf{g}}
\newcommand{\bh}{\mathbf{h}}
\newcommand{\bi}{\mathbf{i}}
\newcommand{\bj}{\mathbf{j}}
\newcommand{\bk}{\mathbf{k}}
\newcommand{\bl}{\mathbf{l}}
%\newcommand{\bm}{\mathbf{m}}
\newcommand{\bn}{\mathbf{n}}
\newcommand{\bo}{\mathbf{o}}
\newcommand{\bp}{\mathbf{p}}
\newcommand{\bq}{\mathbf{q}}
\newcommand{\br}{\mathbf{r}}
\newcommand{\bs}{\mathbf{s}}
\newcommand{\bt}{\mathbf{t}}
\newcommand{\bu}{\mathbf{u}}
\newcommand{\bv}{\mathbf{v}}
\newcommand{\bw}{\mathbf{w}}
\newcommand{\bx}{\mathbf{x}}
\newcommand{\by}{\mathbf{y}}
\newcommand{\bz}{\mathbf{z}}
%\newcommand{\bB}{\mathbf{B}}
%\newcommand{\bZ}{\mathbf{Z}}


\newcommand{\bA}{\mathbf{A}}
\newcommand{\bB}{\mathbf{B}}
\newcommand{\bC}{\mathbf{C}}
\newcommand{\bD}{\mathbf{D}}
\newcommand{\bE}{\mathbf{E}}
\newcommand{\bF}{\mathbf{F}}
\newcommand{\bG}{\mathbf{G}}
\newcommand{\bH}{\mathbf{H}}
\newcommand{\bI}{\mathbf{I}}
\newcommand{\bJ}{\mathbf{J}}
\newcommand{\bK}{\mathbf{K}}
\newcommand{\bL}{\mathbf{L}}
\newcommand{\bM}{\mathbf{M}}
\newcommand{\bN}{\mathbf{N}}
\newcommand{\bO}{\mathbf{O}}
\newcommand{\bP}{\mathbf{P}}
\newcommand{\bQ}{\mathbf{Q}}
\newcommand{\bR}{\mathbf{R}}
\newcommand{\bS}{\mathbf{S}}
\newcommand{\bT}{\mathbf{T}}
\newcommand{\bU}{\mathbf{U}}
\newcommand{\bV}{\mathbf{V}}
\newcommand{\bW}{\mathbf{W}}
\newcommand{\bX}{\mathbf{X}}
\newcommand{\bY}{\mathbf{Y}}
\newcommand{\bZ}{\mathbf{Z}}


\newcommand{\bfa}{{\boldsymbol a}}
\newcommand{\bfb}{{\boldsymbol b}}
\newcommand{\bfc}{{\boldsymbol c}}
\newcommand{\bfd}{{\boldsymbol d}}
\newcommand{\bfe}{{\boldsymbol e}}
\newcommand{\bff}{{\boldsymbol f}}
\newcommand{\bfg}{{\boldsymbol g}}
\newcommand{\bfh}{{\boldsymbol h}}
\newcommand{\bfi}{{\boldsymbol i}}
\newcommand{\bfj}{{\boldsymbol j}}
\newcommand{\bfk}{{\boldsymbol k}}
\newcommand{\bfl}{{\boldsymbol l}}
\newcommand{\bfm}{{\boldsymbol m}}
\newcommand{\bfn}{{\boldsymbol n}}
\newcommand{\bfo}{{\boldsymbol o}}
\newcommand{\bfp}{{\boldsymbol p}}
\newcommand{\bfq}{{\boldsymbol q}}
\newcommand{\bfr}{{\boldsymbol r}}
\newcommand{\bfs}{{\boldsymbol s}}
\newcommand{\bft}{{\boldsymbol t}}
\newcommand{\bfu}{{\boldsymbol u}}
\newcommand{\bfv}{{\boldsymbol v}}
\newcommand{\bfw}{{\boldsymbol w}}
\newcommand{\bfx}{{\boldsymbol x}}
\newcommand{\bfy}{{\boldsymbol y}}
\newcommand{\bfz}{{\boldsymbol z}}

\newcommand{\bfT}{{\boldsymbol T}}
\newcommand{\bfX}{{\boldsymbol X}}

\renewcommand{\Bbb}{\mathbb}
\newcommand{\C}{{\Bbb C}}
\newcommand{\N}{{\Bbb N}}
\newcommand{\Q}{{\Bbb Q}}
\newcommand{\R}{{\Bbb R}}
\newcommand{\Z}{{\Bbb Z}}
\newcommand{\E}{{\Bbb E}}


\newcommand{\bcomment}[1]{\textcolor{blue}{#1}}

\newcommand{\dabs}[1]{\left\|#1\right\|}
\newcommand{\rfrac}[2]{{}^{#1}\!/_{#2}}

\newcommand{\oons}{{\sc OONS}}
\newcommand{\cawe}{{\sc CA-OONS}}
\newcommand{\claoons}{{\sc CLA-OONS}}
\newcommand{\pfonswe}{{\sc PF-OONSwE}}
\newcommand{\fpfons}{{\sc FPF-OONS}}


\newcommand{\cPsi}{\cP_{\mathrm{si}}}
\newcommand{\ocP}{\overline{\cP}}
\newcommand{\regret}{\mathfrak{R}}
\newcommand{\real}{\mathbb{R}}
\newcommand{\inner}[2]{\left\langle #1, #2 \right\rangle}

\newcommand{\adnote}[1]{{\color{blue!80!green!20!red} Adarsh: \emph{#1}}}


\DeclareMathOperator{\tr}{trace}
\renewcommand{\tilde}{\widetilde}
\renewcommand{\hat}{\widehat}
%************************************************************************
%                                                                       *
%            End of preamble and beginning of text.                     *
%                                                                       *
%************************************************************************

\begin{document}


\maketitle

\begin{abstract}
We develop the first parameter-free algorithms for the Stochastically Extended Adversarial (SEA) model, a framework that bridges adversarial and stochastic online convex optimization. Existing approaches for the SEA model require prior knowledge of problem-specific parameters, such as the diameter of the domain $D$ and the Lipschitz constant of the loss functions $G$, which limits their practical applicability. Addressing this, we develop parameter-free methods by leveraging the Optimistic Online Newton Step (OONS) algorithm to eliminate the need for these parameters. We first establish a comparator-adaptive algorithm for the scenario with unknown domain diameter but known Lipschitz constant, achieving an expected regret bound of $\tilde{O}\big(\|u\|_2^2 + \|u\|_2(\sqrt{\sigma^2_{1:T}} + \sqrt{\Sigma^2_{1:T}})\big)$, where $u$ is the comparator vector and $\sigma^2_{1:T}$ and $\Sigma^2_{1:T}$ represent the cumulative stochastic variance and cumulative adversarial variation, respectively. We then extend this to the more general setting where both $D$ and $G$ are unknown, attaining the comparator- and Lipschitz-adaptive algorithm. Notably, the regret bound exhibits the same dependence on $\sigma^2_{1:T}$ and $\Sigma^2_{1:T}$, demonstrating the efficacy of our proposed methods even when both parameters are unknown in the SEA model. 
\end{abstract}

\section{Introduction}
We focus on online convex optimization (OCO)~\cite{cesa2006prediction,hazan2016introduction,orabona2019modern}, a broad framework for sequential decision-making. In each round $t \in [T]$, a learner chooses a point $x_t$ from a convex set $\cX \subseteq \real^d$. The environment then discloses a convex function $f_t : \cX \to \real$, after which the learner incurs a loss $f_t(x_t)$ and updates their decision. The standard way to show the performance is via the \emph{regret}, the total loss relative to a comparator $u\in\cX$, defined as $\regret_T(u) = \sum_{t=1}^T f_t(x_t) - \sum_{t=1}^T f_t(u)$.
\iffalse
\begin{equation*}
\regret_T(u) = \sum_{t=1}^T f_t(x_t) - \sum_{t=1}^T f_t(u).
\end{equation*}
\fi

For convex problems, the regret can be bounded by $O(\sqrt{T})$~\cite{zinkevich2003online}, which is known to be minimax optimal~\cite{abernethy2008optimal}. OCO encompasses two primary frameworks: adversarial OCO~\cite{zinkevich2003online,hazan2007logarithmic}, which aims to minimize regret against arbitrarily chosen loss functions, and stochastic OCO (SCO)~\cite{hazan2007logarithmic,yu2017online}, which minimizes excess risk under i.i.d.\ losses. While both frameworks are well-studied, real-world scenarios typically fall between these theoretical extremes of purely adversarial or stochastic settings. The Stochastically Extended Adversarial (SEA) model proposed in~\citet{sachs2023accelerated} bridges the gap between traditional adversarial and stochastic frameworks in OCO. This hybrid approach serves as an intermediate formulation that captures aspects of both adversarial OCO and SCO settings.

\iffalse
The optimal performance in OCO, SCO, and SEA models traditionally depends on careful step-size tuning, which requires prior knowledge of problem parameters such as the diameter of the decision set and  Lipschitz constants. However, these parameters are often unknown in practice. This limitation has motivated the development of parameter-free algorithms that can achieve comparable regrets without requiring such oracle information. Such algorithms that operate without requiring prior knowledge of problem parameters are termed \emph{``parameter-free''} algorithms. Specifically, parameter-free algorithms include \emph{comparator-adaptive} algorithms (unknown diameter $D$) and \emph{Lipschitz-adaptive} algorithms (unknown Lipschitz constant $G$).


Besides the unknown parameters, the challenge of optimization over unbounded domains has gained attention in recent literature~\cite{cutkosky2018black,chen2021impossible,jacobsen2023unconstrained}. An unbounded decision set $\cX$ allows adversaries to induce arbitrarily large losses for linear loss functions. Traditional approaches typically circumvent this by assuming bounded domains where $\sup_{x,y\in\cX}\|x-y\|_2\leq D$. Consequently, constructing OCO algorithms that operate effectively under unknown parameters and unbounded domain conditions becomes substantially more technically challenging than classical algorithms.
\fi

Optimal performance in OCO, SCO, and SEA models typically relies on careful step-size tuning, which requires prior knowledge of problem parameters such as the diameter of the decision set and Lipschitz constants. However, these parameters are often unknown in practice, motivating the development of \emph{parameter-free} algorithms that achieve comparable regret without requiring such oracle information. Specifically, parameter-free algorithms include \emph{comparator-adaptive} algorithms (unknown diameter $D$) and \emph{Lipschitz-adaptive} algorithms (unknown Lipschitz constant $G$).  A related challenge arises when the decision set $\cX$ is unbounded, allowing adversaries to induce arbitrarily large losses for linear functions. Traditional methods often circumvent this by assuming bounded domains, where $\sup_{x,y\in\cX}\|x-y\|_2\leq D$. Consequently, developing OCO algorithms that remain effective under both unknown parameters and unbounded domains is significantly more challenging than in classical settings~\cite{cutkosky2018black,chen2021impossible,jacobsen2023unconstrained}.


To address these challenges, we propose ``parameter-free'' algorithms for the SEA model, accommodating potentially unbounded decision sets. Using the Optimistic Online Newton Step as our base algorithm, we systematically relax assumptions: first tackling the case of an unknown domain diameter $D$ (potentially infinite) with a known Lipschitz constant $G$, and then extending to the more complex scenario where both $D$ and $G$ are unknown. In the SEA model, at each time step $t$, the learner selects a distribution $\cD_t$ over functions and incurs a loss $f_t(x_t)$, where $f_t$ is sampled from $\cD_t$. The \textit{expected gradient} is denoted as $\nabla F_t(x) = \mathbb{E}_{f_t \sim \mathcal{D}_t}[\nabla f_t(x)]$. %We denote $\mathfrak{G}_t = \sup_{x \in \cX} \| \nabla F_{t}(x) \|$ and use the shorthand $\mathfrak{G}_{1:T}$ to denote $\sum_{t=1}^T \mathfrak{G}_t $.

\iffalse
 {\small 
\begin{table*}[!t]
\centering
\caption{Comparison of the regret bounds of existing results and our three proposed algorithms.}\label{tab:contri}
\resizebox{0.98\textwidth}{!}{
\begin{tabular}{|l|c |c|c|}
\hline
& $D$ known & $G$ known &  Bound  on Expected Regret $\E[\regret_T(u)]$ \\
\hline
Sachs et al.~\citet{sachs2023accelerated} & \multirow{2}{*}{$\surd$} & \multirow{2}{*}{$\surd$} & \multirow{2}{*}{$O\Big(\sqrt{\sigma_{1:T}^2}+\sqrt{\Sigma_{1:T}^2}\Big)$} \\
Chen et al.~\citet{chen2024optimistic} & & & \\
\hline
\oons~(\textbf{Ours}) & $\surd$ & $\surd$ & $\tilde{O}\Big(\sqrt{\sigma_{1:T}^2}+\sqrt{\Sigma_{1:T}^2}\Big)$ \\
\hline
\cawe~(\textbf{Ours})& $\times$ & $\surd$ & $\tilde{O}\Big(\|u\|_2^2 + \|u\|_2(\sqrt{\sigma^2_{1:T}} + \sqrt{\Sigma^2_{1:T}})\Big)$ \\
\hline
\claoons~(\textbf{Ours})& $\times$ & $\times$ & $\tilde{O}\Big(\|u\|_2^2(\sqrt{\sigma_{1:T}^2}+\sqrt{\Sigma_{1:T}^2})+\|u\|_2^4+\sqrt{\sigma_{1:T} + \mathfrak{G}_{1:T} }\Big)$ \\
\hline
\end{tabular}}
\end{table*}}
\fi


\newcolumntype{g}{>{\columncolor{LightCyan}}c}
{\scriptsize 
\begin{table*}[t!]
\centering
\caption{Comparison of the regret bounds of existing results and our proposed algorithms.}
\label{tab:contri}
\resizebox{\columnwidth}{!}{%
\begin{tabular}{gggg}
\hline
\toprule
\rowcolor{white} \textbf{Algorithm} & \textbf{Free of $D$} & \textbf{Free of $G$} & \textbf{Bound on Expected Regret} $\mathbb{E}[\regret_T(u) ]$ \\
\midrule
\rowcolor{white} OFTLR,~OMD & & & \\
\rowcolor{white} \scriptsize{(Sachs et al.~\cite{sachs2023accelerated},~Chen et al.~\cite{chen2024optimistic})} & \multirow{-2}{*}{\No} & \multirow{-2}{*}{\No} & \multirow{-2}{*}{$O\Big(\sqrt{\sigma_{1:T}^2}+\sqrt{\Sigma_{1:T}^2}\Big)$} \\
%\rowcolor{white} OMD & & & \\
%\rowcolor{white} \scriptsize{(Chen et al.~\cite{chen2024optimistic})} & \multirow{-2}{*}{\Yes} & \multirow{-2}{*}{\Yes} & \multirow{-2}{*}{$O\Big(\sqrt{\sigma_{1:T}^2}+\sqrt{\Sigma_{1:T}^2}\Big)$} \\
\rowcolor{white} \textbf{\oons} & & & \\
\rowcolor{white} \scriptsize{(Theorem~\ref{thm:knownDG})} & \multirow{-2}{*}{\No} & \multirow{-2}{*}{\No} & \multirow{-2}{*}{$\tilde{O}\Big(\sqrt{\sigma_{1:T}^2}+\sqrt{\Sigma_{1:T}^2}\Big)$} \\
\textbf{\cawe} & & & \\
\scriptsize{(Theorem~\ref{thm:unknownD})} & \multirow{-2}{*}{\Yes} & \multirow{-2}{*}{\No} & \multirow{-2}{*}{$\tilde{O}\Big(\|u\|_2^2 + \|u\|_2(\sqrt{\sigma^2_{1:T}} + \sqrt{\Sigma^2_{1:T}})\Big)$} \\
\textbf{\claoons} & & & \\
\scriptsize{(Theorem~\ref{thm:fullyfree})} & \multirow{-2}{*}{\Yes} & \multirow{-2}{*}{\Yes} & \multirow{-2}{*}{$\tilde{O}\Big(\|u\|_2^2(\sqrt{\sigma_{1:T}^2}+\sqrt{\Sigma_{1:T}^2})+\|u\|_2^4+\sqrt{\sigma_{1:T} + \mathfrak{G}_{1:T}}\Big)$} \\
\bottomrule
\end{tabular}%
}
\end{table*}}

\paragraph{Main Contributions.} Our main results and contributions are summarized as follows.
\begin{enumerate}[leftmargin=*, label=(\arabic*)]
\item[(1)] We begin by introducing the Optimistic Online Newton Step (OONS) as our foundational algorithm. The \oons\ algorithm is inspired by \citet{chen2021impossible}; however, we incorporate an adaptive step-size $\eta_t$ rather than a fixed step-size throughout the learning process. When the parameters $D$ and $G$ are known, we demonstrate that \oons\ achieves an expected regret bound of $\tilde{O}(\sqrt{\sigma^2_{1:T}} + \sqrt{\Sigma^2_{1:T}})$, matching the state-of-the-art results in terms of dependence on the cumulative stochastic variance $\sigma^2_{1:T}$ and the cumulative adversarial variation $\Sigma^2_{1:T}$~\cite{sachs2023accelerated, chen2024optimistic}. This establishes a solid foundation for our subsequent extensions to parameter-free algorithms.

%We first introduce Optimistic Online Newton Step (OONS) as our base algorithm.   \oons\ algorithm is inspired by \citet{chen2021impossible}, albeit we consider the adaptive step-size $\eta_t$ instead of a fixed step-size through the learning process. With known $D$ and $G$, we prove that \oons~enjoys an $\tilde{O}(\sqrt{\sigma^2_{1:T}} + \sqrt{\Sigma^2_{1:T}})$ expected regret bound, matching the prior state-of-the-art results in terms of dependence on $\sigma^2_{1:T}$ (cumulative stochastic variance) and $\Sigma^2_{1:T}$ (cumulative adversarial variation)~\cite{sachs2023accelerated,chen2024optimistic}. This serves as the foundation for our subsequent parameter-free extensions.
% This foundational analysis provides crucial insights that enable our subsequent parameter-free extensions.
\item[(2)] We introduce the first parameter-free (comparator-adaptive) algorithm for the SEA model that remains effective when the domain diameter $D$ is unknown, provided the Lipschitz constant $G$ is known. This is achieved through a meta-framework wherein each base learner operates within a distinct bounded domain, complemented by the Multi-scale Multiplicative-Weight with Correction (MsMwC) algorithm~\cite{chen2021impossible} for the meta-algorithm's weight updates. This construction yields an expected regret bound of $\tilde{O}(\|u\|_2^2 + \|u\|_2(\sqrt{\sigma^2_{1:T}} + \sqrt{\Sigma^2_{1:T}}))$, where the bound scales with the $\ell_2$-norm of the comparator $u$ without requiring prior knowledge of the domain diameter~$D$.

\item[(3)] We further consider a setting in which {\em both} the domain diameter $D$ and the Lipschitz constant $G$ are unknown. By devising appropriate update rules for the estimation of the domain diameter, we establish an expected regret bound of
$\tilde{O}(\|u\|_2^2(\sqrt{\sigma_{1:T}^2}+\sqrt{\Sigma_{1:T}^2})+\|u\|_2^4+\sqrt{\sigma_{1:T} + \mathfrak{G}_{1:T} })$
where $\sigma_{1:T}$ captures the deviation of the stochastic gradients (excluding squared norms), and $\mathfrak{G}_{1:T}$ denotes the sum of the maximum expected gradients over the sequence.


%We develop the first parameter-free (comparator-adaptive) algorithm for the SEA model that operates effectively when the domain diameter $D$ is unknown but the Lipschitz constant $G$ is known. Through a meta-base framework where each base-learner operates within a different bounded domain, combined with the Multi-scale Multiplicative-weight with Correction (MsMwC) algorithm~\cite{chen2021impossible} for meta-algorithm's weight updates, we achieve an expected regret bound of $\tilde{O}(\|u\|_2^2 + \|u\|_2(\sqrt{\sigma^2_{1:T}} + \sqrt{\Sigma^2_{1:T}}))$. This bound scales with the norm of the comparator $\|u\|_2$ rather than requiring explicit knowledge of the domain diameter $D$. 
%\item[(3)] We consider the comparator and Lipschitz-adaptive setting where {\em both} $D$ and $G$ are unknown. By designed update rules for the domain diameter, we establish an expected regret bound of $\tilde{O}(\|u\|_2^2(\sqrt{\sigma_{1:T}^2}+\sqrt{\Sigma_{1:T}^2})+\|u\|_2^4+\sqrt{\sigma_{1:T} + \mathfrak{G}_{1:T} })$, where $\sigma_{1:T}$ captures the stochastic gradient deviation (without the squared norm) and $\mathfrak{G}_{1:T}$ denotes the sum of maximum expected gradients. %This regret bound is achieved without requiring the creation of any additional experts. 
%\adnote{Can we define it more precisely -- like we do with other quantities in page 3. $\delta_t = \sup_{x \in \mathcal{X}}\mathbb{E}_{f_t \sim \mathcal{D}_t}\left[\|\nabla f_t(x) - \nabla F_t(x)\|_2\right]$. Then we can have $\delta_{1:T} = \mathbb{E}\left[\sum_{t=1}^T \delta_t\right]$}
\end{enumerate}
A summary of our results and the best existing results are included in Table~\ref{tab:contri}. Due to space limitations, we hide the Lipschitz constant $G$ in the $\tilde{O}(\cdot)$-notation in the regret bound of \claoons\ algorithm.

\subsection{Related Work} 
The SEA model~\cite{sachs2023accelerated} is motivated by foundational insights from the \emph{gradual-variation online learning}. The study of gradual variation can be traced back to the works of \citet{hazan2010extracting} and \citet{chiang2012online}, and it has gained significant traction in recent years~\cite{zhao2020dynamic,yan2023universal,zhao2024adaptivity,xie2024gradient,yan2024simple}. Notably, the SEA model has emerged as a practical application of the gradual variation principle~\cite{yan2023universal,xie2024gradient,yan2024simple}. Furthermore, this model serves as a bridge between adversarial OCO and SCO. This intermediate framework is comprehensively understood in the context of expert prediction~\cite{amir2020prediction,ito2021optimal} and the bandit setting~\cite{zimmert2019optimal,lee2021achieving}.

Parameter-free online learning has emerged as a fundamental advancement in machine learning, offering solutions to the critical challenge of parameter tuning in practice. In the baseline scenario, when both the diameter parameter $D$ and the gradient bound $G$ are known, algorithms leveraging Follow the Regularized Leader or Mirror Descent principles achieve the minimax optimal regret bound of  $\regret_T(u)\leq O(GD\sqrt{T})$~\cite{hazan2016introduction}. The field has subsequently progressed to address more practical scenarios where complete parameter knowledge is unavailable. Notably, in the Lipschitz-adaptive setting, it is still possible to attain the same optimal regret bound, differing only by constant factors~\cite{orabona2016coin,cutkosky2019artificial}. Xie et al.~\citet{xie2024gradient} extended these principles to gradient-variation online learning.

In the comparator-adaptive setting, the online learning problem becomes substantially more challenging due to the unknown comparator's magnitude, which could cause the algorithm's predictions to significantly deviate from the optimal solution, leading to a large regret. %This requires the learner to adaptively control $\|x_t\|_2$ based on the unknown comparator norm $\|u\|_2$.
For this challenging scenario, a key result has been established as $\regret_T(u)\leq \tilde{O}(\|u\|_2G\sqrt{T})$~\cite{foster2015adaptive,orabona2016coin,cutkosky2018black,jacobsen2022parameter}. For scenarios where both parameters $D$ and $G$ are unknown, significant progress has been made recently. Cutkosky~\citet{cutkosky2019artificial} developed an algorithm with $\regret_T(u)\leq \tilde{O}(G\|u\|_2^3+\|u\|_2 G\sqrt{T})$, while  Mhammedi~\&~Koolen~\citet{mhammedi2020lipschitz} achieve $\regret_T(u)\leq \tilde{O}(G\|u\|_2^3+G\sqrt{\max_{t\leq T}(\sum_{s=1}^t\|g_s\|_2/\max_{s\leq t}\|g_s\|_2}))$. An alternative approach by \citet{orabona2018scale} presented the regret bound $\regret_T(u)\leq \tilde{O}(\|u\|_2^2G\sqrt{T})$. More recent advances including  \citet{jacobsen2023unconstrained} achieve the regret $\regret_T(u)\leq \tilde{O}(G\|u\|_2\sqrt{T}+L\|u\|_2^2\sqrt{T})$ under the condition that subgradients satisfy $\|g_t\|_2\leq G+L\|x_t\|_2$.  Cutkosky \& Mhammedi~\citet{cutkosky2024fully} further improve it to $\tilde{O}(G\|u\|_2\sqrt{T}+\|u\|_2^2+G^2)$.

Besides  parameter-free algorithms for OCO, \citet{ivgi2023dog} and \citet{khaled2023dowg} studied the parameter-free stochastic gradient descent (SGD) algorithms. Khaled et al.~\citet{khaled2024tuning} introduced the concept of ``tuning-free'' algorithms, which achieve performance comparable to optimally-tuned SGD within polylogarithmic factors, requiring only approximate estimates of the relevant problem parameters.

%Despite this series of works for parameter-free algorithms in OCO, these approaches cannot be directly applied to achieve the optimal regret bound for the SEA model without known parameters. This limitation arises because our desired bounds for the SEA model should be characterized in terms of $\sigma^2_{1:T}$ and $\Sigma^2_{1:T}$ rather than the horizon $T$. While Sachs et al.~\citet{sachs2023accelerated} attempted to address this challenge by proposing an algorithm for the unknown strong convexity parameter, their step-size search range still depends on both $D$ and $G$, limiting its full parameter-free adaptivity. 

Although this series of works on parameter-free algorithms in OCO provides valuable insights, these approaches cannot be directly applied to attain the optimal regret bounds for the SEA model without prior knowledge of parameters. This limitation stems from the fact that the desired bounds for the SEA model should be expressed in terms of the variance-like quantities $\sigma^2_{1:T}$ and $\Sigma^2_{1:T}$, rather than the time horizon $T$. While Sachs et al.~\citet{sachs2023accelerated} have attempted to address this issue by proposing an algorithm that adapts to an unknown strong convexity parameter, their step-size search range still depends on both $D$ and $G$, thereby restricting its fully parameter-free adaptivity.

\section{Problem Setup and Preliminaries}\label{sec:setup}
In this section, we formulate the problem setup of the Stochastically Extended Adversarial (SEA) model, present the existing results, and discuss the key challenges.

\subsection{Problem Setup of the SEA Model}
% \textbf{The Stochastically Extended Adversarial (SEA) Model.} %The SEA model~\cite{sachs2023accelerated} is introduced as an intermediate problem setup between adversarial online convex optimization (adversarial OCO) and stochastic online convex optimization (SCO). 
In iteration $t \in [T]$, the learner selects a decision $x_t$ from a convex feasible domain $\cX \subseteq \mathbb{R}^d$, and nature chooses a distribution $\cD_t$ from a set of distributions over functions. Then, the learner suffers a loss $f_t(x_t)$, where  $f_t$ is a random function sampled from the distribution $\cD_t$. The distributions are allowed to vary over time, and by choosing them appropriately, the SEA model reduces to the adversarial OCO, SCO, or other intermediate settings. Additionally, for each $t \in [T]$, the {\em (conditional) expected function} is defined as $F_{t}(x) =\mathbb{E}_{f_t \sim \cD_t}[f_t(x)]$ and the {\em expected gradient} is defined as $\nabla F_{t}(x) =\mathbb{E}_{f_t \sim \cD_t}[\nabla f_t(x)]$. We define $\mathfrak{G}_t := \sup_{x \in \cX} \| \nabla F_{t}(x) \|_2$ to be the largest norm of the expected gradient, and use the shorthand $\mathfrak{G}_{1:T}$ to denote the sum $\sum_{t=1}^T \mathfrak{G}_t $.

Due to the randomness in the online decision-making process, our goal in the SEA model is to bound the \emph{expected regret} with respect to the randomness in the loss functions $f_t$ drawn from the distribution $\cD_t$ against any fixed comparator $u \in \cX$, defined as $\E[\regret_T(u)] \triangleq \E[ \sum_{t=1}^T f_t(x_t) - \sum_{t=1}^T f_t(u)]$.
\iffalse
\begin{equation*} \label{eqn:expect:regret}
\E[\regret_T(u)] \triangleq \E\left[ \sum_{t=1}^T f_t(x_t) - \sum_{t=1}^T f_t(u) \right].
\end{equation*}
\fi
To capture the characteristics of the SEA model, we introduce the following quantities. For each $t \in [T]$, define the {\em (conditional) variance of the gradients} and {\em cumulative stochastic variance} respectively as 
\begin{equation} \label{eqn:variance:grad}
\sigma_{t}^2 = \sup_{x \in \cX} \E_{f_t \sim \cD_t}  \big[ \| \nabla f_t(x) - \nabla F_{t}(x) \|_2^2 \big],\quad \sigma_{1:T}^2 = \E \left[ \sum_{t=1}^T \sigma_t^2 \right],
\end{equation}
which reflect  the stochasticity of the online process. Additionally, we introduce the concepts of stochastic gradient deviation and cumulative gradient deviation to characterize the stochastic variation of gradients, without the squared norm. The stochastic gradient deviation is defined as $\sigma_{t} = \sup_{x \in \cX} \E_{f_t \sim \cD_t}  \big[ \| \nabla f_t(x) - \nabla F_{t}(x) \|_2 \big]$,
\iffalse
\begin{equation*} \label{eqn:deviation:grad}
\sigma_{t} = \sup_{x \in \cX} \E_{f_t \sim \cD_t}  \big[ \| \nabla f_t(x) - \nabla F_{t}(x) \|_2 \big], 
\end{equation*}
\fi
and the {\em cumulative gradient deviation} is defined as $\sigma_{1:T} = \E \big[ \sum_{t=1}^T \sigma_t \big]$.
% \begin{equation*} \label{eqn:cumu:deviation}
% \sigma_{1:T} = \E \left[ \sum_{t=1}^T \sigma_t \right]~.
% \end{equation*}
The {\em cumulative adversarial variation} is defined as
\begin{equation*} 
\Sigma_{1:T}^2= \E \left[ \sum_{t=1}^T \sup_{x \in \cX} \|\nabla F_t(x)- \nabla F_{t-1}(x)\|_2^2\right],
\end{equation*}
where $\nabla F_0(x) = 0$, reflecting the adversarial difficulty. This work aims to provide expected regret bounds that depend on problem-dependent quantities such as $\sigma^2_{1:T}, \Sigma^2_{1:T}$, and $\mathfrak{G}_{1:T}$ instead of $T$.
%Hence, the expectation of regret depends on $\sigma_{1:T}^2$ and $\Sigma_{1:T}^2$ instead of $T$ in this work.

\iffalse
We also let $\bar{\sigma}_T$ and $\bar{\Sigma}_T$ denote the square root of the average stochastic variance and adversarial variation, respectively; that is, $\bar{\sigma}_T=\sigma^2_{1:T}/T$ and $\bar{\Sigma}_T=\Sigma^2_{1:T}/T$.
\fi


Below, we present several assumptions. Note that our results do not rely on \emph{all} of these assumptions; rather, specific assumptions are required for each result, which will be explicitly stated in the theorem.

% , with one general and mild assumption being that the domain $\mathcal{X}$ includes the origin $0$.
  
% \begin{definition}(Bregman divergence).
% Let $\psi: \cX \to \mathbb{R}$ be a continuously differentiable and strictly convex function. The Bregman divergence $D_{\psi}: \cX^2\rightarrow \mathbb{R}$ is given by the recipe
% $D_{\psi}(x, y) \coloneq \psi(x) - \psi(y) - \langle \nabla \psi(y), x-y \rangle.$
% \end{definition}

\begin{assumption}[Boundedness of gradient norms]
\label{ass:G}
The gradient norms of all loss functions are bounded by $G$, i.e., $\max_{t\in[T]}\max_{x\in\cX}\|\nabla f_t(x)\|_2\leq G$.
\end{assumption}

\begin{assumption}[Boundedness of domain]
\label{ass:D}
The diameter of the convex set $\cX$ (the feasible domain) is bounded by $D$ i.e., $\sup_{x,y\in \cX}\|x-y\|_2\leq D$.
\end{assumption}

\begin{assumption}[Smoothness]
\label{ass:L}
For all $t\in [T]$, the expected function $F_t$ is $L$-smooth over $\cX$, i.e., $\|\nabla F_t(x)-\nabla F_t(y)\|_2\leq L\|x-y\|_2$ for all $x,y\in\cX$.
\end{assumption}

\begin{assumption}[Convexity]
\label{ass:convexF}
For all $t\in[T]$, the expected function $F_t$ is convex on $\cX$.
\end{assumption}
%Our   results do not rely on \emph{all} these assumptions; the specific assumptions required are explicitly stated within each result. One global and mild assumption we need is that  the domain $\cX$ contains the origin $0$. 


\textbf{Notations.} Given a positive definite matrix $A$, the   norm induced by $A$ is  $\|x\|_{A}=\sqrt{x^\top Ax}$.  $\Delta_N$ denotes the $(N-1)$-dimensional simplex. Let $\psi: \cX \to \mathbb{R}$ be a continuously differentiable and strictly convex function, the associated Bregman divergence is defined as $D_{\psi}(x, y) \coloneq \psi(x) - \psi(y) - \langle \nabla \psi(y), x-y \rangle$. The notation $O(\cdot)$ hides constants and $\tilde{O}(\cdot)$ additionally hides  $\mathrm{polylog}$ factors.

% \textbf{Parameter-Free Algorithms for SEA Model.}
\subsection{Existing Results for the SEA Model}

\textbf{Bounded domain and gradient norm.}~~  Sachs et al.~\cite{sachs2023accelerated} established a regret bound for the SEA model using both Optimistic Follow-The-Regularized-Leader (OFTRL) and Optimistic Mirror Descent (OMD), given by $\mathbb{E}[\regret_T(u)]= O(\sqrt{\sigma_{1:T}^2}+\sqrt{\Sigma_{1:T}^2})$, achieved by setting the step-size as $\eta_t=\frac{D^2}{\sum_{s=1}^{t-1} \min\{\frac{\eta_s}{2} \|g_s - m_s\|_2^2, D \|g_s - m_s\|_2\}}$, where $g_t=\nabla f_t(x_t)$ and $m_t=g_{t-1}$. Similarly, Chen et al.~\cite{chen2024optimistic} derived the same bound by Optimistic Online Mirror Descent (OMD) with the step-size $\eta_t=\frac{D}{\sqrt{\delta+4G^2+\bar{V}_{t-1}}}$, where $\bar{V}_{t-1}=\sum_{s=1}^{t-1}\|g_s - m_s\|_2^2$ and $\delta>0$.

In all of the above settings, the optimal step-size $\eta_t$ is dependent on the parameters $D$ (the diameter of decision set $\cX$) and $G$ (Lipschitz constant), so there has been a natural motivation to develop algorithms that achieve similar regret bounds \emph{without} knowing such parameters \emph{a priori}. We term such algorithms as \emph{``parameter-free''} algorithms for the SEA model. 

\textbf{Parameter-free algorithm for the SEA model.}~~ Theorem 5 in \citet{jacobsen2022parameter} demonstrates that the parameter-free mirror descent algorithm can be extended to enjoy a gradient-variation regret of $\regret_T(u) \leq \tilde{O}(\|u\|_2 \sqrt{\sum_{t=1}^T \|\nabla f_t(x_t) - \nabla f_{t-1}(x_t)\|_2^2})$. In fact, this can directly yield an expected regret bound for the SEA model scaling with $\tilde{\sigma}_{1:T}^2 := \mathbb{E}\left[\sum_{t=1}^T \mathbb{E}_{f_t \sim \mathcal{D}_t} \left[\sup_x \|\nabla f_t(x) - \nabla F_t(x)\|_2^2\right]\right]$, i.e.,
\begin{equation}\label{eq:tildesigma}
\mathbb{E}[\regret_T(u)]\leq\tilde{O}\left(\|u\|_2 \Big(\sqrt{\tilde{\sigma}_{1:T}^2} + \sqrt{\Sigma_{1:T}^2} \Big)\right).
\end{equation}
%where $\tilde{\sigma}_{1:T}^2$ is defined as $\tilde{\sigma}_{1:T}^2 := \mathbb{E}\left[\sum_{t=1}^T \mathbb{E}_{f_t \sim \mathcal{D}_t} \left[\sup_x \|\nabla f_t(x) - \nabla F_t(x)\|_2^2\right]\right]$, as introduced in~\cite{chen2024optimistic}.
Akin to $\sigma_{1:T}^2$, $\tilde{\sigma}_{1:T}^2$ defined in \cite{chen2024optimistic} also captures the stochastic nature of the SEA model. Furthermore, in the worst case, the bound in \eqref{eq:tildesigma} reduces to $\tilde{O}(\|u\|_2 \sqrt{T})$, matching the best available problem-independent bound. The outer expectation in the definition of $\tilde{\sigma}_{1:T}^2 $ accounts for the randomness in the choice of the distribution $\cD_t$ at each step. Refer to Appendix~\ref{app:tildesigma} for a self-contained proof of~\eqref{eq:tildesigma}.

\textbf{Key Challenge.}~~ However, we emphasize that our goal is to obtain regret bounds scaling with $\sigma^2_{1:T}$, as defined in~\eqref{eqn:variance:grad}. As pointed out in previous work on the SEA model~\citep[Remark~9]{chen2024optimistic}, $\sigma^2_{1:T}$ is more favorable than $\tilde{\sigma}^2_{1:T}$. First, from a mathematical perspective, the latter is generally larger due to the convexity of the supremum operator. The difference between $\sigma^2_{1:T}$ and $\tilde{\sigma}^2_{1:T}$ can, in fact, be arbitrarily large. The detailed comparison is provided in Appendix~\ref{app:separation}. Second, from an algorithmic perspective, an algorithm with a regret bound involving $\tilde{\sigma}^2_{1:T}$ typically involves an \emph{implicit} update, typically requires an \emph{implicit} update, which operates on the original function and is significantly more costly than standard first-order methods (see Remark~10 in~\cite{chen2024optimistic}). Third, achieving regret bounds scaling with $\sigma^2_{1:T}$ typically requires leveraging the \emph{Regret Bounded by Variation in Utilities (RVU)} property~\cite{syrgkanis2015fast}, which captures the regret to be bounded not only by the gradient variations but also an additional negative stability term. Formally, an algorithm satisfies the RVU property if its regret upper bound is in the form of $\sum_{t=1}^T \langle x^* - x_t, u_t \rangle \leq \alpha + \beta \sum_{t=1}^T \|u_t - u_{t-1}\|^2 - \gamma \sum_{t=1}^T \|x_t - x_{t-1}\|^2$, for some constants \(\alpha, \beta, \gamma > 0\). This structure enables finer control over the regret by explicitly analyzing trajectory stability, establishing profound connections to game theory~\cite{syrgkanis2015fast,zhang2022no} and accelerations in smooth optimization~\cite{LectureNote:AOpt25}. Consequently, the key challenge lies in how to achieve this preferred $\sigma^2_{1:T}$-scaling \emph{without} knowledge of $D$ and $G$ for unbounded domains.


% Importantly, the RVU framework allows the derivation of bounds in terms of $\sigma^2_{1:T}$, rather than the looser $\tilde{\sigma}^2_{1:T}$, thereby offering a more meaningful characterization of learning performance. We will elaborate on this point in  Remark~\ref{remark:dependence-u2} and~\ref{remark:adaptive}. Hence, the core difficulty lies in attaining this preferred $\sigma^2_{1:T}$-scaling \emph{without} knowledge of $D$ and $G$, especially for unbounded domains.


\section{Optimistic Online Newton Step (ONS) for the SEA Model}\label{sec:ons}
In this section, different from the Optimistic follow-the-regularized-leader (OFTRL)~\citep{sachs2023accelerated} and Optimistic mirror descent (OMD)~\citep{chen2024optimistic}, we first introduce the \underline{O}ptimistic \underline{O}nline \underline{N}ewton \underline{S}tep (OONS) algorithm as the base algorithm for the   ``parameter-free'' algorithms to be introduced later. This algorithm is summarized in Algorithm~\ref{alg:ONS}.

\begin{algorithm}[t]
\caption{Optimistic Online Newton Step (OONS)}
\label{alg:ONS}

    \textbf{Input:} learning rate $\eta_t>0$, $x'_1=0$.
    
    \begin{algorithmic}[1]
    \FOR{$t=1,\ldots, T$}
        \STATE Receive optimistic prediction $m_t$ and range hint $z_t$.
        \STATE Update $x_t = \argmin_{x\in\cX}\{\inner{x}{m_t} + D_{\psi_{t}}(x,x_t')\}$  where $\psi_t(x) = \frac{1}{2}\|x\|_{A_t}^2$ and $A_t = 4 z_1^2 I+ \sum_{s=1}^{t-1} \eta_s(\nabla_s \!-\! m_s)(\nabla_s \!-\! m_s)^\top \!+\! 4\eta_t z_t^2 I.$
        
        \STATE Receive ${g_t=\nabla f_t(x_t)}$ and construct ${\nabla_t = g_t + 32 \eta_t \inner{x_t}{g_t - m_t}(g_t - m_t)}$.
        \STATE Update $x_{t+1}' = \argmin_{x\in\cX}\{\inner{x}{\nabla_t} + D_{\psi_t}(x, x_t')\}$.
    \ENDFOR
\end{algorithmic}
\end{algorithm}


The ONS algorithm~\cite{hazan2007logarithmic} is an iterative algorithm that adaptively updates a second-order (Hessian-based) model of the loss, allowing more efficient gradient-based updates and improved regret bounds. \oons\ also maintains two sequences $\{x_t\}_{t=1}^T$ and $\{x'_t\}_{t=1}^T$ like OMD and OFTRL, which is achieved by introducing the optimistic prediction $m_t$. Chen et al.~\citet{chen2021impossible} also considered combining their Multi-scale Multiplicative-weight with Correction (MsMwC) algorithm with this variant of the ONS algorithm. However, the step-size $\eta$ is fixed in their algorithm and the MsMwC algorithm is applied to learn the optimal $\eta_\star$. Different from it, in \oons, we consider adaptive step-sizes $\eta_t$. 

\iffalse
At the start of each iteration, the algorithm receives an optimistic prediction $m_t$ and a range hint $z_t$, which satisfies $\|g_t-m_t\|_2\leq z_t$, where $g_t=\nabla f_t(x_t)$ and $m_t=\nabla f_{t-1}(x_{t-1})$. Define the   {\em surrogate loss}
$ c_t(x) = \inner{x}{g_t}+16\eta_t\inner{x}{g_t - m_t}^2$ and  $\nabla_t =\nabla c_t(x_t)=g_t + 32\eta_t\inner{x_t}{g_t-m_t}(g_t-m_t)$ which   can be considered as a second-order correction term $32\eta_t\inner{x_t}{g_t-m_t}(g_t-m_t)$ added to the vector $g_t$ in the update of $x'_{t+1}$ (Line 5). The introduction of the correction term ensures that the regret bound depends on $\sum_{t=1}^T\inner{u}{g_t-m_t}^2$ instead of $\sum_{t=1}^T\inner{u-x_t}{g_t-m_t}^2$. If the decision set $\cX$ is unbounded, i.e., $D=\infty$, it is difficult to obtain a desirable regret bound from $\sum_{t=1}^T\inner{u-x_t}{g_t-m_t}^2$. The term $\sum_{t=1}^T\inner{x_t}{g_t-m_t}^2$ can be canceled by the additional correction term. We now present an upper bound on the regret of \oons\ below.
\fi

\begin{theorem}\label{thm:thm1}
Suppose that $\|g_t-m_t\|_2\leq z_t$, $z_t$ is non-decreasing in $t$,   $64\eta_t D z_T\leq 1$ for all $t\in[T]$, and $\eta_t$ is non-increasing in $t$. Then, \oons\ guarantees that
\iffalse
\begin{align}\label{eq:thm1}
&\regret_T(u) \leq O\bigg(\frac{r\ln(T\eta_1 z_T/z_1)}{\eta_T}+z_1^2\|u\|_2^2+D(z_T-z_1)\nonumber\\
&\quad+\sum_{t=1}^T\eta_t\langle u,g_t-m_t\rangle^2-z_1^2\sum_{t=2}^T\|x_t-x_{t-1}\|_2^2\bigg),
\end{align}
\fi
\begin{equation}\label{eq:thm1}
\resizebox{\textwidth}{!}{$
\regret_T(u) \leq O\left(\frac{r\ln(T\eta_1 z_T/z_1)}{\eta_T}+z_1^2\|u\|_2^2 + D(z_T - z_1) + \sum_{t=1}^T \eta_t \langle u, g_t - m_t \rangle^2 - z_1^2 \sum_{t=2}^T \|x_t - x_{t-1}\|_2^2 \right)
$}
\end{equation}
where $r$ is the rank of $\sum_{t=1}^T(g_t-m_t)(g_t-m_t)^\top$.
\end{theorem}

Next, we verify that \oons\ also works for the case with known parameters $D, G$, and we can also obtain a similar regret bound as \citet{sachs2023accelerated} and \citet{chen2024optimistic}.
The regret bound of \oons\ for the SEA model with known parameters $D$ and $G$ is presented below. We specify the adaptive step-size for all $t\in[T]$ as 
\begin{align}
\eta_t=\min\bigg\{\frac{1}{64Dz_T}, \frac{1}{D\sqrt{\sum_{s=1}^{t-1}\|g_{s}-m_s\|_2^2}}\bigg\}.
\label{eqn:eta_t_ONS}
\end{align}
Since $G$ is known, we have $\|g_t-m_t\|_2\leq z_t=2G, \forall t\in[T]$ and $\eta_t$ is defined in terms of $z_T$ here.

\begin{theorem}\label{thm:knownDG}
Under Assumptions~\ref{ass:G},~\ref{ass:D},~\ref{ass:L}, and \ref{ass:convexF}, \oons\ with step-size $\eta_t$ given in \eqref{eqn:eta_t_ONS}, $m_t=\nabla f_{t-1}(x_{t-1})$ and $z_t=2G$ for all $t\in[T]$ ensures $\mathbb{E}[\regret_{T}(u)] =\tilde{O}(\sqrt{\sigma_{1:T}^2}+\sqrt{\Sigma_{1:T}^2})$.%  the following regret bound
% \begin{equation*}
% \mathbb{E}[\regret_{T}(u)] =\tilde{O}\Big(\sqrt{\sigma_{1:T}^2}+\sqrt{\Sigma_{1:T}^2}\Big).
% \end{equation*}
\end{theorem}

\begin{remark}
 Theorem~\ref{thm:knownDG} achieves the same (up to logarithmic terms)  dependence on 
$\sigma_{1:T}^2$ and $\Sigma_{1:T}^2$ as in \citet{sachs2023accelerated} and \citet{chen2024optimistic}. The primary reason to use \oons\ as the base algorithm instead of OMD~\cite{chen2024optimistic} is that the final regret bound for OMD typically depends on $D \sqrt{\sum_{t=1}^T \| g_t - m_t \|_2^2}$. In scenarios when $D$ is unknown or potentially infinite, like in Section~\ref{sec:fullyfree}, this might lead to $O(T)$ regret bounds. By contrast, \oons\ leverages adaptive second-order information, which helps remove (or substantially reduce) explicit dependence on $D$. In \eqref{eq:thm1}, the only term relevant to $D$ is $D(z_T - z_1)$, which solely depends on the starting and ending points, $z_1$ and $z_T$. %Since these points can be expressed in terms of $G$, which is independent of $T$. On the other hand, in the OMD bound, this part is expressed as $\sqrt{\sum_{t=1}^T z_t^2}$ with $z_t=\|g_t-m_t\|_2$, which scales with $T$.
\end{remark}

\section{Parameter-free Algorithms for the SEA Model}
In this section, we develop parameter-free algorithms for the SEA model, building on \oons\ which we use as the base algorithm. Moreover, we allow the decision set $\cX$ to be potentially unbounded throughout this section, i.e. $D=\infty$ in Assumption~\ref{ass:D}. In Section~\ref{sec:unknownD}, we present a \emph{comparator-adaptive} algorithm for the SEA model for unknown $D$ but known $G$, and then, we develop the \emph{comparator- and Lipschitz-adaptive} algorithm where both $D$ and $G$ are unknown in Section~\ref{sec:fullyfree}.

\subsection{Comparator-adaptive algorithm}\label{sec:unknownD}
We  now propose the \underline{C}omparator-\underline{A}daptive \underline{O}ptimistic \underline{O}nline \underline{N}ewton \underline{S}tep (\cawe) algorithm for the unknown $D$ (potentially infinite) but known $G$ case by using a meta-base algorithm framework. Recent works in \citet{chen2021impossible} and \citet{jacobsen2023unconstrained}  addressed the challenges associated with unbounded domains by developing  a base-learner framework. Building on this philosophy, we propose \cawe (Algorithm~\ref{alg:unknownD}) where we adopt the MsMwC--Master algorithm~\cite{chen2021impossible} as the meta algorithm.

\iffalse
\begin{algorithm}[t]
\caption{Comparator-adaptive algorithm for the SEA model (\cawe)}
\label{alg:unknownD}
\textbf{Input:}  Lipschitz constant $G$.

\begin{algorithmic}[1]
    \FOR{$t=1,\dotsc,T$}
    \STATE Create $N=\lceil\log T\rceil$ experts. Each expert $j\in[N]$ runs \oons\ with stepsize $\eta_t^j$ from \eqref{eq:PF-OONSwE eta stepsize} in the $D_j$-bounded domain.
    \STATE Each expert $j$ provides $x_t^j$.
    \STATE Run Algorithm~\ref{alg:meta} to obtain $w_t\in\Delta_N$.
    \STATE The final decision is $x_t=\sum_{j=1}^N w_{t,j} x_t^j$.
    \ENDFOR
\end{algorithmic}
\end{algorithm}

\begin{algorithm}[t]
\caption{Meta Algorithm}
\label{alg:meta}
\textbf{Input:}  Additional expert set $\cS$ defined in \eqref{eq:expertset}.\\
\textbf{Initialization:} $p'_1\in\Delta_{\cS}$ such that $p'_{1,k} \propto \beta_k^2$ for all $ k\in\cS$.
\begin{algorithmic}[1]
    \FOR{$t=1,\dotsc,T$}
    \STATE Construct $h_t\in\real^N$ with $h_t^j=\langle \nabla f_{t-1}(x_{t-1}^j), x_t^j \rangle$ and feed it to all MsMwC base algorithms. 
    \STATE Each  expert $k\in\cS$ runs MsMwC algorithm with step-size $\beta_{t,j}^k$ and plays $w_t^k\in\Delta_N$, where $\beta_{t,j}^k=2\beta_k$ for all $t$ and $j$ and $\beta_k=\frac{1}{32\cdot 2^k}$.
    \STATE Receive $w_t^k$ for each $k\in\cS$ and define $H_{t}^k=\inner{w_t^k}{h_t}$.
    \STATE Compute $p_t=\argmin_{p\in\Delta_{\cS}}\inner{p}{H_t}+D_{\phi}(p,p'_t)$, where $ \phi(p)=\sum_{k\in\cS}\frac{p_k}{\beta_k}\ln p_k$.
    
    \STATE Play $w_t=\sum_{k\in\cS}p_{t,k}w_t^k\in\Delta_N$.
    \STATE Receive $\ell_t\in\real^N$ with $\ell_t^j=\langle\nabla f_{t}(x_{t}^j) , x_t^j \rangle$ and define $L_{t}^k=\inner{w_t^k}{\ell_t}$. 
    \STATE Compute 
    $p'_{t+1}=\argmin_{p\in\Delta_{\cS}}\inner{p}{L_t+b_t}+D_{\phi}(p,p'_t)$, where $b_t^k=32\beta_k(L_t^k-H_t^k)^2$.
    \ENDFOR
\end{algorithmic}
\end{algorithm}
\fi

%\setlength{\textfloatsep}{5pt} % Reduce space between floats and text

\begin{algorithm}[!t]
\caption{Comparator-adaptive algorithm for the SEA model (\cawe)}
\label{alg:unknownD}
\textbf{Input:} Lipschitz constant $G$.
\begin{algorithmic}[1]
    \FOR{$t=1,\dotsc,T$}
    \STATE Create $N=\lceil\log T\rceil$ base-learners. Each base-learner $j\in[N]$ runs \oons\ with stepsize $\eta_t^j$.
    \STATE Each base-learner $j$ provides $x_t^j$.
    \STATE Run Algorithm~\ref{alg:meta} to obtain $w_t\in\Delta_N$.
    \STATE The final decision is $x_t=\sum_{j=1}^N w_{t,j} x_t^j$.
    \ENDFOR
\end{algorithmic}
\end{algorithm}

\begin{algorithm}[!t]
\caption{Meta Algorithm}
\label{alg:meta}
\textbf{Input:} Additional expert set $\cS$ defined in \eqref{eq:expertset}.\\
\textbf{Initialization:} $p'_1\in\Delta_{\cS}$ such that $p'_{1,k} \propto \beta_k^2$ for all $ k\in\cS$.
\begin{algorithmic}[1]
    \FOR{$t=1,\dotsc,T$}
    \STATE Construct $h_t\in\real^N$ with $h_t^j=\langle \nabla f_{t-1}(x_{t-1}^j), x_t^j \rangle$.
    \STATE Each expert $k\in\cS$ runs MsMwC with step-size $\beta_{t,j}^k$ and plays $w_t^k\in\Delta_N$.
    \STATE Receive $w_t^k$ for each $k\in\cS$ and compute $H_{t}^k=\inner{w_t^k}{h_t}$.
    \STATE Compute $p_t=\argmin_{p\in\Delta_{\cS}}\inner{p}{H_t}+D_{\phi}(p,p'_t)$.
    \STATE Play $w_t=\sum_{k\in\cS}p_{t,k}w_t^k\in\Delta_N$.
    \STATE Receive $\ell_t\in\real^N$. Define $L_{t}^k=\inner{w_t^k}{\ell_t}$ and $b_t^k=32\beta_k(L_{t}^k-H_{t}^k)$.
    \STATE Compute $p'_{t+1}=\argmin_{p\in\Delta_{\cS}}\inner{p}{L_t+b_t}+D_{\phi}(p,p'_t)$.
    \ENDFOR
\end{algorithmic}
\end{algorithm}


The algorithm uses $N$ base-learners. For \emph{any} base-learner $i \in [N]$, the regret can be decomposed as
\begin{align*}
&\regret_T(u)=\sum_{t=1}^T f_t(x_t) - \sum_{t=1}^T f_t(u)= \underbrace{\sum_{t=1}^T f_t(x_t) - \sum_{t=1}^T f_t(x_t^i)}_{\text{Meta Regret}} + \underbrace{\sum_{t=1}^T f_t(x_t^i) - \sum_{t=1}^T f_t(u)}_{\text{Base Regret}}~.
\end{align*}
Let $\cA_i$ denote the base algorithm for the $i$-th base-learner. We denote  $\regret_T^{\cA_i}(u)=\sum_{t=1}^T f_t(x_t^i) - \sum_{t=1}^T f_t(u)$ as the base regret by taking $\cA_i$ as the base algorithm.  Moreover, the final decision $x_{t}$ is a weighted-average of all the base-learners' decisions: $
x_{t} = \sum_{j=1}^N w_{t,j} x_{t}^j$ with $w_t\in\Delta_N$.
As such,
\begin{equation}\label{eq:splitreg}
\regret_T(u)\leq\regret_T^{\cA_i}(u)+\sum_{t=1}^T\langle \ell_t, w_t-w_\star^i\rangle,
\end{equation}
where $\ell_t\in\mathbb{R}^N$ with $\ell_t^j=\langle \nabla f_t(x_t^j),x_t^j\rangle$ and $w_{\star}^i$ is a  vector in $\Delta_N$ whose $j$-th component is $(w_{\star}^i)_j=1$ if $j=i$ and $0$ otherwise.  Refer to Appendix~\ref{app:splitreg} for the proof of~\eqref{eq:splitreg}.

We first consider the base algorithm. Specifically, for each base-learner $ j \in [N] $, we impose a constraint that it operates within   $\cX_j = \{ x : \|x\|_2 \leq D_j \;\text{and}\; x \in \cX \} $
%(--- $\cX_j = \{ x : \|x\|_2 \leq D_j \wedge x \in \cX \} $, Note: We provide a proof for the case when $\cX = \real^d$, but it can easily be extended to the case when $\cX \subset \real^d$ by changing the definition of $\cX_j$ by including an intersection with $\cX$ ---)
, where $D_j = 2^j$. Then, we define $g_t^j=\nabla f_{t}(x_{t}^j)$ and $m_t^j=\nabla f_{t-1}(x_{t-1}^j)$. Each base-learner $j\in[N]$ runs \oons\ with step-size 
\begin{align}
\label{eq:PF-OONSwE eta stepsize}
\eta_t^j=\min\bigg\{\frac{1}{64D_jz_T},\frac{1}{D_j\sqrt{\sum_{s=1}^{t-1}\|g_{s}^j-m_s^j\|_2^2}}\bigg\},
\end{align}
which depends on $D_j$ instead of $D$ in \oons. Hence, each base-learner $j$ can update $x_t^j$ via \oons\ with step-size $\eta_t^j$. Since the final decision is  $x_t=\sum_{j=1}^N w_{t,j} x_t^j$, we need to adopt a meta-algorithm to learn the weight parameter $w_t\in \Delta_N$.

As mentioned above, we introduce a constraint that each base-learner $j\in[N]$ operates within a $D_j$-bounded domain. We can consider this as a ``multi-scale'' base-learner problem~\cite{bubeck2017online,cutkosky2018black,chen2021impossible} where each base-learner $j$ has a different loss range such that $|\ell_{t}^j|\leq GD_j$ since $\ell_{t}^j= \langle\nabla f_t(x_t^j),x_t^j\rangle$. We choose the Multi-scale Multiplicative-weight with Correction (MsMwC)--Master algorithm (Algorithm 2 in \citet{chen2021impossible}) as the meta-algorithm to learn $w_t$, which is implemented based on the MsMwC~\cite{chen2021impossible}. Details of MsMwC are presented in Appendix~\ref{app:msmwc}.%MsMwC~\cite{chen2021impossible} enhances the classic multiplicative-weights approach by incorporating individualized, time-varying step-sizes to address multi-scale problems. Details of MsMwC are presented in Appendix~\ref{app:msmwc}. The MsMwC--Master algorithm extends this by dynamically combining multiple MsMwC instances with different step-sizes into a unified strategy that adapts to the best-performing configuration. In our case, while the regret bound of MsMwC depends on $\max_{j \in [N]} D_j = 2^N = O(T)$ with $N = \lceil \log T \rceil$, we employ the MsMwC--Master algorithm to refine this dependence to $D_i$ by learning the optimal step-size. 
Specifically, we define a new expert set 
\iffalse
\begin{align}\label{eq:expertset}
\cE=\Big\{(\beta_k,\cM_k):\forall k\in \cS,\beta_k=\frac{1}{32\cdot 2^{k}}\Big\},
\end{align}
\fi
\begin{align}\label{eq:expertset}
\cS=\{k\in \mathbb{Z}:\exists j\in[N], GD_j \!\leq\! 2^{k-2}\!\leq\! GD_j\sqrt{T}\}.
\end{align}

For all $k\in\cS$, the step-size of the MsMwC--Master algorithm is set to $\beta_k=\frac{1}{32\cdot 2^k}$. Each expert $k\in\cS$ runs the MsMwC algorithm with $w'_1$ being uniform over $\cZ(k)$, where $\cZ_k=\{j\in[N]:GD_j\leq 2^{k-2}\}$. Moreover, each base MsMwC algorithm only works in the subset $\cZ(k)$, i.e., $w_t\in \Delta_N$ with $w_{t,j}=0$ for all $j\notin \cZ(k)$. We can view \cawe\ (Algorithm~\ref{alg:unknownD}) as a three-layer structure, where the meta-algorithm (Algorithm~\ref{alg:meta}) itself consists of two layers, which we refer to as \emph{meta top} and \emph{meta middle}. Also, the base layer is \oons\ algorithm. For clarity, we summarize the notations of the three-layer hierarchy for \cawe~in Table~\ref{tab:cawe}.

\iffalse
For each iteration $t\in[T]$, we construct the prediction $h_t\in\real^N=\langle \nabla f_{t-1}(x_{t-1}^j),x_t^j\rangle$ and feed it to all MsMwC base algorithms. After creating additional experts as \eqref{eq:expertset}, each expert $k\in\cS$ runs MsMwC  and plays $w_t^k\in\Delta_N$, where we set $\beta_{t,j}^k=2\beta_k$ for all $t\in [T]$ and $j\in[N]$ in the MsMwC base algorithm. Next, we collect $w_t^k$ for each $k\in\cS$ and define $H_{t,k}=\inner{w_t^k}{h_t}$. Then, we proceed to update $p_t, p'_{t+1}$ via $H_t, L_t$ and the regularizer $\phi(p)$ (Lines 8--12 in Algorithm~\ref{alg:unknownD}) and play $w_t=\sum_{k\in\cS}p_{t,k}w_t^k\in\Delta_N$. Finally, we obtain the final decision as $x_t=\sum_{j=1}^N w_{t,j} x_t^j$. For clarity, we summarize the notations of the three-layer hierarchy for \cawe~in Table~\ref{tab:cawe}.
\fi

{\small 
\begin{table}[t]
\centering
\caption{Three-layer hierarchy of \cawe}
\vspace{1em}
\renewcommand{\arraystretch}{1.4} 
\begin{tabularx}{\textwidth}{c|c|c|c|c|X}
\hline

\hline
\textbf{Layer} & \textbf{Algorithm} & \textbf{Loss} & \textbf{Optimism} & \textbf{Decision} & $\quad\;\;$ \textbf{Output} \\
\hline
\textbf{Top Meta} 
& MsMwC 
&  \((L_t^k)_{k \in \cS}\) 
& \((H_t^k)_{k \in \cS}\) 
& \(p_t \in \Delta_{\cS}\) 
& \(w_t\!=\!\sum_{k} p_{t,k} w_t^k\) \\
\hline
\textbf{Middle Meta} 
& MsMwC 
& \(\ell_t \in \mathbb{R}^N\) 
& \(h_t \in \mathbb{R}^N\)
& \((w_t^k)_{k\in\cS} \in \Delta_N\)
& \(\qquad\quad  w_t^k\) \\
\hline
\textbf{Base} 
& OONS 
& \(\nabla f_t(x_t^j)\) 
& \(\nabla f_{t-1}(x_{t-1}^j)\) 
& \((x_t^j)_{j\in N} \in \mathcal{X}_j\) 
& \(\qquad\quad  x_t^j\) \\
\hline

\hline
\end{tabularx}
\label{tab:cawe}
\end{table}}


In the following, we provide the expected regret guarantee for \cawe.
\begin{theorem}\label{thm:unknownD}
Let $D$ be unknown (potentially infinite). Under Assumptions \ref{ass:G}, \ref{ass:L} and \ref{ass:convexF}, \cawe\ provides the following regret
\begin{equation}\label{eq:thmunknownD}
\mathbb{E}[\regret_T(u)]=\tilde{O}\Big(\|u\|_2^2+\|u\|_2(\sqrt{\sigma_{1:T}^2}+\sqrt{\Sigma_{1:T}^2})\Big).
\end{equation}
\end{theorem}
\iffalse
The proof  considers %is accomplished by considering 
the following three cases.
\begin{enumerate}
\item $\|u\|_2< D_1=2$: We take $i=1$ here and substitute $D_i$ with $D_1=2$.
\item Let $i$ be the smallest integer such that  $\|u\|_2\leq D_i=2^i$. We have $\|u\|_2\leq D_i\leq 2\|u\|_2$ since $D_{i+1}=2D_i$. Then, we can substitute $D_i$ with $2\|u\|_2$.
\item $\|u\|_2>D_{\max}=2^N$: We have
\begin{align*}
\sum_{t=1}^T f_t(x_t)-f_t(u)&\leq \sum_{t=1}^T \inner{\nabla f_t(x_t)}{x_t-u}\\
&\leq 2GT\|u\|_2.
\end{align*}
We take $N=\lceil\log T\rceil$, then $T\leq \|u\|_2$. Therefore,
\begin{align*}
\sum_{t=1}^T f_t(x_t)-f_t(u)\leq 2G\|u\|_2^2.
\end{align*}
\end{enumerate}
\fi
This regret guarantee is referred to as ``comparator-adaptive'' because it depends directly on the norm of the comparator, $\|u\|_2$, rather than explicitly relying on the diameter of the decision set, $D$. Notably, when considering the constrained decision set with a diameter $D$, our regret bound immediately recovers the result $\mathbb{E}[\regret_T(u)]=\tilde{O}(D^2+D(\sqrt{\sigma_{1:T}^2}+\sqrt{\Sigma_{1:T}^2}))$ established in~\citet{sachs2023accelerated,chen2024optimistic}. 


% \begin{remark}
% In \citet{sachs2023accelerated} and \citet{chen2021impossible}, if we retain the dependence on   $D$, the regret bound is  $\mathbb{E}[\regret_T(u)]=\tilde{O}(D^2+D(\sqrt{\sigma_{1:T}^2}+\sqrt{\Sigma_{1:T}^2}))$.
% \iffalse
% \begin{equation*}
% \mathbb{E}[\regret_T(u)]=\tilde{O}\Big(D^2+D(\sqrt{\sigma_{1:T}^2}+\sqrt{\Sigma_{1:T}^2})\Big).
% \end{equation*}
% \fi
% For the ``parameter-free'' algorithm, we cannot set the step-size dependent on $D$ and the term such as $\|u-x_t\|_2$ cannot be bounded by $D$ in the regret bound when $D$ is potentially infinite. However, our result in Theorem~\ref{thm:unknownD} demonstrates that the same regret bound can be achieved even when 
% $D$ is unknown or potentially infinite. Instead of depending explicitly on $D$, our bound depends directly on the norm of the comparator as \eqref{eq:thmunknownD}.
% \end{remark}

% \begin{remark}[]
% \label{remark:offline}

One limitation of our regret bound~\eqref{eq:thmunknownD} is that when particularizing to adversarial OCO, it achieves only an $\tilde{O}(\|u\|_2^2 + \|u\|_2 \sqrt{T})$ worst-case regret bound, which falls short of the best-known $\tilde{O}(\|u\|_2 \sqrt{T})$ regret bound~\cite{orabona2016coin, jacobsen2022parameter, mhammedi2020lipschitz}. However, in the following two remarks, we will justify the $\|u\|_2$-dependence in the gradient-variation regret and emphasize the fundamental challenge of achieving adaptivity from the gradient-variation bound (for smooth functions) to the worst-case bound (for the non-smooth case) when the decision set of online learning is \emph{unconstrained}. 

\begin{remark}[Dependency on $\|u\|_2^2$] 
\label{remark:dependence-u2}
Recent studies have established the connection between gradient-variation online learning and accelerated offline optimization through advanced online-to-batch conversions~\cite{LectureNote:AOpt25,NeurIPS'25:Holder}. Specifically, let $d_0 = \|x_0 - x_* \|_2$ denote the distance of an initial point $x_0$ to the optimum $x_*$. For an $L$-smooth function, gradient-variation online algorithms using first-order information correspond to an accelerated convergence rate of $O(L d_0^2/T^2)$ via the \emph{stabilized} online-to-batch conversion~\cite{kavis2019unixgrad}. For a $G$-Lipschitz function, the problem-independent regret bounds translate to an $O(G d_0/\sqrt{T})$ rate through the standard conversion~\cite{cesa2006prediction}. In this context, we hypothesize that the $\|u\|_2^2$ term may be unavoidable in gradient-variation regret for unconstrained online learning, paralleling how the $d_0^2$ term also appears in the accelerated rate of unconstrained offline optimization. %In this context, we conjecture that the $\|u\|_2^2$ term might be unavoidable in the gradient-variation regret for unconstrained online learning, as the $d_0^2$ term also appears in the accelerated rate of unconstrained offline optimization. 
\end{remark}

\begin{remark}[Adaptivity between gradient-variation bound and worst-case bound]\label{remark:adaptive}
We argue that achieving \emph{adaptivity} between the gradient-variation bound and the problem-independent worst-case bound in \emph{unconstrained} online learning may be as challenging as achieving \emph{universality} in offline optimization over unconstrained domains, where the method must adapt to both smooth and Lipschitz functions. To the best of our knowledge, the best-known universal method for offline unconstrained optimization is by~\cite{kreisler2024accelerated}, which combines \textsc{UnixGrad}~\cite{kavis2019unixgrad} with the \textsc{DoG} step size~\cite{ivgi2023dog}. Nonetheless, this method is complex and still relies on a predefined range of parameters, highlighting both the difficulty of the problem and the fact that it remains only partially solved. Consequently, designing a \emph{single} unconstrained online learning algorithm that adaptively bridges the gradient-variation regret bound for smooth functions and the worst-case bound for Lipschitz functions is non-trivial, which could provide new insights into universal offline optimization methods. We leave this for future work.
\end{remark}
% Therefore, in our view, this adaptivity is as hard as realizing the ``universality" in smooth and non-smooth optimization in unconstrained domain. 
    % Recent studies have drawn a close connection between online learning and offline optimization through online-to-batch conversion techniques. In particular, gradient-variation-based online learning methods—which leverage smoothness assumptions when using first-order algorithms—correspond to smooth offline optimization settings via \emph{weighted} online-to-batch conversions~\cite{LectureNote:AOpt25}. On the other hand, worst-case online learning bounds, typically relying on Lipschitz continuity, correspond to \emph{non-smooth} offline optimization through the standard online-to-batch framework~\cite{cesa2006prediction}. This perspective suggests that an online algorithm capable of interpolating between gradient-variation and worst-case regret bounds implicitly yields an \emph{adaptive} offline optimization method. Such a method would achieve an accelerated convergence rate in smooth settings (e.g., $O(1/T^2)$) and the standard $O(1/\sqrt{T})$ rate in non-smooth cases. However, achieving such adaptivity in the offline setting remains highly non-trivial. The best known parameter-free algorithm that obtains accelerated rates in smooth convex optimization is the recent work from~\cite{kreisler2024accelerated}, which introduces the \textsc{UnixGrad}~\cite{kavis2019unixgrad} combined with the \textsc{DoG}~\cite{ivgi2023dog}. This method represents an advancement in adaptive accelerated optimization and highlights the complexity of bridging smooth and non-smooth regimes in a unified, parameter-free manner.

% \end{remark}

\begin{remark}[On dependence on the time horizon]\label{remark:anytime}
The use of $T$ in \cawe\ (via $N=\lceil \log T\rceil$ experts) is only for convenience and not fundamental. An \emph{anytime} variant is obtained by the standard doubling trick: restart the algorithm at epochs of lengths $1,2,4,\dots$, and in epoch $k$ set $N_k=\lceil \log 2^{k-1}\rceil$. This introduces at most an additional logarithmic factor already hidden in $\tilde{O}(\cdot)$~\citep[Section 4.3]{zhao2021bandit}. A restart-free alternative is a sleeping (awakening) expert grid of learning rates as in the multi–rate construction of \citet{mhammedi2019lipschitz}, which activates only those experts whose scale becomes relevant. 
\end{remark}

\subsection{Comparator- and Lipschitz-adaptive Algorithm}\label{sec:fullyfree}
The algorithm in the previous subsection requires \emph{prior} knowledge of the Lipschitz constant $G$. Due to practical limitations such knowledge may not be available in real applications. A comparator- and Lipschitz-adaptive algorithm would instead {\em adapt} to an unknown Lipschitz constant $G$. %is considered in this subsection. %while maintaining the desired expected regret bound of the form $\mathbb{E}[\regret_T(u)]= \tilde{O}\big(\|u\|_2(\sqrt{\sigma_{1:T}^2}+\sqrt{\Sigma_{1:T}^2})\big)$.

\iffalse
A simple approach to this challenge, proposed by \citet{cutkosky2019artificial}, relies on a gradient-clipping reduction. The method begins by designing an algorithm $\cA$ that achieves appropriate regret when provided with prescient ``hints'' $h_t$ that satisfy $h_t \geq \|g_t\|$ at the start of round $t$. 
\iffalse
Since we cannot practically provide such hints (as $g_t$ has not yet been observed), we instead utilize our best estimate, $h_t = \max_{s<t} \|g_s\|_2$. The algorithm then processes \emph{clipped} subgradients $\tilde{g}_t = g_t \min\{1, \frac{h_t}{\|g_t\|_2}\}$, ensuring that $h_t \geq \|\tilde{g}_t\|_2$, thus maintaining the validity of hints provided to $\cA$. 

The outputs $x_t$ of $\cA$ are constrained to lie within the domain $\cX_t = \{x \in \real^d : \|x\|_2 \leq \sqrt{\sum_{s=1}^{t-1} \|g_s\|_2/G_s}\}$, where $G_t = \max_{s \leq t} \|g_s\|_2$.
\fi
\iffalse
This approach \cite{cutkosky2019artificial} is demonstrated to guarantee the following regret bound 
\begin{align*}
\!\!\regret_T(u) \!\leq \!\regret_T^\cA(u) \!+\! G_T\bigg(\!\|u\|_2\! +\! \sqrt{\sum_{t=1}^T \frac{\|g_t\|_2}{G_t}} \!+\! \|u\|_2^3\!\bigg),\!
\end{align*}
where $\regret_T^\cA(u)$ represents the regret of $\cA$ on the losses $\tilde{g}_t$. 
\fi
Since we cannot practically provide such hints (as $g_t$ has not yet been observed), we instead utilize our best estimate available at this stage, the clipped gradient by adopting the technique from~\cite{cutkosky2019artificial}. Inspired by this approach, we start with an initial guess $B_0$ on the range of $\max_{t}\|g_t-m_t\|_2$, where $g_t=\nabla f_t(x_t)$. We denote $B_t=\max_{0\leq s\leq t}\|g_s-m_s\|_2$ as the range of predicted error up to iteration $t$. Then, we feed the following truncated gradient to the algorithm in each iteration as $\tilde{g}_t=m_t+\frac{B_{t-1}}{B_t}(g_t-m_t)$.
\iffalse
\begin{equation}
\tilde{g}_t=m_t+\frac{B_{t-1}}{B_t}(g_t-m_t).
\end{equation}
\fi
By noting that $\|\tilde{g}_t-m_t\|_2\leq B_{t-1}$, the truncated gradient allows the learner to assume that the range of predicted error in iteration $t$ is known at the beginning of this iteration.
\fi

A simple approach to handling the unknown gradient norms, proposed by \citet{cutkosky2019artificial}, relies on a gradient-clipping reduction. The key idea is to design an algorithm $\cA$ that achieves appropriate regret when given prescient ``hints'' $h_t \geq \|g_t\|_2$ at the start of round $t$. Since such hints are impractical (as $g_t$ is not observed beforehand), we instead approximate them using a clipped gradient, inspired by~\cite{cutkosky2019artificial}. We start with an initial guess $B_0$ on the range of $\max_{t} \|g_t - m_t\|_2$, where $g_t=\nabla f_t(x_t)$. We define $B_t=\max_{0 \leq s \leq t} \|g_s - m_s\|_2$ as the predicted error range up to iteration $t$. The truncated gradient is then defined as $\tilde{g}_t = m_t + \frac{B_{t-1}}{B_t} (g_t - m_t)$. The truncated gradient satisfies $\|\tilde{g}_t - m_t\|_2 \leq B_{t-1}$, allowing the learner to assume that the range of predicted error in iteration $t$ is known at the start.

Next, we initialize the decision set diameter guess as $D_1=1$. For each iteration $t\in[T]$, we first play $x_t$ and receive $g_t=\nabla f_t(x_t)$. To update $D_t$, we consider the condition  $
D_t < \sqrt{\sum_{s=1}^t\frac{\|g_s \|_2}{\max\{1,\max_{k \leq s} \| g_k \|_2 \}}}$. If this condition holds, we update $D_{t+1}$ using the doubling trick.
\iffalse
\begin{equation}\label{eq:updateD}
D_{t+1} = 2\sqrt{\sum_{s=1}^t\frac{\|g_s \|_2}{\max\{1,\max_{k \leq s} \| g_k \|_2 \}}}.
\end{equation}
\fi
This ensures that %the number of updates to $D_t$ is at most $O(\log T)$. Thus, 
we need to update $D_t$ a maximum of $M=O(\log T)$ times. %Similarly, we update $C_t$ a maximum of $O(\log T)$ times.
%Let $M$ be the number of total updates in $D_t$, where $M = O(\log T)$. 
We divide the total $T$ iterations into disjoint subsets of $M$ iterations. If the ``doubling'' occurs at the $t$-th iteration, we update $t_a\leftarrow t$ and reset $x'_{t+1}=0$ and the matrix $A_{t+1}$ in \oons\ as follows
\begin{align}
\label{eq:updatedAt}
A_{t+1} &= 4z_{t_a+1}^2 I \!+\! \sum_{s=t_a+1}^{t} \!\eta_s(\nabla_s - m_s)(\nabla_s - m_s)^\top+4\eta_{t+1}z_{t+1}^2 I.
\end{align}
Then, we feed $\tilde{g}_t$ to \oons\ running in the $D_{t+1}$-bounded domain $\cX_{t+1}=\{x:\|x\|_2\leq D_{t+1}\wedge x \in \cX\}$ and obtain $x_{t+1}$. %We take $z_{t+1}=B_t=\max_{0\leq s\leq t}\|g_s-m_s\|_2$. 
We summarize the ideas in Algorithm~\ref{alg:fullyfree} and term it as \underline{C}omparator and \underline{L}ipschitz-\underline{A}daptive \underline{O}ptimistic \underline{O}nline \underline{N}ewton \underline{S}tep (or \claoons) algorithm.






\begin{algorithm}[t]
\caption{Comparator and Lipschitz-Adaptive (or \claoons) for the SEA model}
\label{alg:fullyfree}

    \textbf{Input:} Initial scale $B_0$.

    \textbf{Initialize:} $D_1=1$.
    
    \begin{algorithmic}[1]
    \FOR{$t=1,\ldots, T$}
        \STATE Run \oons\ in $D_t$-bounded domain and obtain $x_t$. Play $x_t$ and receive $g_t=\nabla f_t(x_t)$.
        \STATE Construct $\tilde{g}_t=m_t+\frac{B_{t-1}}{B_t}(g_t-m_t)$, where $B_t=\max_{0\leq s\leq t}\|g_s-m_s\|_2$.
        \IF{$D_t < \sqrt{\sum_{s=1}^t\frac{\|g_s \|_2}{\max\{1,\max_{k \leq s} \| g_k \|_2\} }}$}
            \STATE Update $D_{t+1} = 2\sqrt{\sum_{s=1}^t\frac{\|g_s \|_2}{\max\{1,\max_{k \leq s} \| g_k \|_2 \}}}$ and reset $A_{t+1}$ as \eqref{eq:updatedAt} and $x'_{t+1}=0$.
        \ENDIF
        \STATE Feed $\tilde{g}_t$ to \oons\ running in the $D_{t+1}$-bounded domain and get $x_{t+1}$, where  $z_{t+1} = B_t$. %= \max_{0 \leq s \leq t} \| g_s - m_s $\|_2$.
    \ENDFOR
\end{algorithmic}
\end{algorithm}

%We present the following regret bound of \fpfons\ for the SEA model as follows.
\begin{theorem}\label{thm:fullyfree}
Let both $D$ (potentially infinite) and $G$ be unknown. Under Assumptions~\ref{ass:G} (but $G$ is unknown),~\ref{ass:L} and~\ref{ass:convexF}, the proposed \fpfons\ algorithm satisfies  %algorithm provides the following regret
\begin{align*}
\mathbb{E}[\regret_{T}(u)]\leq\tilde{O}\Big(\|u\|_2^2(\sqrt{\sigma_{1:T}^2}+\sqrt{\Sigma_{1:T}^2})+G^2\|u\|_2^2+\|u\|_2^4+G\|u\|_2^3+G^2\sqrt{\sigma_{1:T} + \mathfrak{G}_{1:T}}\Big), 
\end{align*}
where $\sigma_{1:T}$ captures the stochastic gradient deviation (without the squared norm) and $\mathfrak{G}_{1:T}$ denotes the sum of maximum expected gradients.
\end{theorem}

\iffalse
\textbf{Proof Sketch.} The maximum number of updates in \fpfons\ is $M=O(\log T)$. We split the total $T$ iterations into $M$ intervals $I_m$ with $m\in[M]$, where the last iteration of $I_m$ (denoted by $t_m$) either equals to $T$ or $D_{t_m+1}\neq D_{t_m}$. The regret can be written as
\begin{align*}
\regret_T(u)&=\sum_{t=1}^T f_t(x_t) - \sum_{t=1}^T f_t(u) \leq \sum_{t=1}^T \inner{g_t}{x_t - u} \\
&= \sum_{m=1}^M \underbrace{\sum_{t \in I_m} \inner{g_t}{x_t - u_t}}_{T_m} + \underbrace{\sum_{t =1}^T \inner{g_t}{u_t - u}}_{T_{\text{extra}}},
\end{align*}
where we define $u_t = \min\{ 1, \frac{D_t}{\| u \|_2} \} u$. We first write $T_m$ as
\begin{align*}
    %\sum_{t \in I_m} \inner{g_t}{x_t - u_t} 
    T_m&= \sum_{t \in I_m} \inner{\tilde{g}_t}{x_t - u_t}+ \sum_{t \in I_m} \inner{g_t - \tilde{g}_t}{x_t - u_t}.
\end{align*}
Note that within interval $I_m$, the domain has a bounded diameter $D_m$. When $t\in I_m$, we take the step-size as
\begin{equation*}
\eta_t = \min\Big\{ \frac{1}{64 D_t z_t}, \frac{1}{\sqrt{\sum_{s=t_1}^{t-1}\| \tilde{g}_s - m_s \|_2^2}} \Big\},
\end{equation*}
where $t_1$ is first index in $I_m$. Then, we feed $\tilde{g}_t$ to \oons\ running in the $D_{t+1}$-bounded domain.%, the regret bound of $\sum_{t \in I_m} \inner{\tilde{g}_t}{x_t - u_t}$ can be obtained similar to the Theorem~\ref{thm:thm1}.  
The details of deriving the regret bound of $\sum_{t \in I_m} \inner{\tilde{g}_t}{x_t - u_t}$ are presented in Appendix~\ref{app:sec42}. Moreover, we have
\begin{align*}
    &\sum_{t \in I_m} \inner{g_t - \tilde{g}_t}{x_t - u_t} \leq (D_t + \| u \|_2 ) \sum_{t \in I_m} \| g_t - \tilde{g}_t \|\\
    &\leq (D_t + \| u \|_2 ) \sum_{t\in I_m} \frac{B_t - B_{t-1}}{B_t} \| g_t - m_t \|_2 \\
    &\leq 2(D_t + \| u \|_2 ) G.
\end{align*}


For  $T_{\text{extra}}$, we observe that $u_t$ is either $u$ or $\frac{D_t}{\| u \|_2} u$. When $u_t \ne u$,  $\| u \|_2 \geq D_t > \sqrt{\sum_{s=1}^t \frac{\| g_s\|_2}{ \max_{k \leq s} \| g_k\|_2}}$. Once $u_t = u$, it stays there. Let $t^*$ be the final round when $u_t \ne u$. Then, we have $T_{\text{extra}} = \sum_{t = 1}^{t^*} \inner{g_t}{u_t - u}$ and 
\begin{align*}
   T_{\text{extra}}% \sum_{t = 1}^T \inner{g_t}{u_t - u} 
   %&= \sum_{t = 1}^{t^*} \inner{g_t}{u_t - u}\\
    &\leq 2 \| u \|_2 \sum_{t = 1}^{t^* - 1} \| g_t \|_2 \!+\! 2 \| u \|_2 G\leq 2 \| u \|_2^3 G \!+\! 2 \| u \|_2 G.
\end{align*}
\fi

\begin{remark}[Discussion and challenges]
%Theorem~\ref{thm:unknownD} in Section~\ref{sec:unknownD} assumes that the diameter $D$ of the decision set is unknown (and potentially infinite) but known $G$, while Theorem~\ref{thm:fullyfree} goes further by assuming that both $D$ and $G$ are unknown. Theorem~\ref{thm:fullyfree} thus provides a more general guarantee without requiring any additional expert oracles or prior knowledge of parameters.
In Theorem~\ref{thm:fullyfree}, the regret includes $\|u\|_2^2(\sqrt{\sigma_{1:T}^2}+\sqrt{\Sigma_{1:T}^2})$. Ideally, we aim to achieve a dependence of $\tilde{O}(\|u\|_2)$, consistent with  \cite{cutkosky2019artificial} and \cite{mcmahan2010adaptive}. However, achieving this within the SEA framework presents significant challenges. As mentioned in Section~\ref{sec:setup}, obtaining regret bounds that scale with $\sigma^2_{1:T}$ in the SEA framework is  difficult. These challenges are compounded in the comparator and Lipschitz-adaptive setting. Below, we outline some of the main technical challenges associated with achieving the desired bound of $\tilde{O}(\|u\|_2(\sqrt{\sigma_{1:T}^2}+\sqrt{\Sigma_{1:T}^2}))$.

%Analyzing methods such as those proposed by~\cite{cutkosky2019artificial,jacobsen2022parameter,jacobsen2023unconstrained} become particularly challenging within the SEA framework. It remains unclear whether a direct extension of these methods to the SEA framework is feasible. The primary difficulties lie in effectively incorporating optimism and devising a tractable analysis of Centered Mirror Descent under the SEA framework, especially with a specific choice of the $\psi$ function.

As stated in Section~\ref{sec:setup}, the methods such as those proposed by~\cite{cutkosky2019artificial,jacobsen2022parameter,jacobsen2023unconstrained} cannot be applied to obtain the expected regret bound in terms of $\sigma^2_{1:T}$. The work~\cite{cutkosky2019combining} also looks promising; however, it remains unclear to us whether the approach proposed in the paper can be directly extended to the SEA framework. Their results, presented in Theorems 3 and 5 of \cite{cutkosky2019combining}, are not Lipschitz-adaptive. Specifically, they operate under the assumptions that $\|g_t\|_2 \leq 1$ and $\|m_t\|_2 \leq 1$. %Also, we cannot derive the expected regret bound in terms of $\sigma^2_{1:T}$.

One could also  use a large number of base-learners to achieve a regret  of  $\tilde{O}(\|u\|_2 \sqrt{r \sum_t \|g_t-m_t\|_2^2}+\|u\|_2^3)$, similar to~\cite{chen2021impossible}. However, this approach presents a subtle yet significant challenge. Following a similar analysis~\cite{chen2021impossible}, we get the following decomposition: $\sum_t \inner{g_t}{x_t-u} =\sum_t \langle g_t,x_t-x_t^{k_*}\rangle+\sum_t \langle g_t, x_t^{k_*}-u\rangle$. By leveraging Theorem 23 in~\cite{chen2021impossible}, we can write $\sum_t \langle g_t, x_t^{k_*}-u\rangle$ as $\sum_t\langle g_t, x_t^{k_*}-u\rangle\leq \tilde{O}(\|u\| \sqrt{r \sum_t \|g_t-m_t\|_2^2}-\sum_t \|x_t^{k_*}-x_{t-1}^{k_*}\|_2^2)$. Observe that expressing the first term, $\sqrt{r \sum_t \|g_t-m_t\|_2^2}$, in terms of $\sigma_{1:T}$ and  $\Sigma_{1:T}$ introduces additional terms involving $\sum_t \|x_t-x_{t-1}\|_2^2$ (See Lemma 2 in~\cite{chen2024optimistic}). The only way to address this term is through the negative term $-\sum_t \|x_t^{k_*}-x_{t-1}^{k_*}\|_2^2$, which becomes tricky. This challenge is reminiscent of the problem encountered by~\cite{zhao2024adaptivity}. Their solution, as outlined in Equation (17) of ~\cite{zhao2024adaptivity}, relies on the bounded domain assumption, which is not applicable in our setting. Consequently, this limitation prevents us from improving the $\|u\|_2^2$ dependence in the leading term by using additional base-learners.


 

The bound in Theorem~\ref{thm:fullyfree} also includes an additive term involving $\sqrt{\mathfrak{G}_{1:T}}$, which reflects the sum of the maximum expected gradients'  norms over $T$ rounds, and arises because the domain is potentially unbounded. Note that this term does not have a $\|u\|$ dependence. Hence,  the comparator having a large norm in an unbounded setting (potentially dependent on~$T$) does not affect its growth.  In the worst case, $\sqrt{\mathfrak{G}_{1:T}}=O(\sqrt{T})$, which underscores that this additive term does not have a significant adverse effect on the regret as $\sqrt{\sigma_{1:T}^2}$ and $\sqrt{\Sigma_{1:T}^2}$ also scale as $\sqrt{T}$~\citep{sachs2023accelerated}.
\end{remark}

%\begin{remark}

%\end{remark}
\section{Conclusions and Future Work}
This paper presents novel parameter-free algorithms for the SEA model, addressing critical challenges in online optimization where traditional approaches require prior knowledge of parameters such as the diameter of the domain $D$ and the Lipschitz constant of the loss functions $G$. Our proposed algorithms: \cawe\ and \claoons\ are designed to operate effectively even when $D$ and $G$ are unknown, demonstrating their adaptability and practicality.

There are several avenues for future research. First, we would like to  improve the regret's dependence on   $\| u \|_2$    when both  $D$  and  $G$  are unknown. %This could be achieved by designing novel algorithms or a more refined analysis of existing ones. 
Another promising direction is to reduce the number of gradient queries in \cawe\ from  $O(\log T)$  to  $O(1)$,  thus enhancing its efficiency. An intriguing question in the comparator-adaptive setting is whether it is possible to design a single, simple online algorithm that simultaneously achieves two types of bounds: $\tilde{O}(\|u\|_2^2 + \|u\|_2(\sqrt{\sigma^2_{1:T}} + \sqrt{\Sigma^2_{1:T}}))$ and $\tilde{O}(\|u\|_2 G \sqrt{T})$. As discussed in Remark~\ref{remark:adaptive}, it remains an open challenge to construct an adaptive parameter-free online algorithm that can interpolate between these bounds. %Such an algorithm would enable adaptive parameter-free offline optimization, achieving both the accelerated convergence rate in smooth settings ($O(1/T^2)$) and the standard $O(1/\sqrt{T})$ rate in non-smooth cases, leveraging the weighted online-to-batch conversion from this ideal online parameter-free algorithm.
An additional open direction is to move beyond \emph{expected} regret and derive \emph{high-probability} (or variance-sensitive) regret guarantees for the SEA model in the parameter-free setting. Current SEA analyses, including \citet{sachs2023accelerated,chen2024optimistic}, bound only $\E[\regret_T(u)]$; developing concentration results that retain the fine $\sigma_{1:T}^2$ and $\Sigma_{1:T}^2$ dependence without incurring suboptimal logarithmic inflation appears non-trivial and is left for future work.
%Besides, another extension would involve adapting our parameter-free framework to strongly convex or exp-concave loss functions. Moreover, another exciting direction is to leverage insights from parameter-free algorithms for the SEA model to guide the design of tuning-free SGD-type algorithms in the presence of adversarial corruptions.

\begin{ack}
This research is funded by the Singapore Ministry of Education
Academic Research Fund Tier 2 under grant number A-8000423-00-00 and
three Singapore Ministry of Education Academic Research Funds Tier 1 under
grant numbers A-8000189-01-00, A-8000980-00-00 and A-8002934-00-00.
\end{ack}


\newpage
\bibliography{references}
\bibliographystyle{unsrt}


%%%%%%%%%%%%%%%%%%%%%%%%%%%%%%%%%%%%%%%%%%%%%%%%%%%%%%%%%%%%

\newpage
\appendix
%\onecolumn



\section{Separation between $\sigma_{1:T}^2$ and $\tilde{\sigma}_{1:T}^2$}
\label{app:separation}

Recall (Section~\ref{sec:setup}) the two stochastic variance measures:
\[
\sigma_t^2=\sup_{x\in\cX}\E\big[\|\nabla f_t(x)-\nabla F_t(x)\|_2^2\big], \qquad
\tilde{\sigma}_t^2=\E\Big[\sup_{x\in\cX}\|\nabla f_t(x)-\nabla F_t(x)\|_2^2\Big],
\]
and their cumulative versions $\sigma_{1:T}^2=\sum_{t=1}^T \sigma_t^2$, $\tilde{\sigma}_{1:T}^2=\sum_{t=1}^T \tilde{\sigma}_t^2$. Always $\sigma_t^2 \le \tilde{\sigma}_t^2$ by Jensen Inequality, but the gap can be \emph{arbitrarily large}. The following simple 1-dimensional construction (with a growing number of disjoint regions) shows an order difference.

\begin{proposition}[Linear separation between $\sigma_{1:T}^2$ and $\tilde{\sigma}_{1:T}^2$]
Let $\cX=\bigcup_{i=1}^n X_i\subset\R$ with disjoint cells $X_i=[i-1,i)$ for $i\in\{1,\ldots,n\}$.
For each round $t\in\{1,\ldots,T\}$ and each cell $i\in\{1,\ldots,n\}$, draw independently
\begin{equation*}
c_{t,i}\sim\mathrm{Bernoulli}\!\left(\tfrac{1}{n}\right)
\quad\text{and}\quad
s_{t,i}\in\{-1,+1\}\ \text{ with }\ \Pr(s_{t,i}=1)=\Pr(s_{t,i}=-1)=\tfrac12,
\end{equation*}
and define the (stochastic) gradient field by
\begin{equation*}
\nabla f_t(x)=c_{t,i}\,s_{t,i}\qquad\text{for all }x\in X_i,
\end{equation*}
with arbitrary values on cell boundaries. Let $F_t$ be the expected loss, defined up to an additive constant by
$\nabla F_t(x)=\E[\nabla f_t(x)]$ (we fix the constant so that $F_t\equiv 0$).
Consider the two gradient–noise proxies
\begin{equation*}
\sigma_t^2 \;:=\; \sup_{x\in\cX}\E\!\left[\bigl\|\nabla f_t(x)-\nabla F_t(x)\bigr\|^2\right], 
\qquad
\tilde\sigma_t^2 \;:=\; \E\!\left[\sup_{x\in\cX}\bigl\|\nabla f_t(x)-\nabla F_t(x)\bigr\|^2\right],
\end{equation*}
and their sums $\sigma_{1:T}^2=\sum_{t=1}^T\sigma_t^2$ and $\tilde\sigma_{1:T}^2=\sum_{t=1}^T\tilde\sigma_t^2$.

Then, for every $n,T\ge 1$,
\begin{equation*}
\sigma_t^2=\frac{1}{n}
\quad\text{and}\quad
\tilde\sigma_t^2 = 1-\Bigl(1-\tfrac1n\Bigr)^{n}\xrightarrow[n\to\infty]{} 1-\tfrac{1}{e},
\end{equation*}
hence, for all large $n$,
\begin{equation*}
\sigma_{1:T}^2=\frac{T}{n}=\Theta(1)
\quad\text{and}\quad
\tilde\sigma_{1:T}^2 = T\!\left(1-\Bigl(1-\tfrac1n\Bigr)^n\right)=\Theta(T).
\end{equation*}
In particular, taking $n=n(T)=T$ yields $\sigma_{1:T}^2=1$ while $\tilde\sigma_{1:T}^2\ge (1-e^{-1})T$, i.e., a linear (in $T$) separation between the two quantities.
\end{proposition}

\begin{proof}
By construction and independence, for any $x\in X_i$,
\begin{equation*}
\nabla F_t(x)=\E[\nabla f_t(x)] = \E[c_{t,i}]\,\E[s_{t,i}] = \tfrac1n\cdot 0 = 0,
\end{equation*}
so $F_t\equiv 0$ (up to an additive constant). Therefore
\begin{equation*}
\sigma_t^2
=\sup_{x\in\cX}\E\big[(\nabla f_t(x))^2\big]
=\sup_{i}\E[c_{t,i}^2 s_{t,i}^2]
=\sup_{i}\E[c_{t,i}]
=\tfrac{1}{n}.
\end{equation*}
Summing over $t$ gives $\sigma_{1:T}^2=T/n$.

For $\tilde\sigma_t^2$, note that $s_{t,i}^2\equiv 1$ and $(\nabla f_t(x))^2=c_{t,i}$ for all $x\in X_i$. Thus
\begin{equation*}
\tilde\sigma_t^2
=\E\Big[\sup_{x\in\cX}(\nabla f_t(x))^2\Big]
=\E\Big[\max_{i} c_{t,i}\Big]
=1-\Pr\!\left(\forall i,\;c_{t,i}=0\right)
=1-\Bigl(1-\tfrac1n\Bigr)^{n}.
\end{equation*}
Since $\bigl(1-\tfrac1n\bigr)^n\le e^{-1}$ for all $n$, we have $\tilde\sigma_t^2\in[\,1-e^{-1},1]$.
Summing over $t$ yields $\tilde\sigma_{1:T}^2=T\bigl(1-(1-\tfrac1n)^n\bigr)\ge (1-e^{-1})T$.
Finally, with $n=T$ we obtain $\sigma_{1:T}^2=1$ and $\tilde\sigma_{1:T}^2\ge (1-e^{-1})T$, which proves the claimed linear separation.
\end{proof}

\paragraph{Remark.}
This example shows that regret bounds expressed in terms of $\tilde\sigma_{1:T}^2=\sum_t \E[\sup_x \|\cdot\|^2]$ can be a factor $\Theta(T)$ looser than bounds in terms of $\sigma_{1:T}^2=\sum_t \sup_x \E[\|\cdot\|^2]$ on the same instance.



\iffalse
\paragraph{Construction.} Let $\cX=\bigcup_{i=1}^n X_i$ where $X_i=[i-1,i)$ (disjoint “cells”). For each round $t$:
\begin{itemize}
\item The expected loss is identically zero: $F_t(x)\equiv 0$, hence $\nabla F_t(x)=0$.
\item For each cell $i\in[n]$, draw independently:
\[
c_{t,i}\in\{0,1\},\quad \Pr(c_{t,i}=1)=\tfrac{1}{n},\qquad
s_{t,i}\in\{-1,1\},\quad \Pr(s_{t,i}=1)=\Pr(s_{t,i}=-1)=\tfrac{1}{2}.
\]
\item For $x\in X_i$, define the stochastic gradient $\nabla f_t(x)=c_{t,i}\, s_{t,i}$ (a scalar sign or $0$). Thus $f_t$ is linear on each $X_i$ with slope $c_{t,i}s_{t,i}$.
\end{itemize}

\paragraph{Evaluating $\sigma_t^2$.} Since $\nabla F_t(x)=0$,
\[
\sigma_t^2=\sup_{x\in\cX}\E\big[(\nabla f_t(x))^2\big]=\sup_{i\in[n]}\E[c_{t,i}^2 s_{t,i}^2]
= \sup_{i\in[n]} \Pr(c_{t,i}=1)=\frac{1}{n}.
\]
Hence $\sigma_{1:T}^2 = T/n$.

\paragraph{Evaluating $\tilde{\sigma}_t^2$.} We have
\[
\tilde{\sigma}_t^2=\E\Big[\sup_{x\in\cX} (\nabla f_t(x))^2\Big]=\E\Big[\max_{i\in[n]} c_{t,i}\Big]
=1-\Pr(\forall i, c_{t,i}=0)=1-(1-1/n)^n \xrightarrow[n\to\infty]{} 1-\tfrac{1}{e}.
\]
So for all large $n$, $\tilde{\sigma}_t^2=\Theta(1)$ and $\tilde{\sigma}_{1:T}^2=\Theta(T)$.

\paragraph{Gap.} Taking $n=n(T)\to\infty$ (even $n=T$ suffices),
\begin{equation*}
\sigma_{1:T}^2 = \Theta(T/n)=\Theta(1),\qquad \tilde{\sigma}_{1:T}^2=\Theta(T),
\end{equation*}
a \emph{linear} (in $T$) separation. Thus any regret bound expressed in terms of $\tilde{\sigma}_{1:T}^2$ can be a factor $\Theta(T)$ larger than one stated in terms of $\sigma_{1:T}^2$ for this instance. This illustrates the practical sharpness advantage of achieving $\sigma_{1:T}^2$-based bounds (as in our results) rather than the looser $\tilde{\sigma}_{1:T}^2$ dependence.
\fi


\section{Omitted Details of Section~\ref{sec:setup}}
\subsection{Proof of~\eqref{eq:tildesigma}}\label{app:tildesigma}
\begin{proof}
    From Theorem~5 in~\cite{jacobsen2022parameter}, we have 
    \begin{align*}
        \regret_T(u) &\leq \tilde{O}\left(\|u\|_2 \sqrt{\sum_{t=1}^T \|\nabla f_t(x_t) - \nabla f_{t-1}(x_t)\|_2^2}\right)\\
        &\leq \tilde{O}\left(\|u\|_2 \sqrt{\sum_{t=1}^T \sup_{x\in \cX}\|\nabla f_t(x) - \nabla f_{t-1}(x)\|_2^2}\right)
    \end{align*}

    From Lemma~8 in \cite{chen2024optimistic}, we also have
    \begin{align*}
        &\sum_{t=1}^T \sup_{x\in\cX}\|\nabla f_t(x) - \nabla f_{t-1}(x)\|_2^2\\&\leq G^2+6\sum_{t=1}^T \sup_{x\in\cX}\|\nabla f_t(x) - \nabla F_{t}(x)\|_2^2 + 4\sum_{t=1}^T \sup_{x\in\cX}\|\nabla F_t(x) - \nabla F_{t-1}(x)\|_2^2.
    \end{align*}

    Taking expectations with Jensen's inequality and the definition of $\tilde{\sigma}_{1:T}^2$ and $\Sigma_{1:T}^2$, we obtain 
    \begin{equation*}
        \mathbb{E}[\regret_T(u)]\leq\tilde{O}\left(\|u\|_2 (\sqrt{\tilde{\sigma}_{1:T}^2} + \sqrt{\Sigma_{1:T}^2} )\right).\qedhere
    \end{equation*}
\end{proof}


\section{Omitted Details of Section~\ref{sec:ons}}

\subsection{Auxiliary Lemmas}
\begin{lemma}[Bregman Proximal Inequality]
\label{lem:3points}
The Bregman Proximal update in the form of $x_{t+1}=\argmin_{x\in\cX}\{\inner{x}{g_t}+D_{\psi}(x,x_t)\}$ satisfies
\begin{equation}\label{eq:bregmanprox}
\inner{g_t}{x_{t+1}-u}\leq D_{\psi}(u,x_t)-D_{\psi}(u,x_{t+1})-D_{\psi}(x_t,x_{t+1}).
\end{equation}
\end{lemma}
\begin{proof}
By the first-order optimality condition at $x_{t+1}$, for any $u\in\cX$, we have
\begin{equation*}
\inner{g_t+\nabla\psi(x_{t+1})-\nabla\psi(x_t)}{u-x_{t+1}}\geq 0,
\end{equation*}
On the RHS of \eqref{eq:bregmanprox}, we expand each term by the definition of Bregman divergence
\begin{equation*}
D_{\psi}(u,x_t)-D_{\psi}(u,x_{t+1})-D_{\psi}(x_t,x_{t+1})=\inner{\nabla\psi(x_{t+1})-\nabla\psi(x_{t})}{u-x_{t+1}}.
\end{equation*}
Hence, the proof is finished by rearranging the terms.
\end{proof}



\begin{lemma}\label{lem:regretexpand}
Let $x_t = \argmin_{x\in\cX}\{\inner{x}{m_t} + D_{\psi_{t}}(x,x_t')\}$ and $x_{t+1}' = \argmin_{x\in\cX}\{\inner{x}{g_t} + D_{\psi_t}(x, x_t')\}$. Then, it holds for any $u$ in $\cX$
\begin{align*}
\sum_{t=1}^T\langle x_t-u,g_t\rangle\leq \sum_{t=1}^T \langle x_t-x'_{t+1},g_t-m_t\rangle+D_{\psi_t}(u,x'_t)-D_{\psi_t}(u,x'_{t+1})
\\-D_{\psi_t}(x_t,x'_{t+1})-D_{\psi_t}(x_t,x'_{t}).
\end{align*}
\end{lemma}
\begin{proof}
We have 
\begin{equation}\label{eq:3points}
\langle x_t-u,g_t\rangle = \langle x_t-x'_{t+1},m_t\rangle+\langle x'_{t+1}-u,g_t\rangle+\langle x_t-x'_{t+1},g_t-m_t\rangle.
\end{equation}
We apply Lemma~\ref{lem:3points}  twice, i.e., $\inner{a-u}{f}\leq D_{\psi}(u,b)-D_{\psi}(u,a)-D_{\psi}(a,b)$ since $a=\argmin_{x\in\cX}\inner{x}{f}+D_{\psi}(x,b)$. Then, we have
\begin{align*}
\inner{x_t-x'_{t+1}}{m_t}&\leq D_{\psi_t}(x'_{t+1},x'_t)-D_{\psi_t}(x'_{t+1},x_t)-D_{\psi_t}(x_{t},x'_t),\\
\inner{x'_{t+1}-u}{g_t}&\leq D_{\psi_t}(u,x'_t)-D_{\psi_t}(x'_{t+1},u)-D_{\psi_t}(x'_{t},x'_{t+1}).
\end{align*}
Substitute these two back to \eqref{eq:3points} and sum over $T$ providing the desired result.
\end{proof}

\begin{lemma}\label{lem:term1}
In \oons, we have 
$0\leq\langle x_{t}-x'_{t+1},\nabla_t-m_t\rangle \leq 2\|\nabla_t-m_t\|_{A_t^{-1}}^2$ and $\sum_{t=1}^T\langle x_{t}-x'_{t+1},\nabla_t-m_t\rangle\leq O\Big(\frac{r\ln(T\eta_1 z_T/z_1)}{\eta_T}\Big)$.
\end{lemma}
\begin{proof}
We define 
\begin{equation*}
F_{m_t}(x)=\langle x,m_t\rangle+D_{\psi_t}(x,x'_t), \qquad F_{\nabla_t}(x)=\langle x,\nabla_t\rangle+D_{\psi_t}(x,x'_t)
\end{equation*}
By \oons, we have
\begin{equation*}
x_t=\argmin_{x\in\cX} F_{m_t}(x), \qquad x'_{t+1}=\argmin_{x\in\cX} F_{\nabla_t}(x).
\end{equation*}
Since $\nabla^2 D_{\psi_t}=A_t$ and by the first-order optimality at $x'_{t+1}$, we have
\begin{equation*}
F_{\nabla_t}(x_t)-F_{\nabla_t}(x'_{t+1})\geq \frac{1}{2}\|x_t-x'_{t+1}\|^2_{A_t}.
\end{equation*}
Also, we can write $F_{\nabla_t}(x_t)-F_{\nabla_t}(x'_{t+1})$ as
\begin{equation*}
F_{\nabla_t}(x_t)-F_{\nabla_t}(x'_{t+1})=\langle x_t-x'_{t+1},\nabla_t\rangle+D_{\psi_t}(x_t,x'_t)-D_{\psi_t}(x'_{t+1},x'_t).
\end{equation*}
Then,
\begin{align*}
\frac{1}{2}\|x_t-x'_{t+1}\|^2_{A_t}&\leq \langle x_t-x'_{t+1},\nabla_t-m_t\rangle+F_{m_t}(x_t)-F_{m_t}(x'_{t+1}),\\
&\leq \langle x_t-x'_{t+1},\nabla_t-m_t\rangle,
\end{align*}
where the second inequality comes from $x_t=\argmin F_{m_t}(x)$.
Therefore, we have
\begin{equation*}\label{eq:part1_1}
\langle x_{t}-x'_{t+1},\nabla_t-m_t\rangle \leq 2\|\nabla_t-m_t\|_{A_t^{-1}}^2.
\end{equation*}
Since $x'_{t+1}$ minimizes $F_{m_t}(x)$ and $x_{t}$ minimizes $F_{\nabla_t}(x)$, we have
\begin{align*}
0\leq F_{\nabla_t}(x_t)-F_{\nabla_t}(x'_{t+1})&=\langle x_t-x'_{t+1},\nabla_t\rangle+D_{\psi_t}(x_t,x'_t)-D_{\psi_t}(x'_{t+1},x'_t),\\
&=\langle x_t-x'_{t+1},\nabla_t-m_t\rangle+F_{m_t}(x_t)-F_{m_t}(x'_{t+1})\\
&\leq \langle x_t-x'_{t+1},\nabla_t-m_t\rangle.
\end{align*}

By the definition of $\nabla_t=g_t+32\eta_t\langle x_t,g_t-m_t\rangle(g_t-m_t)$, we have
\begin{align}
\|\nabla_t-m_t\|_2&=\|g_t-m_t+32\eta_t\langle x_t,g_t-m_t\rangle(g_t-m_t)\|_2\notag\\
&\leq \|g_t-m_t\|_2+32\eta_tD\|g_t-m_t\|_2^2\leq \frac{3}{2} \|g_t-m_t\|_2.\label{eq:norm of g_t - m_t and nabla_t - m_t}
\end{align}
Next, we define 
\begin{equation*}
\bar{A}_t=4z_1^2\cdot I+\sum_{s=1}^t \eta_s (\nabla_s-m_s)(\nabla_s-m_s)^\top.
\end{equation*}
Hence, $A_t\succeq \bar{A}_t$ since $\|\nabla_t-m_t\|_2^2\leq 4\|g_t-m_t\|_2^2\leq 4z_t^2$. Also, we have
\begin{equation*}
(\nabla_t-m_t)(\nabla_t-m_t)^\top=\frac{1}{\eta_t}[\eta_t(\nabla_t-m_t)(\nabla_t-m_t)^\top]=\frac{1}{\eta_t} (\bar{A}_t-\bar{A}_{t-1})
\end{equation*}
Then,
\begin{align*}
\sum_{t=1}^T \|\nabla_t-m_t\|_{A_t^{-1}}^2&\leq \sum_{t=1}^T \|\nabla_t-m_t\|_{\bar{A}_t^{-1}}^2,\\
&=\sum_{t=1}^T \tr\Big((\nabla_t-m_t)(\nabla_t-m_t)^\top \bar{A}_t^{-1}\Big),\\
&\leq \sum_{t=1}^T \frac{1}{\eta_t}\tr\Big(\bar{A}_t^{-1}(\bar{A}_t-\bar{A}_{t-1})\Big),\\
&\leq \sum_{t=1}^T \frac{1}{\eta_t}(\ln|\bar{A}_t|-\ln|\bar{A}_{t-1}|),\\
&\leq \frac{1}{\eta_T}\ln\frac{|\bar{A}_T|}{|\bar{A}_0|}.
\end{align*}
For $|\bar{A}_T|$:
\begin{align*}
|\bar{A}_T|&\leq |4z_1^2 I+\sum_{t=1}^T \eta_t(\nabla_t-m_t)(\nabla_t-m_t)^\top|\\
\ln |\bar{A}_T|&\leq O\Big(r\ln(1+\sum_{t=1}^T \frac{\eta_t}{4z_1^2}\|\nabla_t-m_t\|_2^2)\Big)\\
&\leq O\Big(r\ln(1+\frac{4z_T^2}{4z_1^2}\sum_{t=1}^T \eta_t)\Big)\\
&\leq O\Big(r\ln(1+\frac{\eta_1 z_T^2}{z_1^2}T)\Big),
\end{align*}
where $r$ is the rank of $\sum_{t=1}^T (\nabla_t-m_t)(\nabla_t-m_t)^\top$.

Therefore, we have
\begin{equation*}\label{eq:part1_final}
\sum_{t=1}^T\langle x_{t}-x'_{t+1},\nabla_t-m_t\rangle\leq O\Big(\frac{r\ln(T\eta_1 z_T/z_1)}{\eta_T}\Big).
\end{equation*}
\end{proof}

\begin{lemma}\label{lem:sumetat}
Let $s_t,\forall t\in[T]$ be non-negative. Then, 
\begin{equation*}
\sum_{t=1}^T \frac{s_t}{\sqrt{\sum_{j=1}^{t}s_j}}\leq 2\sqrt{\sum_{t=1}^T s_t}.
\end{equation*}
\end{lemma}
\begin{proof}
Let $S_t=\sum_{j=1}^t s_j$. Then,
\begin{equation*}
\sum_{t=1}^T \frac{s_t}{\sqrt{S_t}}\leq \sum_{t=1}^T\int_{S_{t-1}}^{S_t}\frac{1}{\sqrt{x}}dx=\int_{0}^{S_T} \frac{1}{\sqrt{x}}dx=2\sqrt{S_T}.
\end{equation*}
\end{proof}

\begin{lemma}[{Theorem 5 in \citet{sachs2023accelerated}, Lemma 3 in \citet{chen2024optimistic}}]
\label{lem:ex2sigma}
Under Assumptions \ref{ass:G} and \ref{ass:L}, we have
\begin{align*}
\sum_{t=1}^T \|\nabla f_t(x_t)-\nabla f_{t-1}(x_{t-1})\|_2^2&\leq G^2+4\sum_{t=2}^T\|\nabla F_t(x_{t-1})-\nabla F_{t-1}(x_{t-1})\|_2^2\\
&+8\sum_{t=1}^T\|\nabla f_t(x_t)-\nabla F_t(x_{t})\|_2^2+4L^2\sum_{t=2}^T\|x_t-x_{t-1}\|_2^2.
\end{align*}
\end{lemma}


\subsection{Proof of Theorem~\ref{thm:thm1}}
\begin{proof}
By Lemma~\ref{lem:3points}, we have
\begin{align*}
&\sum_{t=1}^T \langle x_t-u,\nabla_t\rangle\\
&\leq \sum_{t=1}^T \langle x_{t}-x'_{t+1},\nabla_t-m_t\rangle+D_{\psi_t}(u,x'_t)-D_{\psi_t}(u,x'_{t+1})-D_{\psi_t}(x_t,x'_{t+1})-D_{\psi_t}(x_t,x'_{t}),\\
&\leq \sum_{t=1}^T \langle x_{t}-x'_{t+1},\nabla_t-m_t\rangle+D_{\psi_1}(u,x'_1)+\sum_{t=1}^{T-1}D_{\psi_{t+1}}(u,x'_{t+1})-D_{\psi_t}(u,x'_{t+1})\\
&\qquad-\sum_{t=1}^T(D_{\psi_t}(x_t,x'_{t+1})+D_{\psi_t}(x_t,x'_{t}))
\end{align*}
\textbf{Term $D_{\psi_1}(u,x'_1)$}:
Since the initialization of $x'_1=0$ and $A_1=O(z_1^2 I)$, we have
\begin{equation*}
D_{\psi_1}(u,x'_1)=\frac{1}{2}\|u\|^2_{A_1}\leq O(z_1^2\|u\|_2^2).
\end{equation*}
\textbf{Term $\sum_{t=1}^{T-1}D_{\psi_{t+1}}(u,x'_{t+1})-D_{\psi_t}(u,x'_{t+1})$}:
First, we have
\begin{equation*}
A_{t+1}-A_t=\eta_t (\nabla_t-m_t)(\nabla_t-m_t)^\top+4\eta_t(z_{t+1}^2-z_t^2)I.
\end{equation*}
Also,
\begin{align*}
D_{\psi_{t+1}}(u,x'_{t+1})-D_{\psi_t}(u,x'_{t+1})&=\frac{1}{2}(u-x'_{t+1})^\top(A_{t+1}-A_t)(u-x'_{t+1})\\
&=\frac{1}{2}(u-x'_{t+1})^\top(\eta_t (\nabla_t-m_t)(\nabla_t-m_t)^\top+4\eta_t(z_{t+1}^2-z_t^2)I)(u-x'_{t+1})\\
&=\frac{\eta_t}{2}\inner{u-x'_{t+1}}{\nabla_t-m_t}^2+2\eta_t(z_{t+1}^2-z_t^2)\|u-x'_{t+1}\|_2^2.
\end{align*}
\iffalse
By the definition of $\nabla_t=g_t+32\eta_t\langle x_t,g_t-m_t\rangle(g_t-m_t)$, we have
\begin{align*}
\|\nabla_t-m_t\|_2&=\|g_t-m_t+32\eta_t\langle x_t,g_t-m_t\rangle(g_t-m_t)\|_2\\
&\leq \|g_t-m_t\|_2+32\eta_tD\|g_t-m_t\|_2^2\leq \frac{3}{2} \|g_t-m_t\|_2.
\end{align*}
\fi
Then, we have
\begin{align*}
&\sum_{t=1}^{T-1}D_{\psi_{t+1}}(u,x'_{t+1})-D_{\psi_t}(u,x'_{t+1})\\
&\leq \sum_{t=1}^{T-1} \frac{\eta_t}{2}\langle u-x'_{t+1},\nabla_t-m_t\rangle^2+O(\sum_{t=1}^{T-1}\eta_t D^2(z_{t+1}^2-z_t^2)),\\
& \leq \sum_{t=1}^{T-1} \frac{\eta_t}{2}\langle u-x'_{t+1},\nabla_t-m_t\rangle^2+O(\sum_{t=1}^{T-1}D(z^2_{t+1}-z^2_t)/z_T), \qquad \text{(since 
 } \eta_t \leq \frac{1}{64 D z_T}\text{)}\\
&\leq \sum_{t=1}^{T-1} \frac{\eta_t}{2}\langle u-x'_{t+1},\nabla_t-m_t\rangle^2+O(D(z_T-z_1)),\\
& \leq \sum_{t=1}^{T-1} \eta_t\langle u-x_{t},\nabla_t-m_t\rangle^2+\sum_{t=1}^{T-1}\eta_t\langle x_{t}-x'_{t+1},\nabla_t-m_t\rangle^2+O(D(z_T-z_1)),\\
&\leq \sum_{t=1}^{T-1} \eta_t\langle u-x_{t},\nabla_t-m_t\rangle^2+\sum_{t=1}^{T-1}\eta_t\langle x_{t}-x'_{t+1},\nabla_t-m_t\rangle\langle x_{t}-x'_{t+1},\nabla_t-m_t\rangle+O(D(z_T-z_1)),\\
&\leq \sum_{t=1}^{T-1} \eta_t\langle u-x_{t},\nabla_t-m_t\rangle^2+\sum_{t=1}^{T-1}2\eta_tDz_t\langle x_{t}-x'_{t+1},\nabla_t-m_t\rangle+O(D(z_T-z_1)), \qquad \text{(Using Equation~\ref{eq:norm of g_t - m_t and nabla_t - m_t})}\\
&\leq \sum_{t=1}^{T-1} \eta_t\langle u-x_{t},\nabla_t-m_t\rangle^2+O\Big(\frac{r\ln(T\eta_1 z_T/z_1)}{\eta_T}+D(z_T-z_1)\Big),\qquad \text{(Using Lemma~\ref{lem:term1})}\\
&\leq \sum_{t=1}^{T-1} 2\eta_t\langle u,\nabla_t-m_t\rangle^2+2\eta_t\langle x_{t},\nabla_t-m_t\rangle^2+O\Big(\frac{r\ln(T\eta_1 z_T/z_1)}{\eta_T}+D(z_T-z_1)\Big),\\
&\leq \sum_{t=1}^{T-1} 8\eta_t\langle u,g_t-m_t\rangle^2+8\eta_t\langle x_{t},g_t-m_t\rangle^2+O\Big(\frac{r\ln(T\eta_1 z_T/z_1)}{\eta_T}+D(z_T-z_1)\Big). \qquad \text{(since $32\eta_t\inner{x_t}{g_t-m_t}\leq \frac{1}{2}$)}
\end{align*}
\textbf{Term $\sum_{t=1}^T(D_{\psi_t}(x_t,x'_{t+1})+D_{\psi_t}(x_t,x'_{t}))$}:
\begin{align*}
\sum_{t=1}^T(D_{\psi_t}(x_t,x'_{t+1})+D_{\psi_t}(x_t,x'_{t}))&=\sum_{t=1}^T\frac{1}{2}(\|x_t-x'_{t}\|_{A_t}^2+\|x'_{t+1}-x_t\|_{A_t}^2)\\
&\geq \frac{1}{2}\sum_{t=2}^T \|x_t-x'_{t}\|_{A_{t-1}}^2+\frac{1}{2}\sum_{t=2}^{T+1}\|x_{t-1}-x'_t\|_{A_{t-1}}^2\\
&\geq \frac{4z_1^2}{4}\sum_{t=2}^T\|x_t-x_{t-1}\|_2^2=z_1^2\sum_{t=2}^T\|x_t-x_{t-1}\|_2^2,
\end{align*}
where the last inequality comes from $A_t\succeq A_{t-1}\succeq 4z_1^2 I$.

Since $c_t(x)$ is convex in $x$, we have
\begin{align*}
&\sum_{t=1}^T c_t(x_t)-c_t(u)\\
&=\sum_{t=1}^T \langle x_t-u,g_t\rangle+16\eta_t\langle x_t,g_t-m_t\rangle^2-16\eta_t\langle u,g_t-m_t\rangle^2\leq \sum_{t=1}^T \langle \nabla_t,x_t-u\rangle.
\end{align*}
Therefore, the final regret bound here is
\begin{align*}
\regret_T(u)&\leq\sum_{t=1}^T\langle x_t-u,g_t\rangle\\
&\leq O\Big(\frac{r\ln(T\eta_1 z_T/z_1)}{\eta_T}+z_1^2\|u\|_2^2+D(z_T-z_1)+\sum_{t=1}^T\eta_t\langle u,g_t-m_t\rangle^2-z_1^2\sum_{t=2}^T\|x_t-x_{t-1}\|_2^2\Big).
\end{align*}

\end{proof}




\subsection{Proof of Theorem~\ref{thm:knownDG}}
\begin{proof}
In this case, we have $\|g_t-m_t\|_2\leq 2G,\forall t\in[T]$. Then, we set $z_t=2G,\forall t\in[T]$ and step-size $\eta_t$ as
\begin{align*}
\eta_t=\min\Big\{\frac{1}{64Dz_T}, \frac{1}{D\sqrt{\sum_{s=1}^{t-1}\|g_{s}-m_s\|_2^2}}\Big\}\quad \mbox{for all} \quad  t\in[T]. 
\end{align*}

By substituting $\eta_t$ and $z_t$ into \eqref{eq:thm1}, we have
\begin{equation*}
\regret_{T}(u)\leq \tilde{O}\Big(D\sqrt{\sum_{t=1}^{T-1}\|g_{t}-m_t\|_2^2}+G^2D^2+\sum_{t=1}^T \eta_t\inner{u}{g_t-m_t}^2-G^2\sum_{t=2}^T\|x_t-x_{t-1}\|_2^2\Big).
\end{equation*}
By Lemma~\ref{lem:sumetat}, we have 
\begin{equation*}
\sum_{t=1}^T \eta_t\inner{u}{g_t-m_t}^2\leq 2D\sqrt{\sum_{t=1}^{T}\|g_{t}-m_t\|_2^2}
\end{equation*}
Then,
\begin{equation*}
\regret_{T}(u)\leq \tilde{O}\Big(D\sqrt{\sum_{t=1}^{T}\|g_{t}-m_t\|_2^2}-G^2\|x_t-x_{t-1}\|_2^2\Big).
\end{equation*}
By Lemma~\ref{lem:ex2sigma} and the definition of $g_t$ and $m_t$, we have
\begin{align*}
&D\sqrt{\sum_{t=1}^{T}\|g_{t}-m_t\|_2^2}-G^2\sum_{t=2}^T\|x_t-x_{t-1}\|_2^2\\
&\leq DG+2D\sqrt{\sum_{t=2}^T\|\nabla F_t(x_{t-1})-\nabla F_{t-1}(x_{t-1})\|_2^2}+2\sqrt{2}D\sqrt{\sum_{t=1}^T\|\nabla f_t(x_t)-\nabla F_t(x_{t})\|_2^2}\\&\qquad+2LD\sqrt{\sum_{t=2}^T\|x_t-x_{t-1}\|_2^2}-G^2\sum_{t=2}^T\|x_t-x_{t-1}\|_2^2\\
&\leq DG+\frac{D^2L^2}{G^2}+2D\sqrt{\sum_{t=2}^T\|\nabla F_t(x_{t-1})-\nabla F_{t-1}(x_{t-1})\|_2^2}+2\sqrt{2}D\sqrt{\sum_{t=1}^T\|\nabla f_t(x_t)-\nabla F_t(x_{t})\|_2^2}
\end{align*}
Therefore, we have
\begin{equation*}
\mathbb{E}[\regret_T(u)]\leq \tilde{O}(\sqrt{\sigma_{1:T}^2}+\sqrt{\Sigma_{1:T}^2})
\end{equation*}
\end{proof}

\subsection{Computational Complexity of OONS}
The primary computational bottleneck in our proposed \oons\ algorithm (Algorithm~\ref{alg:ONS}) is the management of the $d \times d$ matrix $A_t$. A naive implementation would involve storing this dense matrix and performing a full matrix inversion at each step, leading to prohibitive costs in high-dimensional settings.

\begin{itemize}
    \item \textbf{Storage:} Storing the dense $d \times d$ matrix $A_t$ requires $O(d^2)$ memory.
    \item \textbf{Computation:} A naive matrix inversion $A_t^{-1}$ would cost $O(d^3)$ per step, and the subsequent matrix-vector products would cost $O(d^2)$.
\end{itemize}

In practice, this complexity can be significantly reduced. Since the matrix $A_t$ is constructed by a sum of outer products ($A_t = cI + \sum \eta_s v_s v_s^\top$), its inverse can be efficiently computed and updated at each step using the Sherman-Morrison-Woodbury formula. This reduces the update complexity from $O(d^3)$ to $O(d^2)$ per step.

However, an $O(d^2)$ complexity per step can still be prohibitive in high-dimensional scenarios. To address this, existing research has explored several techniques:
\begin{itemize}
    \item \textbf{Matrix Sketching:} This technique approximates the original $d \times d$ matrix $A_t$ with a much smaller "sketched" matrix, thereby significantly reducing both storage and computational requirements. For instance, Luo et al.~\citet{luo2016efficient} have successfully applied sketching to the Online Newton Step (ONS) algorithm, creating matrix-free updates that avoid direct manipulation of the high-dimensional matrix.
    \item \textbf{Sparsity:} If the gradient vectors are sparse across most iterations, specialized sparse data structures and algorithms can be utilized. This allows update computations to be performed much more efficiently, avoiding the full $O(d^2)$ cost of dense matrix-vector multiplications.
\end{itemize}
Integrating these high-dimensional adaptation techniques into our proposed algorithms for the SEA model and analyzing their theoretical guarantees is an interesting direction for future work.



\section{Omitted Details of Section~\ref{sec:unknownD}}\label{app:sec31}

\subsection{Multi-scale Multiplicative-weight with Correction (MsMwC)}\label{app:msmwc}

We rephrase the MsMwC algorithm~\cite{chen2021impossible} as the following Algorithm~\ref{alg:msmwc}.
\begin{algorithm}[!t]
\caption{Multi-scale Multiplicative-weight with Correction (MsMwC)}\label{alg:msmwc}
\textbf{Input:} $w'_1\in\Delta_N$.
\begin{algorithmic}[1]
\FOR{$t=1,\dotsc,T$}
\STATE Receive the prediction $h_t\in\real^N$.
\STATE Compute $w_t=\argmin_{w\in\Delta_N}\inner{w}{h_t}+D_{\phi}(w,w'_t)$, where $\phi_t(w)=\sum_{j=1}^N \frac{w_j}{\beta_{t,j}}\ln w_j$.
\STATE Play $w_t$, receive $\ell_t$ and construct correction term $a_t\in\real^N$ with $a_{t,j}=32\beta_{t,j}(\ell_t^j-m_t^j)^2$.
\STATE Compute $w'_{t+1}=\argmin_{w\in\Delta_N}\inner{w}{\ell_t+a_t}+D_{\phi}(w,w'_t)$.
\ENDFOR
\end{algorithmic}
\end{algorithm}

\subsection{Proof of equation \eqref{eq:splitreg}}\label{app:splitreg}
The final decision $x_{t}$ is a weighted-average of base-learners' decisions: $
x_{t} = \sum_{j=1}^N w_{t,j} x_{t}^j$.
Then,
\begin{align*}
\regret_T(u)&=\regret_T^{\cA_i}(u)+\sum_{t=1}^T f_t(x_t) - \sum_{t=1}^T f_t(x_t^i)\\
&=\regret_T^{\cA_i}(u)+\sum_{t=1}^T f_t(\sum_{j\in[N]}w_{t,j} x_t^j)-\sum_{t=1}^T f_t(x_t^i)\\
&\leq \regret_T^{\cA_i}(u)+\sum_{t=1}^T \sum_{j\in[N]}w_{t,j} f_t(x_t^j)-\sum_{t=1}^T f_t(x_t^i)\\
&=\regret_T^{\cA_i}(u)+\sum_{t=1}^T \sum_{j\in[N]} f_t(x_t^j)[w_{t,j}-\mathbbm{1}(j=i)]\\
&=\regret_T^{\cA_i}(u)+\sum_{t=1}^T\sum_{j\in[N]} (f_t(x_t^j) - f_t(0))[w_{t,j}-\mathbbm{1}(j=i)]\\
&\leq \regret_T^{\cA_i}(u)+\sum_{t=1}^T \sum_{j\in[N]} \langle \nabla f_t(x_t^j),x_t^j\rangle[w_{t,j}-\mathbbm{1}(j=i)]\\
&=\regret_T^{\cA_i}(u)+\sum_{t=1}^T\langle \ell_t, w_t-w_\star^i\rangle,
\end{align*}
where $\ell_t\in\mathbb{R}^N$ with $\ell_t^j=\langle \nabla f_t(x_t^j),x_t^j\rangle$ and $w_{\star}^i$ is a  vector in $\Delta_N$ whose $j$-th component is $(w_{\star}^i)_j=1$ if $j=i$ and $0$ otherwise.  

\subsection{Auxiliary Lemma}
\begin{lemma} (Theorem 6 in \cite{chen2021impossible})\label{lem:msmwc_master}
Suppose for all $t\in[T]$ and $j\in[N]$, $|\ell_t^j|\leq GD_j$ and $|h_t^j|\leq GD_j$, where $\ell_t^j=\langle \nabla f_t(x_t^j),x_t^j\rangle$ and $h_t^j=\langle \nabla f_{t-1}(x_{t-1}^j),x_t^j\rangle$. Define $\Gamma_j=\ln(\frac{NTD_j}{D_1})$ and the set $\cE$ as 
\begin{align*}
\cE=\Big\{(\beta_k,\cG_k):\forall k\in \cS,\beta_k=\frac{1}{32\cdot 2^{k}}\Big\},
\end{align*}
where $\cG_k$ is the MsMwC algorithm with  $w'_1$ being uniform over $\cZ(k)$, $\cS=\{k\in \mathbb{Z}:\exists j\in[N], GD_j\leq 2^{k-2}\leq GD_j\sqrt{T}\}$ and $\cZ_k=\{j\in[N]:GD_j\leq 2^{k-2}\}$.
We have the following regret bound
\begin{equation*}
\sum_{t=1}^T\langle \ell_t, w_t-w_\star^i\rangle\leq O\Big(D_i\Gamma_i+\sqrt{\Gamma_i\sum_{t=1}^T (\ell_{t}^i-h_t^i)^2}\Big).
\end{equation*}
\end{lemma}
\begin{proof}
The regret $\sum_{t=1}^T\langle \ell_t, w_t-w_\star^i\rangle$ can also be decomposed as
\begin{align*}
\sum_{t=1}^T\langle \ell_t, w_t-w_\star^i\rangle&=\sum_{t=1}^T \inner{\ell_t}{w_t^{k_\star}-w_\star^i}+\sum_{t=1}^T\inner{\ell_t}{w_t-w_t^{k_\star}}\\
&=\sum_{t=1}^T \inner{\ell_t}{w_t^{k_\star}-w_\star^i}+\sum_{t=1}^T\inner{L_t}{p_t-e_{k_\star}},
\end{align*}
where $e_{k_\star}$ is the $k_{\star}$-th standard basis vector and the second equality is from the definition of $L_t$. 

For any $i\in[N]$, there exists a $k_\star$ such that $\eta_{k_\star}\leq \min\Big\{\frac{1}{128GD_i},\sqrt{\frac{\Gamma_i}{\sum_{t=1}^T(\ell_t^i-h_t^i)^2}}\Big\}\leq 2\eta_{k_\star}$. By Lemma~1 and Theorem~4 in \cite{chen2021impossible}, we have
\begin{equation*}
\sum_{t=1}^T\langle \ell_t, w_t-w_\star^i\rangle\leq O\Big(D_i\Gamma_i+\sqrt{\Gamma_i\sum_{t=1}^T (\ell_{t}^i-h_t^i)^2}\Big).
\end{equation*}
\end{proof}

\iffalse
For each expert $i \in \{1, \ldots, N\}$, the base algorithm $\cA_i$ for the $i$-th expert is from \oons.  Also, we adopt the MsMwC-Master algorithm proposed in \citep{chen2021impossible} to learn $w_t$. We define a new expert set $\cE$ as 
\begin{equation}\label{eq:expertset}
\cE=\Big\{(\beta_k,\cG_k):\forall k\in \cS,\beta_k=\frac{1}{32\cdot 2^{k}}, \cG_k \;\text{is the MsMwC with}\; w'_1 \;\text{being uniform over}\; \cZ(k)\Big\},
\end{equation}
where $\cS=\{k\in \mathbb{Z}:\exists j\in[N], GD_j\leq 2^{k-2}\leq GD_j\sqrt{T}\}$ and $\cZ_k=\{j\in[N]:GD_j\leq 2^{k-2}\}$. This algorithm is summarized in Algorithm~\ref{alg:unknownD}.

\begin{algorithm}[!htb]
\caption{Multi-Scale experts--ONS for Unknown $D$ with known $G$}
\label{alg:unknownD}
\textbf{Input:}  Lipschitz constant $G$ and the set $\cE$ defined in \eqref{eq:expertset}.\\
\textbf{Initialization:} $p'_1\in\Delta_{\cE}$ such that $p'_{1,k}\propto \eta_k^2$ for each $k\in\cE$.
\begin{algorithmic}[1]
    \For{$t=1,\dotsc,T$}
    \State Create $N=O(\log T)$ experts. Each expert $i\in[N]$ works in $D_i$ bounded domain and each expert runs Algorithm~\ref{alg:ONS} with stepsize $\eta_t^i$ depending on $D_i$.
    \State Each expert $i$ provides $x_t^i$.
    \State Construct prediction $h_t\in\real^N$ with $h_t^i=\inner{\nabla f_{t-1}(x_{t-1}^i)}{x_t^i}$ and feed it to all base algorithms MsMwC. 
    \State Create additional experts using \eqref{eq:expertset} with learning rate $\beta_t^k\in\real^N, \forall k\in\cE$.
    \State Each  expert $k\in\cE$ runs MsMwC algorithm and plays $w_t^k\in\Delta_N$, where $\beta_{t,i}^k=2\beta_k$ for all $t$ and $i$.
    \State Receive $w_t^k$ for each $k\in\cS$ and define $H_{t,k}=\inner{w_t^k}{h_t}$.
    \State Compute
    \begin{equation}
    p_t=\argmin_{p\in\Delta_{\cE}}\inner{p}{H_t}+D_{\psi}(p,p'_t), \text{where}\; \psi(p)=\sum_{k\in\cE}\frac{p_k}{\beta_k}\ln p_k.
    \end{equation}
    \State Play $w_t=\sum_{k\in\cE}p_{t,k}w_t^k\in\Delta_N$.
    \State Receive $\ell_t\in\real^N$ with $\ell_t^i=\inner{\nabla f_{t}(x_{t}^i)}{x_t^i}$ and define $L_{t,k}=\inner{w_t^k}{\ell_t}$. Compute 
    \begin{equation}
    p'_{t+1}=\argmin_{p\in\Delta_{\cE}}\inner{p}{L_t+b_t}+D_{\psi}(p,p'_t), \text{where}\; b_t^k=32\eta_k(L_t^k-H_t^k)^2.
    \end{equation}
    \State The final decision is $x_t=\sum_{i=1}^N w_{t}^i x_t^i$.
    \EndFor
    
   
\end{algorithmic}
\end{algorithm}
\fi


\subsection{Proof of Theorem~\ref{thm:unknownD}}

\begin{proof}
We begin with considering the first and second cases that  $\|u\|_2\leq D_1$ and $\|u\|_2\leq D_i\leq 2\|u\|_2$.

Here, we define $g_t^j=\nabla f_{t}(x_{t}^j)$ and $m_t^j=\nabla f_{t-1}(x_{t-1}^j)$ for all $t\in[T]$ and $j\in [N]$. For the meta regret, we define $\ell_t\in\mathbb{R}^N$ with $\ell_t^j=\langle \nabla f_t(x_t^j),x_t^j\rangle$ and $h_t\in\mathbb{R}^N$ with $h_t^j=\langle \nabla f_{t-1}(x_{t-1}^j),x_t^j\rangle$. Then, $|\ell_t^j|\leq GD_j$ and $|h_t^j|\leq GD_j$ for all $t\in[T]$ and $j\in[N]$. By applying Lemma~\ref{lem:msmwc_master}, we have
\begin{equation*}
\sum_{t=1}^T\langle \ell_t, w_t-w_\star^i\rangle\leq O\Big(D_i\Gamma_i+\sqrt{\Gamma_i\sum_{t=1}^T (\ell_{t}^i-h_t^i)^2}\Big).
\end{equation*}
By the definition of $\ell_t^i$ and $h_t^i$, we have
\begin{align*}
\sum_{t=1}^T(\ell_{t}^{i}-h_{t}^{i})^2&\leq \sum_{t=1}^T\langle \nabla f_t(x_t^i)-\nabla f_{t-1}(x_{t-1}^i),x_t^i\rangle^2\\
&\leq D_i^2\sum_{t=1}^T\|\nabla f_t(x_t^i)-\nabla f_{t-1}(x_{t-1}^i)\|_2^2=D_i^2 \sum_{t=1}^T\|g_t^i-m_t^i\|_2^2.
\end{align*}
Hence,
\begin{equation}\label{eq:metapart}
\sum_{t=1}^T\langle \ell_t, w_t-w_\star^i\rangle\leq O\Big(D_i\Gamma_i+D_i\sqrt{\Gamma_i\sum_{t=1}^T \|g_{t}^i-m_t^i\|_2^2}\Big).
\end{equation}

Then, we investigate the expert regret part. Here, we set the step-size for the expert $i$ as
\begin{equation}\label{eq:stepsizeunD}
\eta_t^i=\min\Big\{\frac{1}{64D_iz_T},\frac{1}{D_i\sqrt{\sum_{s=1}^{t-1}\|g_{s}^i-m_s^i\|_2^2}}\Big\}.
\end{equation}
By substituting the step-size specified in \eqref{eq:stepsizeunD} to \eqref{eq:thm1}, we have
\begin{align}\label{eq:regbaseAi}
\regret_{T}^{\cA_i}(u)&\leq\tilde{O}\Big(D_i\sqrt{\sum_{t=1}^{T}\|g_{t}^i-m_t^i\|_2^2}+G^2\|u\|_2^2+\frac{\|u\|_2^2}{D_i}\sqrt{\sum_{t=1}^{T}\|g_{t}^i-m_t^i\|_2^2}-G^2\|x_t^i-x_{t-1}^i\|_2^2\Big)\nonumber\\
&\leq \tilde{O}\Big(D_i^2+D_i\sqrt{\sum_{t=1}^{T}\|g_{t}^i-m_t^i\|_2^2}-G^2\|x_t^i-x_{t-1}^i\|_2^2\Big).
\end{align}

Therefore, by combining \eqref{eq:metapart} and \eqref{eq:regbaseAi}, we have 
\begin{align*}
\regret_T(u)&\leq \tilde{O}\Big(D_i^2+D_i\sqrt{\sum_{t=1}^{T}\|g_{t}^i-m_t^i\|_2^2}-G^2\sum_{t=2}^T\|x_t^i-x_{t-1}^i\|_2^2\Big)
\end{align*}
Then, by applying Lemma~\ref{lem:ex2sigma}, we have
\begin{align*}
&D_i\sqrt{\sum_{t=1}^{T}\|g_{t}^i-m_t^i\|_2^2}-G^2\sum_{t=2}^T\|x_t^i-x_{t-1}^i\|_2^2\\
&\leq GD_i+\frac{D_i^2L^2}{G^2}+2D_i\sqrt{\sum_{t=2}^T\|\nabla F_t(x^i_{t-1})-\nabla F_{t-1}(x^i_{t-1})\|_2^2}+2\sqrt{2}D_i\sqrt{\sum_{t=1}^T\|\nabla f_t(x_t^i)-\nabla F_t(x_{t}^i)\|_2^2}.
\end{align*}

Since the expert $i$ runs \oons\ within the set $\cX_i=\{x:\|x\|_2\leq D_i\}\subseteq \cX$, we have
\begin{align*}
\sup_{x\in\cX_i}\E_{f_t \sim \cD_t}  \big[ \| \nabla f_t(x) - \nabla F_{t}(x) \|_2^2 \big]&\leq\sup_{x \in \cX} \E_{f_t \sim \cD_t}  \big[ \| \nabla f_t(x) - \nabla F_{t}(x) \|_2^2 \big]\\
\sup_{x\in\cX_i}   \| \nabla F_t(x) - \nabla F_{t-1}(x) \|_2^2 &\leq\sup_{x \in \cX}  \| \nabla F_t(x) - \nabla F_{t-1}(x) \|_2^2.
\end{align*}

Hence, 
\begin{equation}\label{eq:regretwithDi}
\mathbb{E}[\regret_T(u)]\leq \tilde{O}\Big(D_i^2+\frac{D_i^2L^2}{G^2}+D_i\sqrt{\sigma_{1:T}^2}+D_i\sqrt{\Sigma_{1:T}^2}\Big).
\end{equation}

\textbf{Case 1}: $\|u\|_2\leq D_1$. We take $i=1$ and substitute $D_i$ with $D_1$ into \eqref{eq:regretwithDi}. Hence, we have
\begin{equation*}
\mathbb{E}[\regret_T(u)]\leq \tilde{O}\Big(\sqrt{\sigma_{1:T}^2}+\sqrt{\Sigma_{1:T}^2}\Big).
\end{equation*}

\textbf{Case 2}:
In this case, let $i$ be the smallest integer such that  $\|u\|_2\leq D_i=2^i$. We have $\|u\|_2\leq D_i\leq 2\|u\|_2$ since $D_{i+1}=2D_i$. Then, we substitute $D_i$ with $2\|u\|_2$ into the regret bound \eqref{eq:regretwithDi}. Then,
\begin{equation*}
\mathbb{E}[\regret_T(u)]\leq \tilde{O}\Big(\|u\|_2^2+\|u\|_2(\sqrt{\sigma_{1:T}^2}+\sqrt{\Sigma_{1:T}^2})\Big).
\end{equation*}

\textbf{Case 3}: $\|u\|_2>D_{\max}$.

Next, we consider the case when $\|u\|_2>D_{\max}=2^N$. Then, we have
\begin{align*}
\sum_{t=1}^T f_t(x_t)-f_t(u)&\leq \sum_{t=1}^T \inner{\nabla f_t(x_t)}{x_t-u}\\
&\leq 2GT\|u\|_2.
\end{align*}
We take $N=\lceil\log T\rceil$, then $T\leq \|u\|_2$. Therefore,
\begin{align*}
\sum_{t=1}^T f_t(x_t)-f_t(u)\leq 2G\|u\|_2^2.
\end{align*}

By combining these two cases above, the desirable regret bound is achieved.
\end{proof}


\iffalse
\begin{proof}
\adnote{$N = \mathcal{O}(\log T)$. This means that $|S| = \mathcal{O}(\log^2 T )$.}
For any $i\in [N]$, there exists a $k_{\star}$ such that 
\begin{equation}
\eta_{k_{\star}}\leq \min\Big\{\frac{1}{512GD_i^2},\sqrt{\frac{\Gamma_i}{D_i^2\sum_{t=1}^T(g_{t}^{i}-m_{t}^{i})^2}}\Big\}\leq 2\eta_{k_{\star}}.
\end{equation}
Moreover, $64\eta_{t}^iD_i z_T\leq 512\eta_k D_i GD_i\leq 1$.
From the proof of Theorem~\ref{thm:thm1}, for base algorithm $\cA_{k_\star}$, we have 
\begin{align}
\regret_{\cA_{k_\star}}(p_\star^i)&=\sum_{t=1}^T\inner{\ell_t}{p_t^{k_\star}-p_\star^i}=\sum_{t=1}^T\inner{g_t}{x_t^{k_\star}-x_i}\\
&\leq \tilde{O}\Big(D_i^2+D_i\sqrt{\Gamma_i\sum_{t=1}^T (g_{t}^i-m_t^i)^2}\Big)-32\eta_{k_\star}\sum_{t=1}^T \inner{x_t}{g_{t}-m_t}^2.\label{eq:basep}
\end{align}

Next, by Theorem~4 and Lemma~1 in \citep{chen2021impossible}, we have $32\eta_k|\ell_t-h_t|\leq 1$ and 
$32\sum_{t=1}^T\sum_{j=1}^N \eta_k(\ell_t^j-h_t^j)^2= 32\eta_k\sum_{t=1}^T\inner{x_t} {\ell_t-h_t}^2$, which can be canceled by the last term in \eqref{eq:basep}. Also, $\sum_{k}\eta_k =O(1/D_1^2)$, $\sum_{k}\eta_k^2 =O(1/D_1^4)$ and $\frac{\sum_{k}\eta_k^2}{\eta_{k_\star}^2} =O(D_i^2 T/D_1^2)$.
Therefore, we have
\begin{equation}
\sum_{t=1}^T\langle \ell_t, p_t-p_*^i\rangle\leq \tilde{O}\Big(D_i^2+D_i\sqrt{\Gamma_i\sum_{t=1}^T (\ell_{t}^i-h_t^i)^2}\Big)
\end{equation}
\end{proof}
\fi

\section{Omitted Details of Section~\ref{sec:fullyfree}}\label{app:sec42}
%In this subsection, we consider the case where both parameters diameter of decision set $D$ and Lipschitz constant $G$ are unknown. We summarize our algorithm in Algorithm~\ref{alg:fullyfree}.



\iffalse
For analysis, we also define a sequence $\{ C_t \}_{t=1}^T$. Let $C_1 = 1$ and if $C_t < \sqrt{\sum_{s=1}^t\frac{\|g_t \|_2}{\max_{k \leq s} \| g_k \|_2 }}$
\begin{align}
    C_{t+1} = 2\sqrt{\sum_{s=1}^t\frac{\|g_t \|_2}{\max_{k \leq s} \| g_k \|_2 }}
\end{align}
\fi

In this section, we also denote $g_t=\nabla f_t(x_t)$ and $m_t=\nabla f_{t-1}(x_{t-1})$. We first note that 
\begin{align*}
\max_{t \leq T}D_t < \sqrt{\sum_{s=1}^t\frac{\|g_s \|_2}{\max\{1,\max_{k \leq s} \| g_k \|_2 \}}} \leq \sqrt{T}~.
\end{align*}
Thus, we need to update $D_t$ at most $O(\log T)$ times. Let $M$ be the number of total updates in $D_t$, where $M = \mathcal{O}(\log T)$. We split the $T$ iterations into $M$ intervals $I_m$ with $m\in[M]$, where the last iteration of $I_m$ (denoted by $t_m$) either equals to $T$ or $D_{t_m+1}\neq D_{t_m}$.

\iffalse
\begin{theorem}
Let both $D$ and $G$ be unknown. Under Assumptions 3, Algorithm~\ref{alg:fullyfree} provides the following regret
\begin{equation*}
\mathbb{E}[\regret_{T}(u)]\leq\tilde{O}\Big((\|u\|_2^2+G_T)(\sqrt{\sigma_{1:T}^2}+\sqrt{\Sigma_{1:T}^2})+G_T^2\|u\|_2^2+\|u\|_2^4+G_T\|u\|_2^3\Big),
\end{equation*}
where $G_T=\max_{t\in [T]} \|\nabla f_t(x_t)\|_2$. 
\end{theorem}
\textcolor{red}{Should we define $G_{\max}=\sup_{T\in\mathbb{N}}G_T$. Then replace $G_T$ in the above theorem by $G_{\max}$. We don't want people to mistake $G_T$ as being $T$ dependent. }
\fi

\subsection{Proof of Theorem~\ref{thm:fullyfree}}
\begin{proof}
We have
\begin{align*}
\regret_T(u)&=\sum_{t=1}^T f_t(x_t) - \sum_{t=1}^T f_t(u) \leq \sum_{t=1}^T \inner{g_t}{x_t - u}\nonumber \\
&= \sum_{m=1}^M \underbrace{\sum_{t \in I_m} \inner{g_t}{x_t - u_t}}_{T_m} + \underbrace{\sum_{t =1}^T \inner{g_t}{u_t - u}}_{T_{\text{extra}}},
\end{align*}
where we define $u_t = \min\{ 1, \frac{D_t}{\| u \|_2} \} u$.

We first consider $T_m$ as
\begin{align*}
    \sum_{t \in I_m} \inner{g_t}{x_t - u_t} = \sum_{t \in I_m} \inner{\tilde{g}_t}{x_t - u_t} + \sum_{t \in I_m} \inner{g_t - \tilde{g}_t}{x_t - u_t}.
\end{align*}

Note that iteration $t$ within interval $I_m$, i.e., $t\in I_m$, the domain has a bounded diameter $D_t$. When $t\in I_m$, we take 
\begin{equation*}
\eta_t = \min\{\frac{1}{64 D_t z_t}, \frac{1}{\sqrt{\sum_{s=t_1}^{t-1}\| \tilde{g}_s - m_s \|_2^2}} \}, 
\end{equation*}
where $t_1$ is first index in $I_m$. Also, we denote the last index in $I_m$ as $t_m$, respectively. From the Line 5 or 8 in \fpfons\, we need to reset $x'_{t_1}=0$ and $A_t$ for all $t\in I_m$ at iteration $t_1$ as follows
\begin{align*}
    A_t = 4z_{t_1}^2 I + \sum_{s=t_1}^{t-1} \eta_s(\nabla_s - m_s)(\nabla_s - m_s)^\top + 4\eta_t z_t^2 I.
\end{align*}
Similar to the proof of Theorem~\ref{thm:thm1}, we have
\begin{align*}
&\sum_{t=t_1}^{t_m} \langle x_t-u_t,\nabla_t\rangle\\
&\leq \sum_{t=t_1}^{t_m} \langle x_{t}-x'_{t+1},\nabla_t-m_t\rangle+D_{\psi_{t_1}}(u,x'_{t_1})+\sum_{t=t_1}^{t_m-1}D_{\psi_{t+1}}(u_t,x'_{t+1})-D_{\psi_t}(u_t,x'_{t+1})\\
&\qquad-\sum_{t=t_1}^{t_m}(D_{\psi_t}(x_t,x'_{t+1})+D_{\psi_t}(x_t,x'_{t})),
\end{align*}
where $\nabla_t=\tilde{g}_t + 32\eta_t\inner{x_t}{\tilde{g}_t-m_t}(\tilde{g}_t-m_t)$.

We first consider the term $\sum_{t=t_1}^{t_m}\langle x_{t}-x'_{t+1},\nabla_t-m_t\rangle$. By Lemma~\ref{lem:term1}, we have $0\leq\langle x_{t}-x'_{t+1},\nabla_t-m_t\rangle \leq 2\|\nabla_t-m_t\|_{A_t^{-1}}^2$. 

Also, we have
\begin{align*}
\|\tilde{g}_t-m_t\|_2&=\|m_t+\frac{B_{t-1}}{B_t}(g_t-m_t)-m_t\|_2\\
&\leq \|g_t-m_t\|_2
\end{align*}

By the definition of $\nabla_t=\tilde{g}_t+32\eta_t\langle x_t,\tilde{g}_t-m_t\rangle(\tilde{g}_t-m_t)$ and $\|\tilde{g}_t-m_t\|_2\leq \|g_t-m_t\|_2$, we have
\begin{align*}
\|\nabla_t-m_t\|_2&=\|\tilde{g}_t-m_t+32\eta_t\langle x_t,\tilde{g}_t-m_t\rangle(\tilde{g}_t-m_t)\|_2\\
&\leq \|\tilde{g}_t-m_t\|_2+32\eta_tD_t\|\tilde{g}_t-m_t\|_2^2\\
&\leq \|\tilde{g}_t-m_t\|_2+32\eta_tD_t z_t\|\tilde{g}_t-m_t\|_2\\
&\leq \frac{3}{2} \|\tilde{g}_t-m_t\|_2\leq \frac{3}{2} \|g_t-m_t\|_2.
\end{align*}
For $t\in I_m$, we also redefine 
\begin{equation*}
\bar{A}_t=4z_{t_1}^2\cdot I+\sum_{s=t_1}^{t} \eta_s (\nabla_s-m_s)(\nabla_s-m_s)^\top.
\end{equation*}
Hence, $A_t\succeq \bar{A}_t$ since $\|\nabla_t-m_t\|_2^2\leq 4\|\tilde{g}_t-m_t\|_2^2\leq 4z_t^2$. Also, for $t\in[t_1+1,t_m]$, we have
\begin{equation*}
(\nabla_t-m_t)(\nabla_t-m_t)^\top=\frac{1}{\eta_t}[\eta_t(\nabla_t-m_t)(\nabla_t-m_t)^\top]=\frac{1}{\eta_t} (\bar{A}_t-\bar{A}_{t-1})
\end{equation*}
Then,
\begin{align*}
\sum_{t=t_1}^{t_m} \|\nabla_t-m_t\|_{A_t^{-1}}^2&=\|\nabla_{t_1}-m_{t_1}\|_{A_{t_1}^{-1}}^2+\sum_{t=t_1+1}^{t_m} \|\nabla_t-m_t\|_{A_t^{-1}}^2\\
&= \frac{1}{4z_{t_1}^2}\|\nabla_{t_1}-m_{t_1}\|_2^2+\sum_{t=t_1+1}^{t_m} \|\nabla_t-m_t\|_{A_t^{-1}}^2\\
&\leq 1+\sum_{t=t_1+1}^{t_m} \|\nabla_t-m_t\|_{A_t^{-1}}^2 \qquad\text{(since $\|\nabla_{t_1}-m_{t_1}\|_2^2\leq 4z_{t_1}^2$)}\\
&\leq \sum_{t=t_1+1}^{t_m} \|\nabla_t-m_t\|_{\bar{A}_t^{-1}}^2+1\\
&\leq \sum_{t=t_1+1}^{t_m} \frac{1}{\eta_t}(\ln|\bar{A}_t|-\ln|\bar{A}_{t-1}|)+1\\
&\leq \frac{1}{\eta_{t_m}}\ln\frac{|\bar{A}_{t_m}|}{|\bar{A}_{t_1}|}+1.
\end{align*}

For $|\bar{A}_{t_m}|$:
\begin{align*}
|\bar{A}_{t_m}|\leq |4z_{t_1}^2 I+\sum_{t=t_1}^{t_m} \eta_t(\nabla_t-m_t)(\nabla_t-m_t)^\top|.
\end{align*}
\begin{align*}
\ln |\bar{A}_{t_m}|&\leq O\Big(r\ln(1+\sum_{t=t_1}^{t_m} \frac{\eta_t}{4z_{t_1}^2}\|\nabla_t-m_t\|_2^2)\Big)\\
&\leq O\Big(r\ln(1+\frac{4z_{t_m}^2}{4z_{t_1}^2}\sum_{t=t_1}^{t_m} \eta_t)\Big)\\
&\leq O\Big(r\ln(1+\frac{\eta_{t_1} z_{t_m}^2}{z_{t_1}^2}T)\Big).
\end{align*}
Therefore, we have
\begin{equation*}
\sum_{t=t_1}^{t_m}\langle x_{t}-x'_{t+1},\nabla_t-m_t\rangle\leq O\Big(\frac{r\ln(T\eta_{t_1} z_{t_m}/z_{t_1})}{\eta_{t_m}}\Big).
\end{equation*}
\textbf{Term} $D_{\psi_{t_1}}(u,x'_{t_1})$:
\begin{equation*}
D_{\psi_{t_1}}(u_t,x'_{t_1})=\frac{1}{2}\|u_t\|^2_{A_{t_1}}\leq O(z_{t_1}^2\|u_t\|_2^2).
\end{equation*}
\textbf{Term} $\sum_{t=t_1}^{t_m-1}D_{\psi_{t+1}}(u_t,x'_{t+1})-D_{\psi_t}(u_t,x'_{t+1})$: 

By the definition of $u_t = \min\{ 1, \frac{D_t}{\| u \|_2}\} u$, we have
\begin{equation*}
\|u_t\|_2=\min\{\|u\|_2,D_t\}.
\end{equation*}
Then, we have
\begin{align*}
&\sum_{t=t_1}^{t_m-1}D_{\psi_{t+1}}(u_t,x'_{t+1})-D_{\psi_t}(u_t,x'_{t+1})\\
&\leq\sum_{t=t_1}^{t_m-1} \frac{\eta_t}{2}\inner{u_t-x'_{t+1}}{\nabla_t-m_t}^2+O(\sum_{t=t_1}^{t_m-1} \eta_t\|u_t-x'_{t+1}\|_2^2(z^2_{t+1}-z_t^2))\\
&\leq \sum_{t=t_1}^{t_m-1} \frac{\eta_t}{2}\inner{u_t-x'_{t+1}}{\nabla_t-m_t}^2+O(\sum_{t=t_1}^{t_m-1} \eta_tD_t^2(z^2_{t+1}-z_t^2)) \qquad \text{(since $\|u_t\|_2\leq \|D_t\|_2$)}\\
&\leq \sum_{t=t_1}^{t_m-1} \frac{\eta_t}{2}\inner{u_t-x'_{t+1}}{\nabla_t-m_t}^2+O(\sum_{t=t_1}^{t_m-1} D_t(z^2_{t+1}-z_t^2)/z_t) \qquad \text{(since $\eta_t\leq\frac{1}{64D_tz_t}$)}\\
&\leq \sum_{t=t_1}^{t_m-1} \frac{\eta_t}{2}\inner{u_t-x'_{t+1}}{\nabla_t-m_t}^2+O(D_{t_m}z^2_{t_m})  \qquad \text{(since $z_t\geq z_1=1$)}\\
&\leq \sum_{t={t_1}}^{t_m-1} 8\eta_t\langle u_t,\tilde{g}_t-m_t\rangle^2+8\eta_t\langle x_{t},\tilde{g}_t-m_t\rangle^2+O\Big(\frac{r\ln(T\eta_{t_1} z_{t_m}/z_{t_1})}{\eta_{t_m}}+D_{t_m}z^2_{t_m}\Big).
\end{align*}

\textbf{Term} $\sum_{t=t_1}^{t_m}(D_{\psi_t}(x_t,x'_{t+1})+D_{\psi_t}(x_t,x'_{t}))$:
\begin{align*}
\sum_{t=t_1}^{t_m}(D_{\psi_t}(x_t,x'_{t+1})+D_{\psi_t}(x_t,x'_{t}))&=\sum_{t=t_1}^{t_m}\frac{1}{2}(\|x_t-x'_{t}\|_{A_t}^2+\|x'_{t+1}-x_t\|_{A_t}^2)\\
&\geq \frac{1}{2}\sum_{t=t_1+1}^{t_m} \|x_t-x'_{t}\|_{A_{t-1}}^2+\frac{1}{2}\sum_{t=t_1+1}^{t_m+1}\|x_{t-1}-x'_t\|_{A_{t-1}}^2\\
&\geq \frac{4z_{t_1}^2}{4}\sum_{t=t_1+1}^{t_m}\|x_t-x_{t-1}\|_2^2=z_1^2\sum_{t=t_1+1}^{t_m}\|x_t-x_{t-1}\|_2^2,
\end{align*}
where the last inequality comes from $A_t\succeq A_{t-1}\succeq 4z_{t_1}^2 I$ when $t\in [t_1+1,t_m]$.

Then, we have
\begin{align*}
&\sum_{t \in I_m} \inner{\tilde{g}_t}{x_t - u_t}\\&\leq O\Big(\frac{r\ln(T\eta_{t_1} z_{t_m}/z_{t_1})}{\eta_{t_m}}+z_{t_1}^2\|u_t\|_2^2+D_t z_{t_m}^2+\sum_{t={t_1}}^{t_m-1} \eta_t\langle u_t,\tilde{g}_t-m_t\rangle^2-z_{t_1}^2 \sum_{t_1+1}^{t_m}\|x_t-x_{t-1}\|_2^2\Big)\\
&\leq \tilde{O}\Big(\sqrt{\sum_{t\in I_m}\| \tilde{g}_t - m_t \|_2^2}+z_{t_1}^2\|u\|_2^2+z_{t_m}^2D_t+\|u\|_2^2\sqrt{\sum_{t\in I_m}\| \tilde{g}_t - m_t \|_2^2}-z_{t_1}^2 \sum_{t_1+1}^{t_m}\|x_t-x_{t-1}\|_2^2\Big).
\end{align*}

Furthermore, by $\|\tilde{g}_t-m_t\|_2\leq \|g_t-m_t\|_2$, we have 
\begin{align*}
    \sum_{t \in I_m} \inner{g_t - \tilde{g}_t}{x_t - u_t} &\leq (D_t + \| u \|_2 ) \sum_{t \in I_m} \| g_t - \tilde{g}_t \| \leq (D_t + \| u \|_2 ) \sum_{t\in I_m} \frac{B_t - B_{t-1}}{B_t} \| g_t - m_t \|_2 \\
    &\leq 2(D_t + \| u \|_2 ) G.
\end{align*}

At the last iteration $T$, we have
\begin{align*}
D_T &< \sqrt{\sum_{t=1}^T\frac{\|g_t \|_2}{\max\{1,\max_{k \leq t} \| g_k \|_2 \}}}\\
&\leq \sqrt{\sum_{t=1}^T \|g_t\|_2}\\
&\leq \sqrt{\sum_{t=1}^T \|g_t-\mathbb{E}[g_t]\|_2+\|\nabla F_t(x_t)\|_2}\\
&=\sqrt{\sum_{t=1}^T \|\nabla f_t(x_t)-\nabla F_t(x_t)\|_2+\|\nabla F_t(x_t)\|_2}.
\end{align*}
Therefore, we have
\begin{align*}
&\sum_{m=1}^M\sum_{t \in I_m} \inner{g_t}{x_t - u_t}\leq \tilde{O}\Big(\|u\|_2^2\sqrt{\sum_{t\in [T]}\| g_t - m_t \|_2^2}+G^2D_T+G^2\|u\|_2^2-z_1^2\sum_{t=2}^T\|x_t-x_{t-1}\|_2^2\Big)\\
&\leq\tilde{O}\Big(\|u\|_2^2\sqrt{\sum_{t\in [T]}\| g_t - m_t \|_2^2}+G^2D_T+G^2\|u\|_2^2-\sum_{t=2}^T\|x_t-x_{t-1}\|_2^2\Big). \qquad \text{(since $z_1=B_0=1$)}
\end{align*}

Now, we bound $T_{\text{extra}}$. We observe that $u_t$ is either $u$ or $\frac{D_t}{\| u \|_2} u$. When $u_t \ne u$,  $\| u \|_2 \geq D_t > \sqrt{\sum_{s=1}^t \frac{\| g_s\|_2}{ \max_{k \leq s} \| g_k\|_2}}$. Once $u_t = u$, it stays there. Let $t^*$ be the last round when $u_t \ne u$.
\begin{align*}
    \sum_{t = 1}^T \inner{g_t}{u_t - u} &= \sum_{t = 1}^{t^*} \inner{g_t}{u_t - u} \leq \sum_{t = 1}^{t^* - 1} \inner{g_t}{u_t - u} + 2 \| u \|_2 G\\
    &\leq 2 \| u \|_2 \sum_{t = 1}^{t^* - 1} \| g_t \|_2 + 2 \| u \|_2 G \\
    &\leq 2 \| u \|_2 G \sum_{t = 1}^{t^* - 1} \frac{\| g_t \|_2}{  \max_{k \leq t^* - 1} \| g_k \|_2 } + 2 \| u \|_2 G \\ 
    &\leq 2 \| u \|_2^3 G + 2 \| u \|_2 G.
\end{align*}

Therefore, we have
\begin{equation*}
\regret_{T}(u)\leq\tilde{O}\Big(\|u\|_2^2\sqrt{\sum_{t\in [T]}\| g_t - m_t \|_2^2}+G^2D_T+G^2\|u\|_2^2+G\| u \|_2^3 -\sum_{t=2}^T\|x_t-x_{t-1}\|_2^2\Big)
\end{equation*}
By Lemma~\ref{lem:ex2sigma}, we have
\begin{align*}
&\|u\|_2^2\sqrt{\sum_{t=1}^{T}\|g_{t}-m_t\|_2^2}-\sum_{t=2}^T\|x_t-x_{t-1}\|_2^2\\
&\leq G^2\|u\|_2^2+2\|u\|_2^2\sqrt{\sum_{t=2}^T\|\nabla F_t(x_{t-1})-\nabla F_{t-1}(x_{t-1})\|_2^2}\\
&\quad+2\sqrt{2}\|u\|_2^2\sqrt{\sum_{t=1}^T\|\nabla f_t(x_t)-\nabla F_t(x_{t})\|_2^2}+2L\|u\|_2^2\sqrt{\sum_{t=2}^T\|x_t-x_{t-1}\|_2^2}-\sum_{t=2}^T\|x_t-x_{t-1}\|_2^2\\
&\leq G^2\|u\|_2^2+L^2\|u\|_2^4+2\|u\|_2^2\sqrt{\sum_{t=2}^T\|\nabla F_t(x_{t-1})-\nabla F_{t-1}(x_{t-1})\|_2^2}\\
&\qquad+2\sqrt{2}\|u\|_2^2\sqrt{\sum_{t=1}^T\|\nabla f_t(x_t)-\nabla F_t(x_{t})\|_2^2}
\end{align*}

Therefore, we can conclude that
\begin{align*}
\mathbb{E}[\regret_{T}(u)]&\leq\tilde{O}\Big(\|u\|_2^2(\sqrt{\sigma_{1:T}^2}+\sqrt{\Sigma_{1:T}^2})+G^2\sqrt{\sigma_{1:T} + \mathfrak{G}_{1:T}}+G^2\|u\|_2^2+\|u\|_2^4+G\|u\|_2^3\Big),
\end{align*}
where $\sigma_{1:T}$ captures the stochastic gradient deviation (without the squared norm) and $\mathfrak{G}_{1:T}$ denotes the sum of maximum expected gradients.
\end{proof}

\iffalse
\section{Fully parameter-free algorithm with experts}

\begin{algorithm}[tb]
\caption{Parameter free algorithms with unknown $D$ and $G$}

    \textbf{Input:} Initial scale $z_1=B_0$ and $\tilde{z}=z_1$.

    \textbf{Initialize:} $D_1=1$ and $x'_1=0$.
    
    \begin{algorithmic}[1]
    \FOR{$t=1,\ldots, T$}
        \STATE Run Algorithm~\ref{alg:metafree} in $D_t$ bounded domain and obtain $x_t$. Play $x_t$ and receive $g_t=\nabla f_t(x_t)$.
        \STATE Construct $\tilde{g}_t=m_t+\frac{B_{t-1}}{B_t}(g_t-m_t)$, where $B_t=\max_{0\leq s\leq t}\|g_s-m_s\|_2$.
        \STATE \bcomment{For each $j\in[N]$, construct $\tilde{g}_t^j=m_t^j+\frac{B_{t-1}^j}{B_t^j}(g_t^j-m_t^j)$, where $B_t^j=\max_{0\leq s\leq t}\|g_s^j-m_s^j\|_2$.}
        \IF{$D_t < \sqrt{ \sum_{s=1}^t \frac{\| g_s - m_s \|_2^2}{ \max\{1, \max_{k \leq s} \| g_k - m_k\|_2^2 \} } }$}
            \STATE \begin{align*}
             D_{t+1} = 2 \sqrt{ \sum_{s=1}^t \frac{\| g_s - m_s \|_2^2}{ \max\{1, \max_{k \leq s} \| g_k - m_k \|_2^2 \} } },
            \end{align*}
        reset $A_{t+1}$ as \eqref{eq:updatedAtfree} and $x'_{t+1}=0$.
        \ENDIF
        \IF{$C_t < \sqrt{\sum_{s=1}^t\frac{\|g_s \|_2}{\max_{k \leq s} \| g_k \|_2 }}$}
        \STATE\begin{align*}
            C_{t+1} = 2\sqrt{\sum_{s=1}^t\frac{\|g_s \|_2}{\max_{k \leq s} \| g_k \|_2 }},
        \end{align*}
        reset $A_{t+1}$ as \eqref{eq:updatedAtfree} and $x'_{t+1}=0$.
        \ENDIF
        \STATE \bcomment{Feed $\tilde{g}_t^j,\forall j\in[N]$ to Algorithm~\ref{alg:metafree} running in the $D_{t+1}$ bounded domain and get $x_{t+1}$, where we take $z_{t+1} = \max\{B_t,\max_{j\in [N]} B_t^j $.\}}
        \STATE\bcomment{\IF {$z_{t+1}/\tilde{z}>T$}
        \STATE $\tilde{z}=z_{t+1}.$
        \ENDIF
        }

        
    \ENDFOR
\end{algorithmic}
\end{algorithm}


Create $N=\log 2T^2$ experts. 
{\color{blue}{ For each $k\in[N]$, we have the following step-size
\begin{equation*}
\eta_k=\frac{1}{128D_t\tilde{z} 2^k}.
\end{equation*}
We define a subset $\cS(t)\subseteq [N]$ as follows
\begin{equation*}
\cS(t)=\{ k\in[N]: \frac{1}{128D_t\tilde{z} 2^k}\leq \frac{1}{256D_tz_{t}}\}
\end{equation*}
}}

We have the following Meta Algorithm and Base Algorithm.
\begin{algorithm}[h]
\caption{Meta Algorithm (MsMwC--Master)}\label{alg:metafree}
\textbf{Input:} a expert set of $\cS$.

\textbf{Initialize:} $p'_1\in\Delta_{\cS}$ such that $p'_{1,k}\propto \eta_k^2$ for each $k\in\cS$.
\begin{algorithmic}[1]
\FOR{$t=1,\dotsc,T$}
\STATE Construct prediction $h_t\in\real^N$ with $h_t^j=\langle \nabla f_{t-1}(x_{t-1}^j), x_t^j \rangle$ and feed it to all MsMwC base algorithms.
\STATE Each  expert $k\in\cS$ runs MsMwC algorithm and plays $w_t^k\in\Delta_N$, where $\eta_{t,j}^k=4\eta_k$ for all $t$ and $j$.
\STATE Receive $w_t^k$ for each $k\in\cS$ and define $H_{t}^k=\inner{w_t^k}{h_t}$.
\STATE Compute $p_t=\argmin_{p\in\Delta_{\cS}}\inner{p}{H_t}+D_{\phi}(p,p'_t)$, where $ \phi(p)=\sum_{k\in\cS}\frac{p_k}{\eta_k}\ln p_k$.
\STATE Play $w_t=\sum_{k\in\cS}p_{t,k}w_t^k\in\Delta_N$.
\STATE Receive $\ell_t\in\real^N$ with $\ell_t^j=\langle\nabla f_{t}(x_{t}^j) , x_t^j \rangle$ and define $L_{t}^k=\inner{w_t^k}{\ell_t}$. 
\STATE Compute 
$p'_{t+1}=\argmin_{p\in\Delta_{\cS}}\inner{p}{L_t+b_t}+D_{\phi}(p,p'_t)$, where $b_t^k=32\eta_k(L_t^k-H_t^k)^2$.
\STATE The final decision is $x_t=\sum_{j=1}^N w_{t,j} x_t^j$.

\ENDFOR
\end{algorithmic}
\end{algorithm}

\begin{algorithm}[tb]
\caption{Base Algorithm $\cA$ (ONS)}\label{alg:basefree}
\begin{algorithmic}[1]
    \STATE \textbf{Input:} learning rate $\eta>0$, $x'_1=0$.
    
    \FOR{$t=1,\ldots, T$}
        \STATE Receive optimistic prediction $m_t$.
        \STATE Update $x_t = \argmin_{x\in\cX}\{\inner{x}{m_t} + D_{\psi_{t}}(x,x_t')\}$  where $\psi_t(x) = \frac{1}{2}\|x\|_{A_t}^2$ and
        \begin{align*}
            A_t &= 4z_1^2 I + \eta\sum_{s=1}^{t-1} (\nabla_s - m_s)(\nabla_s - m_s)^\top + 4z_t^2 I.
        \end{align*}
        
        \STATE Receive $g_t=\nabla f_t(x_t)$.
        \STATE Update $x_{t+1}' = \argmin_{x\in\cX}\{\inner{x}{\nabla_t} + D_{\psi_t}(x, x_t')\}$.
    \ENDFOR
\end{algorithmic}
\end{algorithm}

If the ``doubling'' occurs at the $t$-th iteration, we update $t_a\leftarrow t$ and reset $x'_{t+1}=0$ and the matrix $A_{t+1}$ in Algorithm~\ref{alg:ONS} as follows
\begin{equation}\label{eq:updatedAtfree}
A_{t+1} = 4z_{t_a+1}^2 I \!+\eta \sum_{s=t_a+1}^{t} (\nabla_s - m_s)(\nabla_s - m_s)^\top + 4z_{t+1}^2 I.
\end{equation}

Now, we write the regret bound as
\begin{align*}
\regret_T(u)=\regret_T^{\cA_{k_\star}}(u)+\sum_{t=1}^T\langle \ell_t, w_t-w^{k_\star}\rangle
\end{align*}
\begin{align}
\regret_T^{\cA_{k_\star}}(u)&=\sum_{t=1}^T f_t(x_t^{k_\star}) - \sum_{t=1}^T f_t(u) \leq \sum_{t=1}^T \inner{g_t^{k_\star}}{x_t^i - u}\nonumber \\
&= \sum_{m=1}^M \underbrace{\sum_{t \in I_m} \inner{g_t^{k_\star}}{x_t^{k_\star} - u_t}}_{T_m} + \underbrace{\sum_{t =1}^T \inner{g_t^{k_\star}}{u_t - u}}_{T_{\text{extra}}},
\end{align}
where we define $u_t = \min\{ 1, \frac{D_t}{\| u \|_2},  \frac{C_t}{\| u \|_2} \} u$.

\bcomment{
We first consider the $T_m$-expert part.  For better readability, we eliminate the superscript $k$ throughout the following proof. Now we consider $t\in I_m$ and $\tilde{z}$ is not updated in this interval $I_m$, i.e., $z_{t_m}/\tilde{z}\leq T$. In this case, we have
\begin{equation*}
\frac{1}{\max\{256,T\}D_t z_{t_m}}\geq \frac{1}{128D_t\tilde{z}2^N}.
\end{equation*}
Hence, there exists $k_\star\in \cS(t)$ such that
\begin{equation*}
\eta_{k_\star} \leq \min\Big\{\frac{1}{256D_t z_{t_m}},\sqrt{\frac{r\ln T}{\sum_{t\in \tilde{I}_m}\inner{u}{\tilde{g}_t^{k_\star}-m_t^{k_\star}}^2}}\Big\}\leq 2\eta_{k_\star},
\end{equation*}
where $\tilde{I}_m$ is $I_m$ with last index removed.
For all base ONS algorithms, we take $\eta=4\eta_{k_\star}$.
}


Similar to the proof of Theorem~\ref{thm:thm1}, we have
\begin{align*}
&\sum_{t=t_1}^{t_m} \langle x_t-u_t,\nabla_t\rangle\\
&\leq \sum_{t=t_1}^{t_m} \langle x_{t}-x'_{t+1},\nabla_t-m_t\rangle+D_{\psi_{t_1}}(u,x'_{t_1})+\sum_{t=t_1}^{t_m-1}D_{\psi_{t+1}}(u_t,x'_{t+1})-D_{\psi_t}(u_t,x'_{t+1})\\
&\qquad-\sum_{t=t_1}^{t_m}(D_{\psi_t}(x_t,x'_{t+1})+D_{\psi_t}(x_t,x'_{t})),
\end{align*}
where $\nabla_t=\tilde{g}_t + 32\eta\inner{x_t}{\tilde{g}_t-m_t}(\tilde{g}_t-m_t)$.

We first consider the term $\sum_{t=t_1}^{t_m}\langle x_{t}-x'_{t+1},\nabla_t-m_t\rangle$. By Lemma~\ref{lem:term1}, we have $0\leq\langle x_{t}-x'_{t+1},\nabla_t-m_t\rangle \leq 2\|\nabla_t-m_t\|_{A_t^{-1}}^2$. For $t\in I_m$, we also redefine 
\begin{equation*}
\bar{A}_t=4z_{t_1}^2\cdot I+\eta\sum_{s=t_1}^{t} (\nabla_s-m_s)(\nabla_s-m_s)^\top.
\end{equation*}
By the definition of $\nabla_t=\tilde{g}_t+32\eta\langle x_t,\tilde{g}_t-m_t\rangle(\tilde{g}_t-m_t)$ and $\|\tilde{g}_t-m_t\|_2\leq \|g_t-m_t\|_2$, we have
\begin{align*}
\|\nabla_t-m_t\|_2&=\|\tilde{g}_t-m_t+32\eta\langle x_t,\tilde{g}_t-m_t\rangle(\tilde{g}_t-m_t)\|_2\\
&\leq \|\tilde{g}_t-m_t\|_2+32\eta D_t\|\tilde{g}_t-m_t\|_2^2\\
&\leq \|\tilde{g}_t-m_t\|_2+32\eta D_t z_t\|\tilde{g}_t-m_t\|_2\\
&\bcomment{\leq \|\tilde{g}_t-m_t\|_2+128\eta_{k_\star} D_t z_{t_m}\|\tilde{g}_t-m_t\|_2} \\
&\leq \frac{3}{2} \|\tilde{g}_t-m_t\|_2\leq \frac{3}{2} \|g_t-m_t\|_2.
\end{align*}
Hence, $A_t\succeq \bar{A}_t$ since $\|\nabla_t-m_t\|_2^2\leq 4\|\tilde{g}_t-m_t\|_2^2\leq 4z_t^2$. Also, for $t\in[t_1+1,t_m]$, we have
\begin{equation*}
(\nabla_t-m_t)(\nabla_t-m_t)^\top=\frac{1}{\eta}[\eta(\nabla_t-m_t)(\nabla_t-m_t)^\top]=\frac{1}{\eta} (\bar{A}_t-\bar{A}_{t-1})
\end{equation*}
Then,
\begin{align*}
\sum_{t=t_1}^{t_m} \|\nabla_t-m_t\|_{A_t^{-1}}^2&\leq \sum_{t=t_1}^{t_m} \|\nabla_t-m_t\|_{\bar{A}_t^{-1}}^2\\
&\leq \sum_{t=t_1+1}^{t_m} \|\nabla_t-m_t\|_{\bar{A}_t^{-1}}^2\\
&\leq \sum_{t=t_1+1}^{t_m} \frac{1}{\eta}(\ln|\bar{A}_t|-\ln|\bar{A}_{t-1}|)\\
&\leq \frac{1}{\eta}\ln\frac{|\bar{A}_{t_m}|}{|\bar{A}_{t_1}|}.
\end{align*}

For $|\bar{A}_{t_m}|$:
\begin{align*}
|\bar{A}_{t_m}|\leq |4z_{t_1}^2 I+\sum_{s=t_1}^t \eta(\nabla_t-m_t)(\nabla_t-m_t)^\top|=| I+\sum_{s=t_1}^t \frac{\eta}{4z_{t_1}^2}(\nabla_t-m_t)(\nabla_t-m_t)^\top|.
\end{align*}
\begin{align*}
\ln |\bar{A}_{t_m}|&\leq r\ln(1+\sum_{s=t_1}^{t_m} \frac{\eta}{4z_{t_1}^2}\|\nabla_t-m_t\|_2^2),\\
&\leq r\ln(1+\frac{4z_{t_m}^2}{4z_{t_1}^2}\sum_{s=t_1}^{t_m} \eta),\\
&\leq r\ln(1+\frac{\eta z_{t_m}^2}{z_{t_1}^2}T).
\end{align*}
Therefore, we have
\begin{equation}
\sum_{t=t_1}^{t_m}\langle x_{t}-x'_{t+1},\nabla_t-m_t\rangle\leq O\Big(\frac{r\ln(T\eta z_{t_m}/z_{t_1})}{\eta}\Big).
\end{equation}
Term $D_{\psi_{t_1}}(u,x'_{t_1})$:
\begin{equation*}
D_{\psi_{t_1}}(u_t,x'_{t_1})=\frac{1}{2}\|u_t\|^2_{A_{t_1}}\leq O(z_{t_1}^2\|u_t\|_2^2).
\end{equation*}
Term $\sum_{t=t_1}^{t_m-1}D_{\psi_{t+1}}(u_t,x'_{t+1})-D_{\psi_t}(u_t,x'_{t+1})$:
\begin{align*}
&\sum_{t=t_1}^{t_m-1}D_{\psi_{t+1}}(u_t,x'_{t+1})-D_{\psi_t}(u_t,x'_{t+1})\\
&\leq \sum_{t={t_1}}^{t_m-1} 8\eta\langle u_t,\tilde{g}_t-m_t\rangle^2+8\eta\langle x_{t},\tilde{g}_t-m_t\rangle^2+O\Big(\frac{r\ln(T\eta z_{t_m}/z_{t_1})}{\eta}+D_t z_{t_m}\Big).
\end{align*}
Term $\sum_{t=t_1}^{t_m}(D_{\psi_t}(x_t,x'_{t+1})+D_{\psi_t}(x_t,x'_{t}))$:
\begin{align*}
\sum_{t=t_1}^{t_m}(D_{\psi_t}(x_t,x'_{t+1})+D_{\psi_t}(x_t,x'_{t}))\geq z_{t_1}^2 \sum_{t_1+1}^{t_m}\|x_t-x_{t-1}\|_2^2.
\end{align*}
Then, we have
\begin{align*}
&\sum_{t \in I_m} \inner{\tilde{g}_t}{x_t - u_t}\\&\leq O\Big(\frac{r\ln(T\eta z_{t_m}/z_{t_1})}{\eta}+z_{t_1}^2\|u_t\|_2^2+D_t z_{t_m}+\sum_{t={t_1}}^{t_m} \eta\langle u_t,\tilde{g}_t-m_t\rangle^2-z_{t_1}^2 \sum_{t_1+1}^{t_m}\|x_t-x_{t-1}\|_2^2\Big)-8\eta\sum_{t=t_1}^{t_m} \inner{x_t}{\tilde{g}_t-m_t}^2\\
&\leq O\Big(\frac{r\ln(T\eta z_{t_m}/z_{t_1})}{\eta_{k_\star}}+z_{t_1}^2\|u_t\|_2^2+D_t z_{t_m}+\sum_{t={t_1}}^{t_m} \eta_{k_\star}\langle u_t,\tilde{g}_t-m_t\rangle^2-z_{t_1}^2 \sum_{t_1+1}^{t_m}\|x_t-x_{t-1}\|_2^2\Big)-32\eta_{k_\star}\sum_{t=t_1}^{t_m} \inner{x_t}{\tilde{g}_t-m_t}^2
\end{align*}

For the $T_m$-Meta term: Satisfy the condition $32\eta_k|\inner{x_t^k}{\tilde{g}_t-m_t}|\leq 32\eta_kD_tz_t\leq 1$ and
\begin{align*}
\sum_{t \in I_m} \inner{\tilde{g}_t}{x_t - x_t^{k_\star}}&=\sum_{t \in I_m}\inner{\tilde{\ell}_t}{w_t - w_{k_\star}}\\
&=\frac{1}{\eta_{k_\star}}\ln (\frac{\sum_{k}\eta_k^2}{\eta^2_{k_\star}})+\frac{\sum_k \eta_k}{\sum_k \eta_k^2}+32\eta_{k_\star}\sum_{t=t_1}^{t_m}\inner{x_t^{k_\star}}{\tilde{g}_t-m_t}^2.
\end{align*}
Here, $\sum_k \eta_k=\Theta(\frac{1}{\tilde{z}})$, $\sum_k \eta_k^2=\Theta(\frac{1}{\tilde{z}^2})$ and $\frac{\sum_{k}\eta_k^2}{\eta^2_{k_\star}}=O(\ln T)$.

\bcomment{Shuche: here is not correct from the first equation to the second. The term is $\sum_{t={t_1}}^{t_m} \eta_{k_\star}\langle u_t,\tilde{g}_t-m_t\rangle^2$ not $\sum_{t={t_1}}^{t_m} \eta_{k_\star}\langle u_t,\tilde{g}_t^{k_\star}-m_t^{k_\star}\rangle^2$.}

Therefore, we have
\begin{align*}
\sum_{t \in I_m} \inner{\tilde{g}_t}{x_t^{k_\star} - u_t}&\leq O\Big(\frac{r\ln(T\eta z_{t_m}/z_{t_1})}{\eta_{k_\star}}+z_{t_1}^2\|u_t\|_2^2+D_t z_{t_m}+\sum_{t={t_1}}^{t_m} \eta_{k_\star}\langle u_t,\tilde{g}_t-m_t\rangle^2-z_{t_1}^2 \sum_{t_1+1}^{t_m}\|x_t-x_{t-1}\|_2^2\Big)\\
&\leq \tilde{O}\Big( \|u\|_2\sqrt{\sum_{t\in I_m}\| \tilde{g}_t^{k_\star} - m_t^{k_\star} \|_2^2}+ z_{t_1}^2\|u_t\|_2^2+D_t z_{t_m}-z_{t_1}^2 \sum_{t_1+1}^{t_m}\|x_t^{k_\star}-x_{t-1}^{k_\star}\|_2^2\Big)
\end{align*}

{\color{red}
\begin{align}
\|u\|_2\sqrt{\sum_{t\in I_m}\| \tilde{g}_t^{k_\star} - m_t^{k_\star} \|_2^2} \leq \|u\|_2 \sqrt{\left( \sum_{t \in I_m} c_1 \| g_t^{k^*} - \nabla F_t() \| + \ldots + \sum_{t \in I_m}\| x_t^{k^*} - x_{t-1}^{k*}\|^2 \right)}
\end{align}
Eventually,
\begin{align}
\|u\|_2 \sum_{m =1}^M \sqrt{\left( \sum_{t \in I_m} c_1 \| g_t^{k^*} - \nabla F_t() \| \right)}
\end{align}
}

\bcomment{How to deal with $-z_{t_1}^2 \sum_{t_1+1}^{t_m}\|x_t^{k_\star}-x_{t-1}^{k_\star}\|_2^2$, in each interval we have different expert set $\cS(t)$ and correspond to the different $k_\star$.
]\\
}

Next, we show that this desired bound also holds when there are several restarts in the interval $I_m$.



\section{Strongly convex case}
\begin{algorithm}[h]
\caption{Parameter free algorithms with unknown $D$ and $G$}

    \textbf{Input:} Initial scale $z_1=B_0$ and $\tilde{z}=z_1$.

    \textbf{Initialize:} $D_1=1$ and $x'_1=0$.
    
    \begin{algorithmic}[1]
    \FOR{$t=1,\ldots, T$}
        \STATE Run Algorithm~\ref{alg:metafree} in $D_t$ bounded domain and obtain $x_t$. Play $x_t$ and receive $g_t=\nabla f_t(x_t)$.
        \STATE Construct $\tilde{g}_t=m_t+\frac{B_{t-1}}{B_t}(g_t-m_t)$, where $B_t=\max_{0\leq s\leq t}\|g_s-m_s\|_2$.
        \STATE \bcomment{For each $j\in[N]$, construct $\tilde{g}_t^j=m_t^j+\frac{B_{t-1}^j}{B_t^j}(g_t^j-m_t^j)$, where $B_t^j=\max_{0\leq s\leq t}\|g_s^j-m_s^j\|_2$.}
        \IF{$D_t < \sqrt{ \sum_{s=1}^t \frac{\| g_s - m_s \|_2^2}{ \max\{1, \max_{k \leq s} \| g_k - m_k\|_2^2 \} } }$}
            \STATE \begin{align*}
             D_{t+1} = 2 \sqrt{ \sum_{s=1}^t \frac{\| g_s - m_s \|_2^2}{ \max\{1, \max_{k \leq s} \| g_k - m_k \|_2^2 \} } },
            \end{align*}
        reset $A_{t+1}$ as \eqref{eq:updatedAtfree} and $x'_{t+1}=0$.
        \ENDIF
        \IF{$C_t < \sqrt{\sum_{s=1}^t\frac{\|g_s \|_2}{\max_{k \leq s} \| g_k \|_2 }}$}
        \STATE\begin{align*}
            C_{t+1} = 2\sqrt{\sum_{s=1}^t\frac{\|g_s \|_2}{\max_{k \leq s} \| g_k \|_2 }},
        \end{align*}
        reset $A_{t+1}$ as \eqref{eq:updatedAtfree} and $x'_{t+1}=0$.
        \ENDIF
        \STATE \bcomment{Feed $\tilde{g}_t^j,\forall j\in[N]$ to Algorithm~\ref{alg:metafree} running in the $D_{t+1}$ bounded domain and get $x_{t+1}$, where we take $z_{t+1} = \max\{B_t,\max_{j\in [N]} B_t^j $.\}}
        \STATE\bcomment{\IF {$z_{t+1}/\tilde{z}>T$}
        \STATE $\tilde{z}=z_{t+1}.$
        \ENDIF
        }

        
    \ENDFOR
\end{algorithmic}
\end{algorithm}

We have the following Meta Algorithm and Base Algorithm.
\begin{algorithm}[h]
\caption{Meta Algorithm (MsMwC)}\label{alg:metafree}
\textbf{Input:} $w'_1\in\Delta_N$.
\begin{algorithmic}[1]
\FOR{$t=1,\dotsc,T$}
\STATE Receive $x_t^1,\dotsc,x_t^N$ from expert algorithms and the prediction $h_t\in\real^N$.
\STATE Compute $w_t=\argmin_{w\in\Delta_N}\inner{w}{M_t}+D_{\phi}(w,w'_t)$, where $\phi(w)=\sum_{j=1}^N \frac{w_j}{\eta_{j}}\ln w_j$.
\STATE Play $x_t=\sum_{j=1}^N w_t x_t^j$, receive $\ell_t$ and construct correction term $a_t\in\real^N$.
\STATE Compute $w'_{t+1}=\argmin_{w\in\Delta_N}\inner{w}{\ell_t+a_t}+D_{\phi}(w,w'_t)$.
\ENDFOR
\end{algorithmic}
\end{algorithm}

\begin{algorithm}[h]
\caption{Base Algorithm $\cA$ (ONS)}\label{alg:basefree}
\begin{algorithmic}[1]
    \STATE \textbf{Input:} learning rate $\eta>0$, $x'_1=0$.
    
    \FOR{$t=1,\ldots, T$}
        \STATE Receive optimistic prediction $m_t$.
        \STATE Update $x_t = \argmin_{x\in\cX}\{\inner{x}{m_t} + D_{\psi_{t}}(x,x_t')\}$  where $\psi_t(x) = \frac{A_t}{2}\|x\|^2$ and
        \begin{align*}
            A_t =\eta(8z_1^2+\sum_{s=1}^{t-1}\|\nabla_s-m_s\|_2^2).
        \end{align*}
        
        \STATE Receive $g_t=\nabla f_t(x_t)$.
        \STATE Update $x_{t+1}' = \argmin_{x\in\cX}\{\inner{x}{\nabla_t} + D_{\psi_t}(x, x_t')\}$.
    \ENDFOR
\end{algorithmic}
\end{algorithm}

We denote $g_t=\nabla f_t(x_t)$, $m_t=\nabla f_{t-1}(x_{t-1})$, $M_t$ such that $M_t^j=\inner{m_t}{x_t^j}$, $L_t$ such that $L_t^j=\inner{\tilde{g}_t}{x_t^j}$, $a_t$ such that $a_t^j=32\eta_j(L_t^j-(M_t^j+\inner{w_t}{L_t-M_t}))^2$. Also, $c_t$ is defined as
\begin{equation*}
c_t(x)=\inner{\tilde{g}_t}{x}+64\eta\|x-x_t\|^2\|\tilde{g}_t-m_t\|^2.
\end{equation*}
Now, the regret can be decomposed as
\begin{align*}
\regret_T(u)&=\sum_{t=1}^T f_t(x_t) - \sum_{t=1}^T f_t(u) \leq \sum_{t=1}^T \inner{g_t}{x_t - u}-\sum_{t=1}^T \frac{\lambda}{2}\|x_t-u\|_2^2\leq \sum_{t=1}^T \inner{g_t}{x_t - u}-\sum_{t=1}^T \lambda\|x_t-u_t\|_2^2 \\
&= \sum_{m=1}^M \sum_{t \in I_m} \inner{g_t}{x_t - u_t} + \sum_{t =1}^T \inner{g_t}{u_t - u}-\sum_{t=1}^T \lambda\|x_t-u_t\|_2^2\\
& = \sum_{m=1}^M\sum_{t \in I_m} \inner{\tilde{g}_t}{x_t - u_t} + \sum_{m=1}^M\sum_{t \in I_m} \inner{g_t - \tilde{g}_t}{x_t - u_t}+ \sum_{t =1}^T \inner{g_t}{u_t - u}-\sum_{t=1}^T \lambda\|x_t-u_t\|_2^2.\\
&= \sum_{m=1}^M \underbrace{\sum_{t \in I_m} \inner{\tilde{g}_t}{x_t - x_t^{k_\star}}}_{T_m-\text{Meta}} + \sum_{m=1}^M \underbrace{\sum_{t \in I_m} \inner{\tilde{g}_t}{x_t^{k_\star}-u_t}}_{T_m-\text{Expert}}+ \underbrace{\sum_{t =1}^T \inner{g_t}{u_t - u}}_{T_{\text{extra}}}+ \sum_{m=1}^M\sum_{t \in I_m} \inner{g_t - \tilde{g}_t}{x_t - u_t}-\sum_{t=1}^T \lambda\|x_t-u_t\|_2^2.
\end{align*}

We first consider $T_m$-Expert part:
\begin{align*}
&\sum_{t=t_1}^{t_m} \langle x_t-u_t,\nabla_t\rangle\\
&\leq \sum_{t=t_1}^{t_m} \langle x_{t}-x'_{t+1},\nabla_t-m_t\rangle+D_{\psi_{t_1}}(u,x'_{t_1})+\sum_{t=t_1}^{t_m-1}D_{\psi_{t+1}}(u_t,x'_{t+1})-D_{\psi_t}(u_t,x'_{t+1})
\end{align*}
We also have 
\begin{equation*}
0\leq\langle x_{t}-x'_{t+1},\nabla_t-m_t\rangle \leq \frac{2}{A_t}\|\nabla_t-m_t\|^2
\end{equation*}
Since $A_t =4z_1^2+\eta\sum_{s=1}^{t-1}\|\nabla_s-m_s\|_2^2+4z_t^2$, for $t\in I_m$, we also redefine 
\begin{equation*}
\bar{A}_t=4z_{t_1}^2+\eta\sum_{s=t_1}^{t} \|\nabla_s-m_s\|^2.
\end{equation*}
Then, by $1-x\leq \ln(1/x)$,
\begin{equation*}
\frac{1}{\eta}\sum_{t=t_1}^{t_m}\frac{2}{A_t}\eta\|\nabla_t-m_t\|^2\leq \frac{2}{\eta}\sum_{t=t_1}^{t_m} (1-\frac{A_{t-1}}{A_t})\leq \frac{2}{\eta}\ln\frac{\bar{A}_{t_m}}{\bar{A}_{t_1}}\leq O(\frac{\ln T}{\eta})
\end{equation*}

Term $\sum_{t=t_1}^{t_m-1}D_{\psi_{t+1}}(u_t,x'_{t+1})-D_{\psi_t}(u_t,x'_{t+1})$:
\begin{align*}
&\sum_{t=t_1}^{t_m-1}D_{\psi_{t+1}}(u_t,x'_{t+1})-D_{\psi_t}(u_t,x'_{t+1})\\
& \leq \sum_{t=t_1}^{t_m-1} \eta\langle u_t-x_{t},\nabla_t-m_t\rangle^2+\sum_{t=t_1}^{t_m-1}\eta\langle x_{t}-x'_{t+1},\nabla_t-m_t\rangle^2+O() \\
&\leq \sum_{t={t_1}}^{t_m-1} 8\eta\langle u_t,\tilde{g}_t-m_t\rangle^2+8\eta\langle x_{t},\tilde{g}_t-m_t\rangle^2+O\Big(\frac{r\ln(T\eta z_{t_m}/z_{t_1})}{\eta}+D_t z_{t_m}\Big).
\end{align*}
\fi

\iffalse
\newpage
\section*{NeurIPS Paper Checklist}

\iffalse
%%% BEGIN INSTRUCTIONS %%%
The checklist is designed to encourage best practices for responsible machine learning research, addressing issues of reproducibility, transparency, research ethics, and societal impact. Do not remove the checklist: {\bf The papers not including the checklist will be desk rejected.} The checklist should follow the references and follow the (optional) supplemental material.  The checklist does NOT count towards the page
limit. 

Please read the checklist guidelines carefully for information on how to answer these questions. For each question in the checklist:
\begin{itemize}
    \item You should answer \answerYes{}, \answerNo{}, or \answerNA{}.
    \item \answerNA{} means either that the question is Not Applicable for that particular paper or the relevant information is Not Available.
    \item Please provide a short (1–2 sentence) justification right after your answer (even for NA). 
   % \item {\bf The papers not including the checklist will be desk rejected.}
\end{itemize}

{\bf The checklist answers are an integral part of your paper submission.} They are visible to the reviewers, area chairs, senior area chairs, and ethics reviewers. You will be asked to also include it (after eventual revisions) with the final version of your paper, and its final version will be published with the paper.

The reviewers of your paper will be asked to use the checklist as one of the factors in their evaluation. While "\answerYes{}" is generally preferable to "\answerNo{}", it is perfectly acceptable to answer "\answerNo{}" provided a proper justification is given (e.g., "error bars are not reported because it would be too computationally expensive" or "we were unable to find the license for the dataset we used"). In general, answering "\answerNo{}" or "\answerNA{}" is not grounds for rejection. While the questions are phrased in a binary way, we acknowledge that the true answer is often more nuanced, so please just use your best judgment and write a justification to elaborate. All supporting evidence can appear either in the main paper or the supplemental material, provided in appendix. If you answer \answerYes{} to a question, in the justification please point to the section(s) where related material for the question can be found.

IMPORTANT, please:
\begin{itemize}
    \item {\bf Delete this instruction block, but keep the section heading ``NeurIPS Paper Checklist"},
    \item  {\bf Keep the checklist subsection headings, questions/answers and guidelines below.}
    \item {\bf Do not modify the questions and only use the provided macros for your answers}.
\end{itemize} 
 

%%% END INSTRUCTIONS %%%
\fi

\begin{enumerate}

\item {\bf Claims}
    \item[] Question: Do the main claims made in the abstract and introduction accurately reflect the paper's contributions and scope?
    \item[] Answer: \answerYes{} % Replace by \answerYes{}, \answerNo{}, or \answerNA{}.
    \item[] Justification: The abstract and introduction accurately reflect the paper’s contributions and scope.
    \item[] Guidelines:
    \begin{itemize}
        \item The answer NA means that the abstract and introduction do not include the claims made in the paper.
        \item The abstract and/or introduction should clearly state the claims made, including the contributions made in the paper and important assumptions and limitations. A No or NA answer to this question will not be perceived well by the reviewers. 
        \item The claims made should match theoretical and experimental results, and reflect how much the results can be expected to generalize to other settings. 
        \item It is fine to include aspirational goals as motivation as long as it is clear that these goals are not attained by the paper. 
    \end{itemize}

\item {\bf Limitations}
    \item[] Question: Does the paper discuss the limitations of the work performed by the authors?
    \item[] Answer: \answerYes{} % Replace by \answerYes{}, \answerNo{}, or \answerNA{}.
    \item[] Justification: The limitations of our regret bound are discussed in Remarks 4.2, 4.3, and 4.5.
    \item[] Guidelines:
    \begin{itemize}
        \item The answer NA means that the paper has no limitation while the answer No means that the paper has limitations, but those are not discussed in the paper. 
        \item The authors are encouraged to create a separate "Limitations" section in their paper.
        \item The paper should point out any strong assumptions and how robust the results are to violations of these assumptions (e.g., independence assumptions, noiseless settings, model well-specification, asymptotic approximations only holding locally). The authors should reflect on how these assumptions might be violated in practice and what the implications would be.
        \item The authors should reflect on the scope of the claims made, e.g., if the approach was only tested on a few datasets or with a few runs. In general, empirical results often depend on implicit assumptions, which should be articulated.
        \item The authors should reflect on the factors that influence the performance of the approach. For example, a facial recognition algorithm may perform poorly when image resolution is low or images are taken in low lighting. Or a speech-to-text system might not be used reliably to provide closed captions for online lectures because it fails to handle technical jargon.
        \item The authors should discuss the computational efficiency of the proposed algorithms and how they scale with dataset size.
        \item If applicable, the authors should discuss possible limitations of their approach to address problems of privacy and fairness.
        \item While the authors might fear that complete honesty about limitations might be used by reviewers as grounds for rejection, a worse outcome might be that reviewers discover limitations that aren't acknowledged in the paper. The authors should use their best judgment and recognize that individual actions in favor of transparency play an important role in developing norms that preserve the integrity of the community. Reviewers will be specifically instructed to not penalize honesty concerning limitations.
    \end{itemize}

\item {\bf Theory assumptions and proofs}
    \item[] Question: For each theoretical result, does the paper provide the full set of assumptions and a complete (and correct) proof?
    \item[] Answer: \answerYes{} % Replace by \answerYes{}, \answerNo{}, or \answerNA{}.
    \item[] Justification: The problem setting is clearly introduced in Section 2 and all theoretical results are accompanied by proofs.
    \item[] Guidelines:
    \begin{itemize}
        \item The answer NA means that the paper does not include theoretical results. 
        \item All the theorems, formulas, and proofs in the paper should be numbered and cross-referenced.
        \item All assumptions should be clearly stated or referenced in the statement of any theorems.
        \item The proofs can either appear in the main paper or the supplemental material, but if they appear in the supplemental material, the authors are encouraged to provide a short proof sketch to provide intuition. 
        \item Inversely, any informal proof provided in the core of the paper should be complemented by formal proofs provided in appendix or supplemental material.
        \item Theorems and Lemmas that the proof relies upon should be properly referenced. 
    \end{itemize}

    \item {\bf Experimental result reproducibility}
    \item[] Question: Does the paper fully disclose all the information needed to reproduce the main experimental results of the paper to the extent that it affects the main claims and/or conclusions of the paper (regardless of whether the code and data are provided or not)?
    \item[] Answer: \answerNA{} % Replace by \answerYes{}, \answerNo{}, or \answerNA{}.
    \item[] Justification: This paper is a fully theoretical paper without an experimental section.
    \item[] Guidelines:
    \begin{itemize}
        \item The answer NA means that the paper does not include experiments.
        \item If the paper includes experiments, a No answer to this question will not be perceived well by the reviewers: Making the paper reproducible is important, regardless of whether the code and data are provided or not.
        \item If the contribution is a dataset and/or model, the authors should describe the steps taken to make their results reproducible or verifiable. 
        \item Depending on the contribution, reproducibility can be accomplished in various ways. For example, if the contribution is a novel architecture, describing the architecture fully might suffice, or if the contribution is a specific model and empirical evaluation, it may be necessary to either make it possible for others to replicate the model with the same dataset, or provide access to the model. In general. releasing code and data is often one good way to accomplish this, but reproducibility can also be provided via detailed instructions for how to replicate the results, access to a hosted model (e.g., in the case of a large language model), releasing of a model checkpoint, or other means that are appropriate to the research performed.
        \item While NeurIPS does not require releasing code, the conference does require all submissions to provide some reasonable avenue for reproducibility, which may depend on the nature of the contribution. For example
        \begin{enumerate}
            \item If the contribution is primarily a new algorithm, the paper should make it clear how to reproduce that algorithm.
            \item If the contribution is primarily a new model architecture, the paper should describe the architecture clearly and fully.
            \item If the contribution is a new model (e.g., a large language model), then there should either be a way to access this model for reproducing the results or a way to reproduce the model (e.g., with an open-source dataset or instructions for how to construct the dataset).
            \item We recognize that reproducibility may be tricky in some cases, in which case authors are welcome to describe the particular way they provide for reproducibility. In the case of closed-source models, it may be that access to the model is limited in some way (e.g., to registered users), but it should be possible for other researchers to have some path to reproducing or verifying the results.
        \end{enumerate}
    \end{itemize}


\item {\bf Open access to data and code}
    \item[] Question: Does the paper provide open access to the data and code, with sufficient instructions to faithfully reproduce the main experimental results, as described in supplemental material?
    \item[] Answer: \answerNA{} % Replace by \answerYes{}, \answerNo{}, or \answerNA{}.
    \item[] Justification: This paper is a fully theoretical paper without an experimental section.
    \item[] Guidelines:
    \begin{itemize}
        \item The answer NA means that paper does not include experiments requiring code.
        \item Please see the NeurIPS code and data submission guidelines (\url{https://nips.cc/public/guides/CodeSubmissionPolicy}) for more details.
        \item While we encourage the release of code and data, we understand that this might not be possible, so “No” is an acceptable answer. Papers cannot be rejected simply for not including code, unless this is central to the contribution (e.g., for a new open-source benchmark).
        \item The instructions should contain the exact command and environment needed to run to reproduce the results. See the NeurIPS code and data submission guidelines (\url{https://nips.cc/public/guides/CodeSubmissionPolicy}) for more details.
        \item The authors should provide instructions on data access and preparation, including how to access the raw data, preprocessed data, intermediate data, and generated data, etc.
        \item The authors should provide scripts to reproduce all experimental results for the new proposed method and baselines. If only a subset of experiments are reproducible, they should state which ones are omitted from the script and why.
        \item At submission time, to preserve anonymity, the authors should release anonymized versions (if applicable).
        \item Providing as much information as possible in supplemental material (appended to the paper) is recommended, but including URLs to data and code is permitted.
    \end{itemize}


\item {\bf Experimental setting/details}
    \item[] Question: Does the paper specify all the training and test details (e.g., data splits, hyperparameters, how they were chosen, type of optimizer, etc.) necessary to understand the results?
    \item[] Answer: \answerNA{} % Replace by \answerYes{}, \answerNo{}, or \answerNA{}.
    \item[] Justification: This paper is a fully theoretical paper without an experimental section.
    \item[] Guidelines:
    \begin{itemize}
        \item The answer NA means that the paper does not include experiments.
        \item The experimental setting should be presented in the core of the paper to a level of detail that is necessary to appreciate the results and make sense of them.
        \item The full details can be provided either with the code, in appendix, or as supplemental material.
    \end{itemize}

\item {\bf Experiment statistical significance}
    \item[] Question: Does the paper report error bars suitably and correctly defined or other appropriate information about the statistical significance of the experiments?
    \item[] Answer: \answerNA{} % Replace by \answerYes{}, \answerNo{}, or \answerNA{}.
    \item[] Justification: This paper is a fully theoretical paper without an experimental section.
    \item[] Guidelines:
    \begin{itemize}
        \item The answer NA means that the paper does not include experiments.
        \item The authors should answer "Yes" if the results are accompanied by error bars, confidence intervals, or statistical significance tests, at least for the experiments that support the main claims of the paper.
        \item The factors of variability that the error bars are capturing should be clearly stated (for example, train/test split, initialization, random drawing of some parameter, or overall run with given experimental conditions).
        \item The method for calculating the error bars should be explained (closed form formula, call to a library function, bootstrap, etc.)
        \item The assumptions made should be given (e.g., Normally distributed errors).
        \item It should be clear whether the error bar is the standard deviation or the standard error of the mean.
        \item It is OK to report 1-sigma error bars, but one should state it. The authors should preferably report a 2-sigma error bar than state that they have a 96\% CI, if the hypothesis of Normality of errors is not verified.
        \item For asymmetric distributions, the authors should be careful not to show in tables or figures symmetric error bars that would yield results that are out of range (e.g. negative error rates).
        \item If error bars are reported in tables or plots, The authors should explain in the text how they were calculated and reference the corresponding figures or tables in the text.
    \end{itemize}

\item {\bf Experiments compute resources}
    \item[] Question: For each experiment, does the paper provide sufficient information on the computer resources (type of compute workers, memory, time of execution) needed to reproduce the experiments?
    \item[] Answer: \answerNA{} % Replace by \answerYes{}, \answerNo{}, or \answerNA{}.
    \item[] Justification: This paper is a fully theoretical paper without an experimental section.
    \item[] Guidelines:
    \begin{itemize}
        \item The answer NA means that the paper does not include experiments.
        \item The paper should indicate the type of compute workers CPU or GPU, internal cluster, or cloud provider, including relevant memory and storage.
        \item The paper should provide the amount of compute required for each of the individual experimental runs as well as estimate the total compute. 
        \item The paper should disclose whether the full research project required more compute than the experiments reported in the paper (e.g., preliminary or failed experiments that didn't make it into the paper). 
    \end{itemize}
    
\item {\bf Code of ethics}
    \item[] Question: Does the research conducted in the paper conform, in every respect, with the NeurIPS Code of Ethics \url{https://neurips.cc/public/EthicsGuidelines}?
    \item[] Answer: \answerYes{} % Replace by \answerYes{}, \answerNo{}, or \answerNA{}.
    \item[] Justification: The authors confirm with the NeurIPS Code of Ethics.
    \item[] Guidelines:
    \begin{itemize}
        \item The answer NA means that the authors have not reviewed the NeurIPS Code of Ethics.
        \item If the authors answer No, they should explain the special circumstances that require a deviation from the Code of Ethics.
        \item The authors should make sure to preserve anonymity (e.g., if there is a special consideration due to laws or regulations in their jurisdiction).
    \end{itemize}


\item {\bf Broader impacts}
    \item[] Question: Does the paper discuss both potential positive societal impacts and negative societal impacts of the work performed?
    \item[] Answer: \answerNA{} % Replace by \answerYes{}, \answerNo{}, or \answerNA{}.
    \item[] Justification: The work is mainly theoretical, thus there is no identifiable societal impact.
    \item[] Guidelines:
    \begin{itemize}
        \item The answer NA means that there is no societal impact of the work performed.
        \item If the authors answer NA or No, they should explain why their work has no societal impact or why the paper does not address societal impact.
        \item Examples of negative societal impacts include potential malicious or unintended uses (e.g., disinformation, generating fake profiles, surveillance), fairness considerations (e.g., deployment of technologies that could make decisions that unfairly impact specific groups), privacy considerations, and security considerations.
        \item The conference expects that many papers will be foundational research and not tied to particular applications, let alone deployments. However, if there is a direct path to any negative applications, the authors should point it out. For example, it is legitimate to point out that an improvement in the quality of generative models could be used to generate deepfakes for disinformation. On the other hand, it is not needed to point out that a generic algorithm for optimizing neural networks could enable people to train models that generate Deepfakes faster.
        \item The authors should consider possible harms that could arise when the technology is being used as intended and functioning correctly, harms that could arise when the technology is being used as intended but gives incorrect results, and harms following from (intentional or unintentional) misuse of the technology.
        \item If there are negative societal impacts, the authors could also discuss possible mitigation strategies (e.g., gated release of models, providing defenses in addition to attacks, mechanisms for monitoring misuse, mechanisms to monitor how a system learns from feedback over time, improving the efficiency and accessibility of ML).
    \end{itemize}
    
\item {\bf Safeguards}
    \item[] Question: Does the paper describe safeguards that have been put in place for responsible release of data or models that have a high risk for misuse (e.g., pretrained language models, image generators, or scraped datasets)?
    \item[] Answer: \answerNA{} % Replace by \answerYes{}, \answerNo{}, or \answerNA{}.
    \item[] Justification: The work is mainly theoretical.
    \item[] Guidelines:
    \begin{itemize}
        \item The answer NA means that the paper poses no such risks.
        \item Released models that have a high risk for misuse or dual-use should be released with necessary safeguards to allow for controlled use of the model, for example by requiring that users adhere to usage guidelines or restrictions to access the model or implementing safety filters. 
        \item Datasets that have been scraped from the Internet could pose safety risks. The authors should describe how they avoided releasing unsafe images.
        \item We recognize that providing effective safeguards is challenging, and many papers do not require this, but we encourage authors to take this into account and make a best faith effort.
    \end{itemize}

\item {\bf Licenses for existing assets}
    \item[] Question: Are the creators or original owners of assets (e.g., code, data, models), used in the paper, properly credited and are the license and terms of use explicitly mentioned and properly respected?
    \item[] Answer: \answerNA{} % Replace by \answerYes{}, \answerNo{}, or \answerNA{}.
    \item[] Justification: The paper does not use existing assets.
    \item[] Guidelines:
    \begin{itemize}
        \item The answer NA means that the paper does not use existing assets.
        \item The authors should cite the original paper that produced the code package or dataset.
        \item The authors should state which version of the asset is used and, if possible, include a URL.
        \item The name of the license (e.g., CC-BY 4.0) should be included for each asset.
        \item For scraped data from a particular source (e.g., website), the copyright and terms of service of that source should be provided.
        \item If assets are released, the license, copyright information, and terms of use in the package should be provided. For popular datasets, \url{paperswithcode.com/datasets} has curated licenses for some datasets. Their licensing guide can help determine the license of a dataset.
        \item For existing datasets that are re-packaged, both the original license and the license of the derived asset (if it has changed) should be provided.
        \item If this information is not available online, the authors are encouraged to reach out to the asset's creators.
    \end{itemize}

\item {\bf New assets}
    \item[] Question: Are new assets introduced in the paper well documented and is the documentation provided alongside the assets?
    \item[] Answer: \answerNA{} % Replace by \answerYes{}, \answerNo{}, or \answerNA{}.
    \item[] Justification: The paper does not release new assets.
    \item[] Guidelines:
    \begin{itemize}
        \item The answer NA means that the paper does not release new assets.
        \item Researchers should communicate the details of the dataset/code/model as part of their submissions via structured templates. This includes details about training, license, limitations, etc. 
        \item The paper should discuss whether and how consent was obtained from people whose asset is used.
        \item At submission time, remember to anonymize your assets (if applicable). You can either create an anonymized URL or include an anonymized zip file.
    \end{itemize}

\item {\bf Crowdsourcing and research with human subjects}
    \item[] Question: For crowdsourcing experiments and research with human subjects, does the paper include the full text of instructions given to participants and screenshots, if applicable, as well as details about compensation (if any)? 
    \item[] Answer: \answerNA{} % Replace by \answerYes{}, \answerNo{}, or \answerNA{}.
    \item[] Justification: The paper does not involve crowdsourcing nor research with human subjects.
    \item[] Guidelines:
    \begin{itemize}
        \item The answer NA means that the paper does not involve crowdsourcing nor research with human subjects.
        \item Including this information in the supplemental material is fine, but if the main contribution of the paper involves human subjects, then as much detail as possible should be included in the main paper. 
        \item According to the NeurIPS Code of Ethics, workers involved in data collection, curation, or other labor should be paid at least the minimum wage in the country of the data collector. 
    \end{itemize}

\item {\bf Institutional review board (IRB) approvals or equivalent for research with human subjects}
    \item[] Question: Does the paper describe potential risks incurred by study participants, whether such risks were disclosed to the subjects, and whether Institutional Review Board (IRB) approvals (or an equivalent approval/review based on the requirements of your country or institution) were obtained?
    \item[] Answer: \answerNA{} % Replace by \answerYes{}, \answerNo{}, or \answerNA{}.
    \item[] Justification: The paper does not involve crowdsourcing nor research with human subjects.
    \item[] Guidelines:
    \begin{itemize}
        \item The answer NA means that the paper does not involve crowdsourcing nor research with human subjects.
        \item Depending on the country in which research is conducted, IRB approval (or equivalent) may be required for any human subjects research. If you obtained IRB approval, you should clearly state this in the paper. 
        \item We recognize that the procedures for this may vary significantly between institutions and locations, and we expect authors to adhere to the NeurIPS Code of Ethics and the guidelines for their institution. 
        \item For initial submissions, do not include any information that would break anonymity (if applicable), such as the institution conducting the review.
    \end{itemize}

\item {\bf Declaration of LLM usage}
    \item[] Question: Does the paper describe the usage of LLMs if it is an important, original, or non-standard component of the core methods in this research? Note that if the LLM is used only for writing, editing, or formatting purposes and does not impact the core methodology, scientific rigorousness, or originality of the research, declaration is not required.
    %this research? 
    \item[] Answer: \answerNA{} % Replace by \answerYes{}, \answerNo{}, or \answerNA{}.
    \item[] Justification: The core method development in this research does not involve LLMs as any important, original, or non-standard components.
    \item[] Guidelines:
    \begin{itemize}
        \item The answer NA means that the core method development in this research does not involve LLMs as any important, original, or non-standard components.
        \item Please refer to our LLM policy (\url{https://neurips.cc/Conferences/2025/LLM}) for what should or should not be described.
    \end{itemize}

\end{enumerate}
\fi







\end{document}
